\section{Standardy migracji --- secondkey-standards}\label{sec:standardy}

Repozytorium \path{letsgolegacy.secondkey-standards} jest komponentem C4 Second Key:
pakietem standardów architektonicznych i~migracyjnych, które są zapisane raz i~z~jednego
źródła trafiają w~dwie postacie -- do agenta migracyjnego jako skill oraz do bramki
(Portcullis, rozdział~\ref{sec:portcullis}) jako konfiguracja analyzerów (czytana w~buildzie
z~pakietem analyzerów, nie w~skanie CLI Portcullis; sekcja~\ref{sec:standardy-wagi-kanaly}).
Rozdział opisuje format standardu, wszystkie dwadzieścia reguł, generator, który je waliduje
i~kompiluje, wersjonowanie, procedurę zmiany standardu, to, jak wagi standardów działają
w~kanałach Portcullis, skill dla GitHub Copilot upgrade, drift check dla repozytoriów
korzystających ze standardów, procedurę użycia z~modernize-dotnet, testy, CI i~stan prac.

Stan według README: faza 01 w~budowie; standardy, ich walidator, generator, skill dla agenta
modernizacyjnego GitHub Copilot i~drift check dla repozytoriów konsumentów istnieją. Żadna
wersja nie została jeszcze wydana (otagowana).

\zrodla{\path{letsgolegacy.secondkey-standards/README.md}}

\subsection{Po co istnieje}\label{sec:standardy-cel}

W~łańcuchu Second Key (rozdział~\ref{sec:architektura}) krok 3 to „Standards”: komponent C4
dostarcza standardy agentowi migracyjnemu jako skills, a~wynikiem kroku jest pull request
agenta. Krok 4 to bramka C5, która sprawdza zmienione wiersze tego pull requestu. Oba kroki
muszą mówić o~tych samych regułach z~tą samą wagą. Plan pracy repozytorium formułuje to
wprost: jedno źródło standardów w~Markdown, z~metadanymi reguł, kompilowane do wszystkiego,
co musi się co do nich zgadzać -- katalogu Agent Skills ładowanego przez agenta, pakietu
konfiguracji analyzerów czytanego przez bramkę oraz pliku \repopath{.editorconfig}.
\finding{Rozbieżność między tym, co powiedziano agentowi, a~tym, co egzekwuje bramka, jest
błędem CI, a~nie znaleziskiem.} W~praktyce CI tego repozytorium pilnuje zgodności
wygenerowanych plików ze źródłami (\repopath{generate -{}-check}); nie porównuje ich
z~poziomami, które emituje Portcullis, więc rozbieżność tego drugiego rodzaju istnieje
i~opisuje ją sekcja~\ref{sec:standardy-wagi-kanaly}.

W~tabeli artefaktów rdzenia katalog \repopath{standards/} ma format „Markdown → skills
i~NuGet”, produkuje go C4, a~czytają agent migracyjny i~bramka C5. Według planu pracy
treść standardów pochodzi z~repozytorium \path{architecture-standards} (konstytucja
architektury i~przewodniki estate) oraz ze standardów specyficznych dla migracji z~.NET
Framework do nowoczesnego .NET; reguły cytują z~tego repozytorium jednak tylko konstytucję
(siedem reguł \repopath{SK-ARCH-} powtarza jej zasady P4, P5, P9, P10 i~P15), a~żadna nie
cytuje przewodników. \path{architecture-standards} to osobne repozytorium, nie ten zbiór
dwudziestu reguł.
Drift check przejmuje ideę repozytorium \repopath{context-pin} -- wersję standardów
przypiętą w~każdym repozytorium -- bez osobnej usługi: wystarcza pakiet NuGet, tag git
i~jeden krok CI.

Zbiór obejmuje dwadzieścia reguł dla migracji .NET Framework → .NET~10; są to standardy
ogólne, nie wyprowadzone z~systemu klienta (standardy klienta to faza~02). Trzynaście jest
specyficznych dla migracji (prefiks \repopath{SK-MIG-}), siedem przenosi na grunt migracji
zasady konstytucji architektury (prefiks \repopath{SK-ARCH-}) i~linkuje zasadę, którą
powtarza. Tylko jedenaście reguł ma przypisaną diagnostykę, czyli sprawdzenie wykonywalne:
cztery migracyjne reguły Portcullis i~pięć architektonicznych reguł Portcullis (dziewięć
diagnostyk Portcullis) oraz dwie reguły globalizacji z~analyzerami .NET SDK (CA1310;
CA1304, CA1305, CA1311), które uruchamia build, a~nie \repopath{portcullis scan}. Pozostałych
dziewięć reguł jest wyłącznie instrukcjami dla agenta. Regułą, na której opierają się
pozostałe, jest SK-MIG-011: zachowaj zachowanie i~oflaguj podejrzewany defekt zamiast go
naprawiać.

\begin{xltabular}{\linewidth}{@{}>{\raggedright\arraybackslash}p{4.6cm} >{\raggedright\arraybackslash}X@{}}
\toprule
Ścieżka & Zawartość \\
\midrule
\endfirsthead
\toprule
Ścieżka & Zawartość \\
\midrule
\endhead
\path{standards/} & Reguły: jeden plik Markdown na regułę, z~metadanymi. \path{standards/README.md} definiuje format \\
\path{catalog/pack.json} & Wersja pakietu, nazwy skilla i~pakietu NuGet, zamknięte słowniki, względem których walidowane są reguły \\
\path{catalog/skill.template.md} & Pisana ręcznie treść skilla; generator wstawia tabele reguł i~bramki \\
\path{generated/} & Generowane, nigdy nie edytowane ręcznie: skill, konfiguracja analyzerów, projekt pakietu NuGet, manifest \\
\path{src/SecondKey.Standards.Generator} & Narzędzie \repopath{secondkey-standards} (.NET~10): waliduje reguły, generuje \path{generated/}, sprawdza konsumenta pod kątem dryfu \\
\path{actions/drift-check/} & Drift check jako jeden krok CI dla repozytorium konsumenta \\
\path{tests/} & Testy narzędzia, w~tym testy trzymające reguły i~drzewo generowane tego repozytorium przy kontrakcie, oraz repozytoria-fixture dla drift checka \\
\path{scripts/} & \path{setup.sh}, \path{verify-package.sh}, \path{scan-secrets.sh}, hook \path{hooks/pre-commit} \\
\path{CHANGELOG.md} & Co zmieniło się w~każdej wersji \\
\path{docs/USING-WITH-MODERNIZE-DOTNET.md} & Instalacja skilla w~repozytorium docelowym i~uruchomienie agenta modernizacyjnego GitHub Copilot \\
\path{docs/adr/} & Decyzje (ADR 0001--0005) i~ich uzasadnienia \\
\bottomrule
\end{xltabular}

Rejestr decyzji (\path{docs/adr/README.md}) podaje tytuły; statusy są w~nagłówkach samych
ADR, a~każdy ADR ma datę 2026-09-26:

\begin{xltabular}{\linewidth}{@{}>{\raggedright\arraybackslash}p{1.1cm} >{\raggedright\arraybackslash}X >{\raggedright\arraybackslash}p{4.2cm}@{}}
\toprule
ADR & Decyzja & Status \\
\midrule
\endfirsthead
\toprule
ADR & Decyzja & Status \\
\midrule
\endhead
0001 & Format źródłowy: jeden plik Markdown na regułę, metadane w~front matter, jedno słownictwo wag dla agenta i~bramki (sekcja~\ref{sec:standardy-format}) & przyjęty (C4-a) \\
0002 & Generator jako narzędzie .NET~10, które waliduje standardy od pierwszego pull requestu (sekcja~\ref{sec:standardy-generator}) & przyjęty (C4-a) \\
0003 & Drzewo generowane, \repopath{generate -{}-check}, bramka wersji i~pakiet NuGet zawierający wyłącznie konfigurację (sekcja~\ref{sec:standardy-wersje}) & przyjęty (C4-b) \\
0004 & Standardy jako własny skill agenta modernizacyjnego GitHub Copilot: format, lokalizacja, nazwa, wykrywanie (sekcja~\ref{sec:standardy-skill}) & przyjęty (R4); podjęcie skilla w~prawdziwym przebiegu ma potwierdzić P6; poprawka z~2026-09-26 \\
0005 & Drift check: pin czytany z~referencji pakietu i~ze skilla, wydania z~tagów git, jeden krok CI, bez usługi (sekcja~\ref{sec:standardy-drift}) & przyjęty (C4-c) \\
\bottomrule
\end{xltabular}

Indeks wyjaśnia, po co ADR zapisują odrzucone warianty: dokument, który mówi
„rozważyliśmy X i~odrzuciliśmy go, bo Y”, jest wart więcej niż taki, który wymienia
polecenia (P14).

Rysunek~\ref{fig:standardy} pokazuje, jak jedno źródło reguł trafia do agenta
migracyjnego i~do buildu kandydata. Drift check (\path{actions/drift-check}, krawędź
przerywana) czyta oba piny repozytorium konsumenta -- \repopath{metadata.version}
zainstalowanego skilla i~referencję pakietu -- i~porównuje je z~najwyższym tagiem
\repopath{v<MAJOR.MINOR.PATCH>} repozytorium standardów (sekcja~\ref{sec:standardy-drift}).
Rysunek pokazuje zamierzony przepływ, a~nie stan sprawdzenia: tagów wydań jeszcze nie ma
(drift check zwraca wtedy \repopath{no-release}), a~podjęcie skilla przez agenta nie zostało
dotąd sprawdzone (P6 jest planowane).

\begin{figure}[htbp]
  \centering
  \includegraphics[width=\linewidth]{\dgm{a9-standardy}}
  \caption{Standardy: jedno źródło, wyjścia generatora i~ich konsumenci.}
  \label{fig:standardy}
\end{figure}

Zależności są pinowane centralnie w~\path{Directory.Packages.props} (central package
management z~\path{CentralPackageTransitivePinningEnabled}); budżet zależności jest
opublikowany w~README, a~każda nowa linia w~tym pliku jest decyzją, o~którą można zapytać
w~przeglądzie.

\begin{tabularx}{\linewidth}{@{}>{\raggedright\arraybackslash}p{4.6cm} >{\raggedright\arraybackslash}p{3.2cm} >{\raggedright\arraybackslash}X@{}}
\toprule
Projekt & Pakiety & Powód \\
\midrule
\path{SecondKey.Standards.Generator} & \path{System.CommandLine} 2.0.12, \path{YamlDotNet} 18.1.0 & Powierzchnia poleceń; front matter YAML z~numerami wierszy (ADR 0002) \\
\path{SecondKey.Standards.Generator.Tests} & \path{xunit.v3} 4.0.1 & Testy na Microsoft.Testing.Platform \\
\path{SecondKey.Standards} (pakiet generowany) & brak & Tylko konfiguracja: bez assembly i~bez zależności \\
\bottomrule
\end{tabularx}

Licencja (\path{LICENSE}, README): wszystkie prawa zastrzeżone. Repozytorium jest
opublikowane po to, żeby można było obejrzeć jego pracę; poza prawami, które warunki
korzystania z~GitHub dają każdemu użytkownikowi (oglądanie i~fork repozytoriów
publicznych), nie jest udzielana licencja na używanie, kopiowanie, modyfikowanie, łączenie,
publikowanie, dystrybucję, sublicencjonowanie ani sprzedaż oprogramowania i~dokumentacji.
Model licencjonowania jest decyzją otwartą; do jej podjęcia ta informacja dotyczy każdego
pliku, który nie stanowi inaczej. Pakiet \repopath{SecondKey.Standards} zawiera ten sam
plik \path{LICENSE}. Reguły \repopath{SK-ARCH-} powtarzają zasady
\path{architecture-standards} (MIT, ten sam autor), a~skrypty skanu sekretów są
przeniesione z~tamtego repozytorium.

Konwencje repozytorium (\path{AGENTS.md}): reguła mówi, co zmigrowany kod musi lub czego
nie wolno mu robić, dlaczego, i~pokazuje przykład niezgodny oraz zgodny; twierdzenia
techniczne cytują źródło pierwotne (dokumentację Microsoft, konstytucję) w~sekcji
\repopath{References}; treść reguł używa wyłącznie linków bezwzględnych, bo jest kopiowana
do \path{references/} skilla; diagnostykę Portcullis dopisuje się do \path{catalog/pack.json}
dopiero wtedy, gdy bramka ją dostarcza; kod samego narzędzia stosuje publikowane reguły
(porównania porządkowe, jawna kultura); repozytorium jest publiczne, więc tylko język
angielski, bez sekretów i~bez informacji biznesowych. Zasady architektury nie są tu
wyprowadzane na nowo -- konstytucja pozostaje źródłem rozstrzygającym, a~reguły
\repopath{SK-ARCH-} ją powtarzają i~linkują.

Co edytuje się przy jakiej zmianie (\path{AGENTS.md}); niczego pod \path{generated/} nie
pisze się ręcznie:

\begin{tabularx}{\linewidth}{@{}>{\raggedright\arraybackslash}p{5.4cm} >{\raggedright\arraybackslash}X >{\raggedright\arraybackslash}p{2.0cm}@{}}
\toprule
Zmiana & Plik do edycji & Potem \\
\midrule
Reguła & \path{standards/<id>-<slug>.md} & \repopath{generate} \\
Wersja pakietu standardów, nazwa lub opis skilla, pakiet NuGet & \path{catalog/pack.json} & \repopath{generate} \\
Ręcznie pisana treść skilla & \path{catalog/skill.template.md} & \repopath{generate} \\
Opis tego, co zmieniła wersja & \path{CHANGELOG.md} & -- \\
\bottomrule
\end{tabularx}

Przed edycją reguły czyta się \path{standards/README.md}, który definiuje format oraz
znaczenie każdego pola i~każdej wagi. Decyzje i~ich powody są w~\path{docs/adr/}; zmiana,
która odwraca decyzję ADR, aktualizuje ten ADR.

Repozytorium deklaruje \path{architecture-standards} tak samo jak rdzeń
(rozdział~\ref{sec:rdzen}): \path{.claude/settings.json} rejestruje marketplace
\texttt{architecture-standards} i~włącza plugin \texttt{architecture-core}. Pliki higieny
powołują się na \path{REPO-BASELINE.md} z~\path{architecture-standards}: §1 --
\path{.editorconfig}, \path{.gitattributes}, \path{.github/CODEOWNERS},
\path{.github/dependabot.yml}, \path{Directory.Build.props},
\path{Directory.Packages.props} i~CodeQL; §2 -- skan sekretów w~hooku pre-commit i~w~CI;
§3 -- onboarding jednym poleceniem (\path{scripts/setup.sh}); §4 -- skrypty odtwarzające
job CI 1:1 (\path{scripts/scan-secrets.sh}, \path{scripts/verify-package.sh}).

\zrodla{\path{letsgolegacy.secondkey-standards/README.md},
\path{letsgolegacy.secondkey-standards/AGENTS.md},
\path{letsgolegacy.secondkey-standards/LICENSE},
\path{letsgolegacy.secondkey-standards/docs/WORKPLAN.md},
\path{letsgolegacy.secondkey-standards/docs/adr/README.md},
\path{letsgolegacy.secondkey-standards/standards/} (sekcje \emph{References} reguł),
\path{letsgolegacy.secondkey-standards/Directory.Packages.props},
\path{letsgolegacy.secondkey-standards/Directory.Build.props},
\path{letsgolegacy.secondkey-standards/.claude/settings.json},
\path{letsgolegacy.secondkey-standards/.editorconfig},
\path{letsgolegacy.secondkey-standards/.gitattributes},
\path{letsgolegacy.secondkey-standards/.github/CODEOWNERS},
\path{letsgolegacy.secondkey-standards/.github/dependabot.yml},
\path{letsgolegacy.secondkey-standards/.github/workflows/codeql.yml},
\path{letsgolegacy.secondkey-standards/.github/workflows/secret-scan.yml},
\path{letsgolegacy.secondkey-standards/.gitleaks.toml},
\path{letsgolegacy.secondkey-standards/scripts/setup.sh},
\path{letsgolegacy.secondkey-standards/scripts/hooks/pre-commit},
\path{letsgolegacy.secondkey-standards/scripts/scan-secrets.sh},
\path{letsgolegacy.secondkey-standards/scripts/verify-package.sh},
\path{letsgolegacy.secondkey/README.md},
\path{letsgolegacy.secondkey/docs/WORKPLAN.md}}

\subsection{Format źródłowy standardu (ADR 0001)}\label{sec:standardy-format}

C4 kompiluje jedno źródło standardów dla dwóch konsumentów, którzy nie mogą się ze sobą nie
zgadzać: skilla ładowanego przez agenta migracyjnego i~konfiguracji analyzerów egzekwowanej
przez bramkę. Źródło musi więc nieść dla każdej reguły zarówno prozę dla agenta, jak
i~metadane dla bramki, w~postaci, którą człowiek przejrzy w~pull requeście, a~program
sprawdzi.

\subsubsection{Plik reguły}

Każda reguła to jeden plik \path{standards/<id>-<slug>.md}, np.
\path{SK-MIG-005-explicit-string-comparison.md}. Plik \path{standards/README.md} jest
dokumentacją, nie regułą, i~generator go pomija. Kształt pliku (wartości ilustracyjne,
przepisane z~\path{standards/README.md}):

\begin{lstlisting}
---
id: SK-MIG-000
title: One imperative line the agent can act on
severity: warning
category: globalization
appliesTo:
  - "*.cs"
portcullisRule: PORTCULLIS_MIG_EXAMPLE   # optional
analyzers:                               # optional
  - CA1310
principle: P5                            # optional
---

The statement: what migrated code must or must not do, in a paragraph or two.

## Rationale
## Non-compliant
## Compliant
## Migration                 (optional)
## Flag instead of fixing    (optional)
## References                (optional)
\end{lstlisting}

\subsubsection{Front matter}

Plik musi zaczynać się w~pierwszym wierszu blokiem \repopath{-{}-{}-}, a~blok musi być
zamknięty drugim wierszem \repopath{-{}-{}-}. Metadane są czytane przez YamlDotNet; błąd składni
YAML jest raportowany z~numerem wiersza w~pliku. Klucz spoza listy poniżej jest błędem, a~nie
jest ignorowany. Wartości skalarne muszą mieścić się w~jednym wierszu i~nie mogą być puste.

\begin{xltabular}{\linewidth}{@{}>{\raggedright\arraybackslash}p{2.9cm} >{\raggedright\arraybackslash}p{1.9cm} >{\raggedright\arraybackslash}X@{}}
\toprule
Klucz & Wymagany & Znaczenie i~walidacja \\
\midrule
\endfirsthead
\toprule
Klucz & Wymagany & Znaczenie i~walidacja \\
\midrule
\endhead
\path{id} & tak & Segmenty wielkimi literami, zakończone trzycyfrowym numerem (\repopath{SK-MIG-001}); wzorzec \path{^[A-Z][A-Z0-9]*(-[A-Z][A-Z0-9]*)*-[0-9]{3}$}. Stabilny: agenci cytują go w~komentarzach kodu, a~wyniki bramki są do niego śledzone. Nazwa pliku musi brzmieć \path{<id>.md} albo \path{<id>-<slug>.md}, gdzie slug to małe litery i~cyfry rozdzielone pojedynczymi myślnikami \\
\path{title} & tak & Jedno zdanie w~trybie rozkazującym, najwyżej 120 znaków. Tę treść widzi agent w~tabeli reguł skilla \\
\path{severity} & tak & \repopath{error}, \repopath{warning} albo \repopath{suggestion} (sekcja~\ref{sec:standardy-wagi}) \\
\path{category} & tak & Jedna z~kategorii z~\path{catalog/pack.json}: \repopath{hosting}, \repopath{layering}, \repopath{concurrency}, \repopath{configuration}, \repopath{globalization}, \repopath{dependency-injection}, \repopath{outbound-http}, \repopath{data-access}, \repopath{behaviour}, \repopath{observability} \\
\path{appliesTo} & tak & Lista wzorców plików \repopath{.editorconfig} (\repopath{"*.cs"}, \repopath{"*.csproj"}), co najmniej jeden. Wzorce trzeba ujmować w~cudzysłów: niecytowana \repopath{*} rozpoczyna alias YAML, a~komunikat błędu walidatora to podpowiada. Wzorzec nie może zawierać spacji, zaczynać się od \repopath{/} ani zawierać \repopath{[}, \repopath{]}, \repopath{\#}, \repopath{;}. Wygenerowany \repopath{.editorconfig} zawęża każdą diagnostykę do tych wzorców \\
\path{portcullisRule} & nie & Diagnostyka Portcullis, która sprawdza regułę. Format \path{PORTCULLIS_<SLUG>} z~podkreślnikami, bo Roslyn odrzuca identyfikator diagnostyki z~myślnikiem. Musi być na liście \repopath{portcullisRules} w~\path{catalog/pack.json} -- to kontrakt z~bramką \\
\path{analyzers} & nie & Identyfikatory analyzerów .NET SDK (\repopath{CAxxxx}, \repopath{IDExxxx}), które sprawdzają regułę; bez powtórzeń. Dostarcza je SDK, więc build z~konfiguracją standardów może je egzekwować bez nowych narzędzi (są domyślnie wyłączone, a~\repopath{portcullis scan} ich nie uruchamia). Identyfikatory Portcullis należą do \repopath{portcullisRule} \\
\path{principle} & nie & Zasada konstytucji, którą reguła powtarza lub na której się opiera: \repopath{P1}--\repopath{P15}, także z~literowym uzupełnieniem, np. \repopath{P2a} \\
\bottomrule
\end{xltabular}

Walidacja obejmująca cały zbiór: identyfikatory są unikalne (duplikat jest zgłaszany na
drugim pliku), a~każda diagnostyka jest przypisana tylko jednej regule -- dwie reguły dałyby
jej dwie wagi. Brak katalogu \path{standards/} („\texttt{the standards directory does not
exist}”) albo katalog bez plików reguł („\texttt{no rule files found; each rule is one
standards/<id>-<slug>.md file}”) jest problemem, a~nie zaliczoną walidacją.
\path{README.md} jest pomijany bez względu na wielkość liter nazwy. Przed parsowaniem
generator usuwa z~pliku reguły znacznik BOM UTF-8 i~zamienia końce linii CRLF na LF, więc
reguła zapisana w~edytorze z~CRLF znaczy to samo.

\subsubsection{Wagi}\label{sec:standardy-wagi}

Jedno słownictwo dla agenta i~bramki. Nazwy są wagami \repopath{.editorconfig}, na które
ustawiane są przypisane diagnostyki.

\begin{tabularx}{\linewidth}{@{}>{\raggedright\arraybackslash}p{2.4cm} >{\raggedright\arraybackslash}X >{\raggedright\arraybackslash}X@{}}
\toprule
Waga & Dla agenta & Dla bramki \\
\midrule
\repopath{error} & Zmigrowany kod nie może naruszać reguły. Jeśli nie da się jej spełnić bez zmiany zachowania, agent zatrzymuje się i~flaguje -- nigdy nie dostarcza naruszenia po cichu & Diagnostyka jest błędem builda \\
\repopath{warning} & Naprawia to, czego dotyka migracja; każde pozostawione naruszenie jest zapisane jako flaga zachowania z~powodem & Diagnostyka jest raportowana \\
\repopath{suggestion} & Stosuje regułę tam, gdzie zmiana jest mechaniczna i~neutralna dla zachowania & Raportowana jako sugestia; nigdy nie blokuje \\
\bottomrule
\end{tabularx}

Waga \repopath{suggestion} nie zostawia śladu w~standardowym logu \texttt{dotnet build}:
w~logu błędów kompilatora (\repopath{ErrorLog}) ma poziom \emph{note} (sprawdzone na CA1310
oraz, z~pakietem analyzerów zbudowanym ze źródeł Portcullis, na
\path{PORTCULLIS_P10_CUSTOM_BASE_CLASS} i~\path{PORTCULLIS_P15_MISSING_SERVICE_DEFAULTS}).
Dla SK-ARCH-006 i~SK-ARCH-007 sam build nie wypisuje więc żadnego ostrzeżenia.

Reguły bez przypisanej diagnostyki (bez \repopath{portcullisRule} i~bez
\repopath{analyzers}) nie mają sprawdzenia wykonywalnego: są instrukcjami dla agenta,
a~w~zakresie zachowania ma je weryfikować replay i~porównanie Second Key, nie analyzer.

\subsubsection{Treść reguły}

\begin{itemize}
  \item \textbf{Stwierdzenie} -- wszystko przed pierwszym nagłówkiem \repopath{\#\#}. Normatywne,
        krótkie, sprawdzalne. Brak stwierdzenia jest błędem.
  \item \textbf{Rationale} (wymagana, niepusta) -- dlaczego. Według README powód jest tą
        częścią reguły, która przetrwa zderzenie z~przypadkiem, którego reguła nie przewidziała.
  \item \textbf{Non-compliant} i~\textbf{Compliant} (wymagane) -- co najmniej jeden blok kodu
        w~ogrodzeniu (fence) w~każdej z~tych sekcji; blok bez niepustego wiersza nie liczy się
        jako przykład. Przykłady używają kształtów sklepu (zamówienia, ceny, rabaty), bo tam
        migracje się psują.
  \item \textbf{Migration} (opcjonalna) -- jak przejść od kształtu legacy do zgodnego bez
        zmiany zachowania.
  \item \textbf{Flag instead of fixing} (opcjonalna) -- różnice zachowania, które agent
        zapisuje zamiast je rozwiązywać.
  \item \textbf{References} (opcjonalna) -- źródła pierwotne.
\end{itemize}

Nagłówek poziomu pierwszego (\repopath{\#}) jest zabroniony: tytuł pochodzi z~front matter
i~jest podany raz. Nagłówek \repopath{\#\#} spoza zamkniętej listy jest błędem, podobnie jak
powtórzona sekcja. Parser śledzi ogrodzenia kodu (\repopath{`{}`{}`} lub \repopath{\textasciitilde\textasciitilde\textasciitilde}, co najmniej trzy znaki, do trzech
spacji wcięcia), więc \repopath{\#\# heading} wewnątrz bloku kodu jest treścią, a~nie sekcją;
niezamknięty blok jest zgłaszany w~wierszu, który go otworzył. Link względny w~treści
(\repopath{](plik.md)}) jest błędem, bo w~wygenerowanym skillu wskazywałby donikąd -- należy
użyć bezwzględnego URL albo zacytować inną regułę jej identyfikatorem. Dozwolone są cele
linków zaczynające się od \repopath{http://}, \repopath{https://}, \repopath{mailto:}
i~\repopath{\#} (kotwica w~dokumencie); wierszy wewnątrz bloków kodu parser pod tym kątem
nie sprawdza.

\subsubsection{Decyzje ADR 0001 i~ich powody}

\begin{enumerate}
  \item \textbf{Jeden plik Markdown na regułę, z~front matter YAML.} Reguła jest jednostką
        przeglądu, wersjonowania i~cytowania; jeden plik na regułę sprawia, że diff dotyczący
        jednej reguły czyta się jako jeden. Front matter trzyma metadane obok prozy, którą
        opisują.
  \item \textbf{Identyfikatory \repopath{PREFIX-NNN}.} Numer utrzymuje identyfikator stabilnym,
        gdy tytuł jest przeredagowany; prefiks mówi, skąd reguła pochodzi. Walidator
        akceptuje dowolny prefiks z~wielkich liter, więc standardy klienta (faza 02) mogą
        mieć własny.
  \item \textbf{Jedno słownictwo wag dla agenta i~bramki} (\repopath{error}, \repopath{warning},
        \repopath{suggestion} -- nazwy z~\repopath{.editorconfig}). Reguła nie może znaczyć
        „musi” dla agenta i~„może” dla bramki.
  \item \textbf{Zamknięte słowniki w~\path{catalog/pack.json}}: kategorie oraz identyfikatory
        diagnostyk Portcullis, na które reguła może się mapować. Lista Portcullis jest
        kontraktem z~C5: reguła wskazująca identyfikator spoza tej listy nie przechodzi
        walidacji, zamiast po cichu mapować się na nic. Listę utrzymuje się ręcznie, a~to, że
        bramka dostarcza wpisany identyfikator, sprawdza człowiek
        (sekcja~\ref{sec:standardy-ograniczenia}). Cztery identyfikatory
        migracyjne ustala backlog między repozytoriami
        (\path{PORTCULLIS_MIG_SYSTEM_WEB}, \path{PORTCULLIS_MIG_HTTPCONTEXT_CURRENT},
        \path{PORTCULLIS_MIG_SYNC_OVER_ASYNC}, \path{PORTCULLIS_MIG_CONFIGURATION_MANAGER});
        identyfikatory zasad pochodzą z~istniejącego rejestru reguł Portcullis.
  \item \textbf{\repopath{portcullisRule} ma jedną wartość; diagnostyka należy do jednej
        reguły.} Tam, gdzie jedna zasada ma dwie reguły Portcullis (P4), powstają dwa
        standardy (SK-ARCH-001, SK-ARCH-002). Diagnostyka przypisana dwukrotnie wymagałaby
        dwóch wag, więc walidator to odrzuca.
  \item \textbf{Pole \repopath{analyzers}, ponad listę z~ticketu.} Ticket wskazuje
        \repopath{portcullisRule} jako mapowanie; reguły globalizacji -- pułapka migracyjna, wokół
        której zbudowane jest demo -- nie mają reguły Portcullis, ale .NET SDK dostarcza dla
        nich analyzery (CA1310, CA1304, CA1305, CA1311), domyślnie wyłączone. Mapowanie ich
        z~tego samego źródła wprowadza reguły NLS → ICU do konfiguracji bramki bez nowych
        narzędzi (zasada produktu nr 3), więc agent i~build z~tą konfiguracją zgadzają się także
        co do nich (\repopath{portcullis scan} analyzerów SDK nie uruchamia,
        sekcja~\ref{sec:standardy-wagi-kanaly}). Pole przyjmuje tylko identyfikatory SDK.
  \item \textbf{\repopath{appliesTo} to lista wzorców plików \repopath{.editorconfig}.} Jest
        jednocześnie dokumentacją dla agenta i~zakresem sekcji wygenerowanego
        \repopath{.editorconfig}.
  \item \textbf{Treść ma stały zestaw sekcji.} Zamknięty zestaw wyłapuje nagłówek z~literówką,
        który inaczej wypadłby z~każdego wygenerowanego artefaktu.
  \item \textbf{Protokół flag jest zdefiniowany w~regule SK-MIG-011, a~nie w~narzędziach.}
        „Flag instead of fixing” jest centralnym zachowaniem, o~które proszą standardy, więc
        jego format -- komentarz \repopath{SECONDKEY-FLAG <rule-id>:} i~wiersz w~pliku
        \path{docs/migration/behaviour-flags.md} repozytorium docelowego -- jest wersjonowany
        razem z~regułami i~dociera do agenta tym samym kanałem.
  \item \textbf{Reguły architektoniczne powtarzają, nie zastępują.} Każda reguła
        \repopath{SK-ARCH-} mówi, czego jedna zasada wymaga od migracji, i~linkuje kotwicę
        zasady w~konstytucji (\path{00-REFERENCE-ARCHITECTURE.md#p4}); test sprawdza ten link.
        Wybrane zasady to te, których dotyka migracja z~.NET Framework: P4 (schemat, dane
        seedujące), P5 (sekrety), P9 (Program.cs, warstwy kontrolerów), P10 (rozszerzanie przez
        DI) i~P15 (obserwowalność). P2, P3, P6--P8 i~P11--P14 są poza zasięgiem migracji albo
        są już niesione przez pole \repopath{principle} reguły migracyjnej (P2 w~SK-MIG-009,
        P13 w~SK-MIG-013, P14 w~SK-MIG-011).
\end{enumerate}

\subsubsection{Wybrane wagi i~powody}

\begin{tabularx}{\linewidth}{@{}>{\raggedright\arraybackslash}p{5.0cm} >{\raggedright\arraybackslash}p{2.4cm} >{\raggedright\arraybackslash}X@{}}
\toprule
Reguły & Waga & Powód \\
\midrule
System.Web, \repopath{HttpContext.Current}, sync-over-async, \repopath{ConfigurationManager} & \repopath{error} & Każda z~nich albo zawodzi w~czasie działania w~sposób, którego nie widzi żaden test (ustawienie \repopath{null}, wyczerpana pula wątków), albo zostawia migrację niedokończoną \\
Porównania napisów, kultura (CA1310, CA1304, CA1305, CA1311) & \repopath{warning} & Jawne wywołanie jest mechaniczne i~neutralne dla zachowania, ale kod legacy ma wiele miejsc wywołań; raportowane wszędzie, naprawiane tam, gdzie migracja dotyka kodu \\
Tryb globalizacji, zachowanie EF, zachowanie ogólnie, kontrakt HTTP, testy & \repopath{error} & To zachowanie, które migracja ma zachować; żaden analyzer go nie sprawdza, więc kontrolą jest instrukcja agenta, a~weryfikacją -- replay \\
DI, \repopath{IHttpClientFactory}, przeniesiony \repopath{DbContext} w~kontrolerach, dane seedujące & \repopath{warning} & Dług do usunięcia, ale nie kosztem zmiany zachowania w~zmianie migracyjnej \\
Program.cs jako manifest, klasy bazowe, okablowanie obserwowalności & \repopath{suggestion} & Struktura, nie zachowanie \\
\bottomrule
\end{tabularx}

\subsubsection{Konsekwencje i~odrzucone alternatywy}

Walidator (\repopath{secondkey-standards validate}) sprawdza wszystko powyższe i~raportuje
każdy problem z~plikiem i~wierszem; CI uruchamia go na każdym pull requeście. Kryterium
akceptacji C4-a jest testem (\repopath{RepositoryStandardsTests}), a~nie listą kontrolną
recenzenta. Dodanie diagnostyki Portcullis do reguły to zmiana w~dwóch plikach
(\path{catalog/pack.json} i~reguła), wykonywana dopiero, gdy bramka ją dostarcza. Reguły są
pisane dla dwóch czytelników naraz: człowieka przeglądającego standard i~agenta, który go
stosuje.

Odrzucone warianty:
\begin{itemize}
  \item \textbf{Jeden duży dokument z~kotwicami} -- łatwiej czytać od początku do końca, ale
        reguły nie dałoby się wersjonować, walidować ani cytować osobno, a~generator musiałby
        wydobywać strukturę z~prozy.
  \item \textbf{Metadane w~osobnym katalogu JSON lub YAML, proza w~Markdown} -- dwa pliki na
        regułę, które mogą się nie zgadzać; front matter trzyma je w~jednym.
  \item \textbf{Wagi liczbowe albo osobna waga dla bramki} -- dwa słowniki to dokładnie ten
        dryf, któremu komponent ma zapobiegać.
\end{itemize}

Zmiana reguły według \path{standards/README.md}: edycja lub dodanie pliku (stwierdzenie
sprawdzalne, przykłady minimalne); \repopath{validate} musi przejść; pull request cytuje
dowód -- dokumentację, kod lub zaobserwowane zachowanie, na którym reguła się opiera.
Szablon zgłoszenia \path{.github/ISSUE_TEMPLATE/rule-gap.yml} („Gap in a migration
standard”) wymaga identyfikatora reguły (albo informacji, że reguły brakuje), rodzaju luki
(\emph{Missing} -- pułapka migracyjna, której nie pokrywa żadna reguła, \emph{Wrong} --
reguła przeczy udokumentowanemu zachowaniu .NET, \emph{Drift} -- agent usłyszał co innego,
niż egzekwuje bramka, \emph{Unclear} -- reguła jest zapisana, ale w~praktyce nie da się jej
rozstrzygnąć) oraz dowodu: kodu, zaobserwowanego zachowania albo dokumentacji; reguła
zmieniona bez dowodu to błąd, któremu standardy mają zapobiegać. Pole „What the rule should
say” jest opcjonalne. \path{.github/ISSUE_TEMPLATE/config.yml} dopuszcza puste zgłoszenia
i~kieruje do konstytucji architektury: luka w~zasadzie należy tam, luka w~regule
migracyjnej -- tutaj.

\zrodla{\path{letsgolegacy.secondkey-standards/docs/adr/0001-standards-source-format.md},
\path{letsgolegacy.secondkey-standards/standards/README.md},
\path{letsgolegacy.secondkey-standards/catalog/pack.json},
\path{letsgolegacy.secondkey-standards/src/SecondKey.Standards.Generator/Parsing/RuleParser.cs},
\path{letsgolegacy.secondkey-standards/src/SecondKey.Standards.Generator/Parsing/MarkdownBody.cs},
\path{letsgolegacy.secondkey-standards/src/SecondKey.Standards.Generator/Parsing/FrontMatter.cs},
\path{letsgolegacy.secondkey-standards/src/SecondKey.Standards.Generator/Parsing/SourceText.cs},
\path{letsgolegacy.secondkey-standards/src/SecondKey.Standards.Generator/Model/Patterns.cs},
\path{letsgolegacy.secondkey-standards/src/SecondKey.Standards.Generator/StandardsLoader.cs},
\path{letsgolegacy.secondkey-standards/.github/ISSUE_TEMPLATE/rule-gap.yml},
\path{letsgolegacy.secondkey-standards/.github/ISSUE_TEMPLATE/config.yml}}

\subsection{Przegląd standardów}\label{sec:standardy-przeglad}

Tabela zbiorcza wszystkich dwudziestu reguł w~wersji 0.1.0. Tytuły są streszczone po polsku;
oryginalny tytuł (to, co agent widzi w~tabeli reguł skilla) podaje opis każdej reguły
w~sekcjach~\ref{sec:standardy-arch} i~\ref{sec:standardy-mig}. Kolumna „Diagnostyka” podaje
identyfikator, którym bramka lub .NET SDK sprawdza regułę; kreska oznacza regułę bez
analyzera: jest tylko instrukcją dla agenta, a~w~zakresie zachowania ma ją weryfikować replay
i~porównanie Second Key.

{\small
\begin{xltabular}{\linewidth}{@{}>{\raggedright\arraybackslash}p{2.3cm} >{\raggedright\arraybackslash}X >{\raggedright\arraybackslash}p{1.95cm} >{\raggedright\arraybackslash}p{3.9cm} >{\raggedright\arraybackslash}p{1.15cm}@{}}
\toprule
Reguła & Treść & Waga & Diagnostyka & Zasada \\
\midrule
\endfirsthead
\toprule
Reguła & Treść & Waga & Diagnostyka & Zasada \\
\midrule
\endhead
SK-ARCH-001 & Istniejący schemat jest kontraktem; migracje, nigdy \texttt{EnsureCreated} & \texttt{error} & \path{PORTCULLIS_P4_ENSURE_CREATED_OUTSIDE_TEST} & P4 \\
SK-ARCH-002 & Dane seedujące EF6 w~wersjonowanym kroku, nie w~\texttt{HasData} & \texttt{warning} & \path{PORTCULLIS_P4_SEED_DATA_IN_MODEL} & P4 \\
SK-ARCH-003 & Żadne dane uwierzytelniające z~\texttt{web.config} nie trafiają do commitowanego pliku & \texttt{error} & -- & P5 \\
SK-ARCH-004 & \texttt{Program.cs} jako manifest; okablowanie z~\texttt{Global.asax} i~\texttt{App\_Start} w~metodach rozszerzających & \texttt{suggestion} & -- & P9 \\
SK-ARCH-005 & Żadnego nowego \texttt{DbContext} w~kontrolerze; przeniesione przypadki zapisane & \texttt{warning} & \path{PORTCULLIS_P9_CONTROLLER_NO_DBCONTEXT} & P9 \\
SK-ARCH-006 & Punkty rozszerzeń jako interfejsy w~DI, nie nowe klasy bazowe & \texttt{suggestion} & \path{PORTCULLIS_P10_CUSTOM_BASE_CLASS} & P10 \\
SK-ARCH-007 & Istniejące service defaults w~nowym hoście; żadnych nowych w~trakcie migracji & \texttt{suggestion} & \path{PORTCULLIS_P15_MISSING_SERVICE_DEFAULTS} & P15 \\
SK-MIG-001 & Typy hostingu System.Web zastąpione przez ASP.NET Core & \texttt{error} & \path{PORTCULLIS_MIG_SYSTEM_WEB} & -- \\
SK-MIG-002 & Dane żądania przekazywane jawnie zamiast \texttt{HttpContext.Current} & \texttt{error} & \path{PORTCULLIS_MIG_HTTPCONTEXT_CURRENT} & -- \\
SK-MIG-003 & Asynchroniczność do końca łańcucha; nigdy blokowanie na zadaniu & \texttt{error} & \path{PORTCULLIS_MIG_SYNC_OVER_ASYNC} & -- \\
SK-MIG-004 & Ustawienia przez \texttt{IOptions<T>} lub \texttt{IConfiguration}, nigdy \texttt{ConfigurationManager} & \texttt{error} & \path{PORTCULLIS_MIG_CONFIGURATION_MANAGER} & P5 \\
SK-MIG-005 & Jawne \texttt{StringComparison} w~każdym porównaniu i~sortowaniu zależnym od kultury & \texttt{warning} & \path{CA1310} & -- \\
SK-MIG-006 & Nazwana kultura przy formatowaniu, parsowaniu i~zmianie wielkości liter; źródło kultury legacy zachowane & \texttt{warning} & \path{CA1304}, \path{CA1305}, \path{CA1311} & -- \\
SK-MIG-007 & Bez zmiany trybu globalizacji; flaga dla każdego wyniku, który może zmienić NLS → ICU & \texttt{error} & -- & -- \\
SK-MIG-008 & Zależności przez konstruktor, nie statyczne singletony ani service locator & \texttt{warning} & -- & P10 \\
SK-MIG-009 & Wychodzący HTTP przez \texttt{IHttpClientFactory}, nigdy klient per wywołanie & \texttt{warning} & -- & P2 \\
SK-MIG-010 & Zachowanie zapytań, ładowania i~walidacji przy odejściu od EF6 & \texttt{error} & -- & P4 \\
SK-MIG-011 & Zachowaj zachowanie; flaguj podejrzewany defekt zamiast go naprawiać & \texttt{error} & -- & P14 \\
SK-MIG-012 & Kontrakt HTTP bez zmian: trasy, kody statusu, nagłówki, ciasteczka, kształt JSON & \texttt{error} & -- & -- \\
SK-MIG-013 & Istniejące testy i~ich oczekiwania zostają; nieprzechodzący test jest znaleziskiem & \texttt{error} & -- & P13 \\
\bottomrule
\end{xltabular}
}

Rozkład wag: jedenaście reguł \repopath{error}, sześć \repopath{warning}, trzy
\repopath{suggestion}. Jedenaście reguł ma diagnostykę; razem daje to trzynaście
identyfikatorów diagnostyk w~wygenerowanej konfiguracji (SK-MIG-006 mapuje trzy analyzery).
Kategorie i~wzorce \repopath{appliesTo} podaje opis każdej reguły.

\zrodla{\path{letsgolegacy.secondkey-standards/generated/manifest.json},
\path{letsgolegacy.secondkey-standards/generated/skills/applying-dotnet-migration-standards/SKILL.md},
\path{letsgolegacy.secondkey-standards/README.md},
\path{letsgolegacy.secondkey-standards/standards/README.md}}

\subsection{Standardy architektoniczne SK-ARCH-001--007}\label{sec:standardy-arch}

Każda reguła \repopath{SK-ARCH-} powtarza jedną zasadę konstytucji architektury dla przypadku
migracji i~linkuje jej kotwicę w~\path{00-REFERENCE-ARCHITECTURE.md}; konstytucja pozostaje
źródłem rozstrzygającym. Opis każdej reguły ma cztery części: regułę (stwierdzenie
normatywne), powód („dlaczego”), sposób spełnienia (sekcja \emph{Migration} reguły) oraz
kontrolę zgodności z~przypadkami, które agent flaguje zamiast naprawiać. Zdanie „Bramka
zgłasza … jako …” w~opisie reguły podaje wagę, którą ustawia konfiguracja standardów
w~buildzie; skan CLI Portcullis raportuje poziomy domyślne
(sekcja~\ref{sec:standardy-wagi-kanaly}).

\subsubsection{SK-ARCH-001 -- istniejący schemat jest kontraktem}\label{sec:standardy-sk-arch-001}

Tytuł źródłowy: \emph{Treat the existing schema as the contract; migrate it, never
EnsureCreated it}. Waga \repopath{error}, kategoria \repopath{data-access},
\repopath{appliesTo}: \repopath{"*.cs"}, zasada P4.

\paragraph{Reguła.} Baza danych, na której działa system legacy, jest częścią kontraktu.
Zmigrowany system nigdy nie wywołuje \repopath{Database.EnsureCreated()} poza testami
i~nigdy nie pozwala, by konwencja modelu zmieniła schemat. Model EF Core jest mapowany na
istniejące tabele, pierwsza migracja EF Core opisuje istniejący schemat, a~późniejsze zmiany
schematu wchodzą jako przeglądane migracje -- przez \repopath{MigrateAsync} w~hosted service
albo jako skrypty. Inicjalizatorów bazy EF6 (\repopath{CreateDatabaseIfNotExists},
\path{DropCreateDatabaseIfModelChanges}, \path{MigrateDatabaseToLatestVersion}) nie
odtwarza się przez \repopath{EnsureCreated}.

\paragraph{Dlaczego.} \repopath{EnsureCreated} tworzy schemat z~bieżącego modelu, jeśli baza
nie istnieje, i~nic nie robi, jeśli istnieje: nie ma historii migracji ani ścieżki
aktualizacji, więc schemat zamarza przy pierwszym uruchomieniu. P4 dopuszcza je tylko dla
ścieżki in-memory i~testów -- po tym, jak system referencyjny użył go na prawdziwej bazie
i~zamroził swój schemat produkcyjny. W~migracji zawodzi jeszcze ciszej: skierowane na kopię
bazy legacy nie robi nic, a~niezgodność nowego modelu ze starym schematem wychodzi później
jako nieudane zapytanie.

\paragraph{Jak spełnić.} Pierwszą migrację EF Core generuje się z~modelu zmapowanego na
istniejący schemat i~sprawdza względem tego schematu: musi opisywać tabele takimi, jakie są.
W~istniejących bazach tę migrację oznacza się jako zastosowaną (wiersz bazowy
w~\path{__EFMigrationsHistory}) zamiast ją uruchamiać; nowe środowiska uruchamiają ją
normalnie. Testy korzystające z~providera in-memory albo z~jednorazowej bazy mogą zachować
\repopath{EnsureCreated}. Przykład zgodny z~reguły: klasa \repopath{MigrationHostedService}
(\repopath{BackgroundService}) wywołuje \path{db.Database.MigrateAsync(stoppingToken)} po
starcie aplikacji, tak aby sondy zdrowia odpowiadały.

\paragraph{Zgodność i~flagi.} Bramka zgłasza \path{PORTCULLIS_P4_ENSURE_CREATED_OUTSIDE_TEST}
jako błąd. Agent flaguje zamiast naprawiać: różnicę między zmapowanym modelem a~istniejącym
schematem, której nie da się usunąć mapowaniem (zamiast generować migrację zmieniającą
schemat), oraz inicjalizator EF6, który usuwał lub odtwarzał bazę, jeśli jakieś środowisko
na tym polegało.

\zrodla{\path{letsgolegacy.secondkey-standards/standards/SK-ARCH-001-migrate-never-ensure-schema.md}}

\subsubsection{SK-ARCH-002 -- dane seedujące poza modelem}\label{sec:standardy-sk-arch-002}

Tytuł źródłowy: \emph{Move EF6 seed data to a versioned seeding step, not HasData in the
model}. Waga \repopath{warning}, kategoria \repopath{data-access}, \repopath{appliesTo}:
\repopath{"*.cs"}, zasada P4.

\paragraph{Reguła.} Dane referencyjne, które seedował EF6 -- w~metodzie \repopath{Seed}
inicjalizatora albo w~nadpisaniu \repopath{Seed} konfiguracji migracji -- przechodzą do
wersjonowanego skryptu seedującego albo do idempotentnej usługi seedującej, uruchamianej raz
na środowisko. Nie przechodzą do wywołań \repopath{HasData} w~\repopath{OnModelCreating}.

\paragraph{Dlaczego.} \repopath{HasData} osadza cały zbiór danych w~migawce modelu każdej
migracji. Konstytucja zapisuje koszt: taksonomia kategorii seedowana w~ten sposób dała pliki
designera migracji po około 18\,000 wierszy każdy, co uczyniło każdą zmianę schematu
nieprzeglądalną -- i~zmiany schematu wchodziły bez przeglądu. Migracje opisują
schemat; dane referencyjne są danymi i~mają własny, wersjonowany krok. Przeniesienie musi też
zachować semantykę seedowania legacy: \repopath{Seed} migracji EF6 uruchamiał się przy każdym
zastosowaniu migracji i~zwykle używał \repopath{AddOrUpdate}, więc wstawiał brakujące wiersze
i~nadpisywał edytowane. To, na którym z~tych efektów system polegał, jest zachowaniem.

\paragraph{Jak spełnić.} Wartości seedujące przenosi się bez zmian, z~tymi samymi kluczami,
które mają dane legacy. Z~kodu \repopath{Seed} legacy ustala się, czy tylko wstawiał, czy
także nadpisywał: \repopath{AddOrUpdate}, który przy każdym wdrożeniu resetował edycje
operatorów, jest zachowaniem do zachowania albo do oflagowania. Przykład zgodny:
\repopath{CountrySeeder.SeedAsync} dodaje tylko brakujące kraje (sprawdzenie po
\repopath{TwoLetterIsoCode}) i~jest uruchamiany raz na środowisko przez wdrożenie.

\paragraph{Zgodność i~flagi.} Bramka zgłasza \path{PORTCULLIS_P4_SEED_DATA_IN_MODEL} jako
ostrzeżenie. Flagowany jest \repopath{Seed} legacy, który nadpisywał wiersze edytowane przez
operatorów albo seedował dane testowe do produkcji.

\zrodla{\path{letsgolegacy.secondkey-standards/standards/SK-ARCH-002-seed-data-outside-model.md}}

\subsubsection{SK-ARCH-003 -- sekrety nie trafiają do commitowanych plików}\label{sec:standardy-sk-arch-003}

Tytuł źródłowy: \emph{Move no credential from web.config into a committed file}. Waga
\repopath{error}, kategoria \repopath{configuration}, \repopath{appliesTo}:
\repopath{"*.json"}, \repopath{"*.config"}, \repopath{"*.cs"}, zasada P5.

\paragraph{Reguła.} Żadne dane uwierzytelniające z~\repopath{web.config} -- hasło
w~connection stringu, klucz API, hasło SMTP, machine key -- nie przechodzą do pliku
w~repozytorium: \repopath{appsettings.json}, \path{appsettings.<Environment>.json},
\repopath{launchSettings.json}, kodu źródłowego ani komentarza. Sekrety pochodzą ze
środowiska: \repopath{dotnet user-secrets} na maszynie dewelopera, magazyn sekretów platformy
w~każdym wdrożonym środowisku. Commitowany plik zachowuje klucz z~pustą wartością.

\paragraph{Dlaczego.} \repopath{web.config} w~produkcji bywał trzymany poza repozytorium,
szyfrowany albo przepisywany transformacją wdrożeniową. \repopath{appsettings.json} jest
commitowany, kopiowany do każdego wyniku builda i~obrazu kontenera i~pozostaje na zawsze
w~historii. Migracja, która mechanicznie kopiuje \repopath{<connectionStrings>} do tego
pliku, publikuje hasło każdemu, kto może czytać repozytorium. P5 mówi wprost, że żaden
sekret nie jest literałem w~commitowanym pliku konfiguracyjnym i~że egzekwuje to skaner
sekretów w~CI i~w~hooku pre-commit, a~nie przegląd.

\paragraph{Jak spełnić.} Dla każdego wpisu \repopath{<connectionStrings>} i~każdej tajnej
wartości \repopath{<appSettings>} tworzy się klucz w~\repopath{appsettings.json} z~pustą
wartością i~dokumentuje, skąd przychodzi wartość. Nazwy kluczy muszą pozostać stabilne:
zmienna środowiskowa ustawiana przez operatora jest z~nich wyprowadzana
(\repopath{ConnectionStrings\_\_Shop}). Jeśli repozytorium docelowe nie ma skanera sekretów,
jest to flaga -- P5 traktuje takie repozytorium jako zawierające dane uwierzytelniające,
dopóki nie udowodni się inaczej. Przykład zgodny z~reguły:

\begin{lstlisting}
# A developer machine: the value lives outside the repository
dotnet user-secrets set "ConnectionStrings:Shop" "<connection string>"

# A deployed environment: an environment variable filled from the secret store
ConnectionStrings__Shop=<from the secret store>
\end{lstlisting}

\paragraph{Zgodność i~flagi.} Reguła nie ma diagnostyki bramki; jest tylko instrukcją dla
agenta. Flagowane są dane uwierzytelniające, które już są w~historii repozytorium legacy:
wymagają rotacji, która jest decyzją i~pracą człowieka, a~usunięcie ich z~pliku nie cofa
ujawnienia.

\zrodla{\path{letsgolegacy.secondkey-standards/standards/SK-ARCH-003-no-secrets-in-committed-config.md}}

\subsubsection{SK-ARCH-004 -- Program.cs jako manifest}\label{sec:standardy-sk-arch-004}

Tytuł źródłowy: \emph{Keep Program.cs a manifest; move Global.asax and App\_Start wiring
into extension methods}. Waga \repopath{suggestion}, kategoria \repopath{layering},
\repopath{appliesTo}: \repopath{"*.cs"}, zasada P9.

\paragraph{Reguła.} Nowy \repopath{Program.cs} czyta się jak lista możliwości aplikacji.
Okablowanie, które mieszkało w~\repopath{Global.asax} (\repopath{Application\_Start}),
w~\path{App_Start/*Config.cs} (trasy, filtry, bundle), w~klasie \repopath{Startup} albo
w~module kontenera, przechodzi do metod rozszerzających należących do funkcji, którą
okablowuje -- \path{builder.Services.AddCatalog()}, \path{app.MapShopRoutes()} --
z~zachowaniem kolejności, w~jakiej wykonywały się rejestracje i~middleware.

\paragraph{Dlaczego.} Migracja wytwarza całe okablowanie aplikacji naraz. Wstawione inline
robi z~\repopath{Program.cs} plik na kilkaset wierszy, przed którym ostrzega P9 -- przykład
w~konstytucji porównuje manifest na 130 wierszy z~odpowiednikiem inline na 399 wierszy,
trudniejszym do zmiany. Nazwane metody rozszerzające nic nie kosztują, gdy okablowanie i~tak
jest przepisywane, a~jedyną rzecz, która jest tu zachowaniem -- kolejność -- czynią widoczną
w~kilku wierszach.

\paragraph{Jak spełnić.} Jedna metoda rozszerzająca na funkcję albo na klasę konfiguracji
legacy, w~projekcie, który jest jej właścicielem. Kolejność rejestracji zachowuje się tam,
gdzie ma znaczenie (ostatnia rejestracja usługi wygrywa), a~kolejność middleware --
dokładnie. W~przykładzie zgodnym \repopath{app.UseShopPipeline()} ustawia middleware
w~kolejności zdarzeń \repopath{Global.asax}, a~\path{app.MapShopRoutes()} odtwarza
\repopath{RouteConfig.RegisterRoutes} z~tymi samymi szablonami w~tej samej kolejności.

\paragraph{Zgodność i~flagi.} Reguła nie ma diagnostyki; jako \repopath{suggestion} jest
stosowana tam, gdzie zmiana jest mechaniczna. Flagowane jest okablowanie legacy, którego
kolejności nie da się odtworzyć, na przykład moduł polegający na kolejności, w~jakiej
kontener skanował assembly.

\zrodla{\path{letsgolegacy.secondkey-standards/standards/SK-ARCH-004-program-cs-is-a-manifest.md}}

\subsubsection{SK-ARCH-005 -- kontrolery pozostają warstwą transportu}\label{sec:standardy-sk-arch-005}

Tytuł źródłowy: \emph{Introduce no DbContext into a controller; record the ones carried
over}. Waga \repopath{warning}, kategoria \repopath{layering}, \repopath{appliesTo}:
\repopath{"*.cs"}, zasada P9.

\paragraph{Reguła.} Kontrolery pozostają transportem: wiążą dane, autoryzują i~delegują.
Migracja nie wprowadza \repopath{DbContext} do kontrolera, który go nie używał. Tam, gdzie
kontroler legacy już korzystał bezpośrednio z~kontekstu danych, migracja zostawia dostęp do
danych tam, gdzie był -- przeniesienie go do usługi to refaktoryzacja na osobną zmianę --
i~zapisuje wynik bramki jako przeniesiony dług.

\paragraph{Dlaczego.} P9 układa usługę warstwami: kontrolery nad orkiestratorami nad
usługami domenowymi nad dostępem do danych, a~kontroler trzymający \repopath{DbContext}
zwija te warstwy do jednej klasy. Migrację kusi dokładnie to -- wstrzyknięcie kontekstu jest
najkrótszą drogą od \repopath{new ShopContext()} -- i~równie mocno kusi ją „naprawienie”
warstw przy okazji. Oba ruchy są błędne z~tego samego powodu: pierwszy dodaje dług, drugi
miesza refaktoryzację ze zmianą, której całym twierdzeniem jest to, że zachowanie się nie
przesunęło. Bramka zgłasza wzorzec na zmienionych wierszach, więc przeniesione przypadki
wychodzą w~przeglądzie; zapisanie ich utrzymuje dług widocznym bez wykonywania refaktoryzacji
teraz.

\paragraph{Jak spełnić.} Kolaboratorzy każdego kontrolera zostają tacy, jak byli: kontroler,
który wywoływał usługę, wywołuje tę samą usługę. Kontroler legacy, który tworzył kontekst
przez \repopath{new}, dostaje go wstrzykniętego, z~tym samym czasem życia na żądanie, oraz
flagę, np.:

\begin{lstlisting}
// SECONDKEY-FLAG SK-ARCH-005: controller reads ShopDbContext directly (carried over, not refactored)
\end{lstlisting}

\paragraph{Zgodność i~flagi.} Bramka zgłasza \path{PORTCULLIS_P9_CONTROLLER_NO_DBCONTEXT}
jako ostrzeżenie. Flagowany jest każdy kontroler, który po migracji nadal trzyma
\repopath{DbContext}.

\zrodla{\path{letsgolegacy.secondkey-standards/standards/SK-ARCH-005-controllers-stay-transport.md}}

\subsubsection{SK-ARCH-006 -- rozszerzanie przez interfejsy, nie klasy bazowe}\label{sec:standardy-sk-arch-006}

Tytuł źródłowy: \emph{Add extension points as interfaces registered in DI, not new base
classes}. Waga \repopath{suggestion}, kategoria \repopath{dependency-injection},
\repopath{appliesTo}: \repopath{"*.cs"}, zasada P10.

\paragraph{Reguła.} Wspólne zachowanie, które migracja musi dodać -- mechanika, którą każdy
kontroler dostawał wcześniej od frameworka, hook wywoływany przez kilka usług -- jest
interfejsem albo filtrem zarejestrowanym w~DI, a~nie nową klasą bazową. Klasy bazowe, które
kod legacy już ma, zostają bez zmian: przekazywane przez nie atrybuty, filtry i~nadpisania są
zachowaniem.

\paragraph{Dlaczego.} P10 zastępuje „klasę bazową biblioteki rdzenia” interfejsem i~jedną
rejestracją: nowy algorytm, dostawca albo krok to klasa i~jedna linia DI, bez frameworka do
spełnienia. Migracje łatwo wymyślają klasy bazowe -- miejsce na usługi, które wcześniej
pochodziły ze statyka, na helpery, które żyły w~\repopath{HttpContext} -- a~każda z~nich
staje się hierarchią, w~którą musi się wpasować każdy późniejszy kontroler. Odwrotny błąd
jest równie kosztowny: spłaszczenie bazowego kontrolera legacy zmienia, które filtry
i~atrybuty obowiązują dla każdej akcji, która je dziedziczyła.

\paragraph{Jak spełnić.} Klasy bazowe legacy i~ich atrybuty przenosi się tak, jak są. Nowe
wspólne zachowanie trafia do filtrów, middleware albo małych wstrzykiwanych usług; przykład
zgodny rejestruje filtr raz
(\path{options.Filters.Add<StoreClosedFilter>()}), a~kontrolery deklarują w~konstruktorze
to, czego używają.

\paragraph{Zgodność i~flagi.} Bramka zgłasza \path{PORTCULLIS_P10_CUSTOM_BASE_CLASS} jako
sugestię. Flagowana jest klasa bazowa legacy, której zachowania ASP.NET Core nie odtworzy
przez dziedziczenie, na przykład kontroler bazowy polegający na kolejności
\repopath{OnActionExecuting}, która się zmieniła.

\zrodla{\path{letsgolegacy.secondkey-standards/standards/SK-ARCH-006-extend-through-interfaces.md}}

\subsubsection{SK-ARCH-007 -- okablowanie obserwowalności}\label{sec:standardy-sk-arch-007}

Tytuł źródłowy: \emph{Wire the solution's shared service defaults into the new host; invent
none mid-migration}. Waga \repopath{suggestion}, kategoria \repopath{observability},
\repopath{appliesTo}: \repopath{"*.cs"}, zasada P15.

\paragraph{Reguła.} Jeśli rozwiązanie ma już wspólne ustawienia usług -- metodę
\repopath{AddServiceDefaults()}, która okablowuje OpenTelemetry, health checki i~service
discovery -- zmigrowany host ją wywołuje, więc ślady, metryki i~endpointy zdrowia istnieją od
pierwszego wdrożenia. Jeśli rozwiązanie jej nie ma, migracja jej nie tworzy: zapisuje lukę
na zmianę uzupełniającą.

\paragraph{Dlaczego.} P15 czyni obserwowalność decyzją podejmowaną przy budowie, a~nie
dodatkiem na koniec; dla zmigrowanego systemu telemetria jest też sposobem, w~jaki operator
porównuje stary i~nowy system, gdy oba działają w~produkcji. Nowy \repopath{Program.cs},
który buduje hosta bez ustawień rozwiązania, daje usługę niewidoczną wtedy, gdy źle działa.
Tworzenie projektu service defaults w~trakcie migracji jest błędem odwrotnym: nowa
infrastruktura, nowe zależności i~nowe endpointy w~zmianie, której twierdzeniem jest, że
przesunęła się tylko platforma.

\paragraph{Jak spełnić.} Każdy zmigrowany host wywołuje istniejące ustawienia rozwiązania:
\path{builder.AddServiceDefaults()} oraz \path{app.MapDefaultEndpoints()}
(\repopath{/health} i~\repopath{/alive}). Cele i~poziomy logowania legacy przenosi się bez
zmian; przejście na nowy pipeline logowania to osobna zmiana.

\paragraph{Zgodność i~flagi.} Bramka zgłasza \path{PORTCULLIS_P15_MISSING_SERVICE_DEFAULTS}
jako sugestię. Flagowany jest zmigrowany host bez okablowania obserwowalności, gdy
rozwiązanie nie ma wspólnych ustawień -- jako zmiana uzupełniająca, której wymaga P15
(i~P2a: „every service calls AddServiceDefaults()”).

\zrodla{\path{letsgolegacy.secondkey-standards/standards/SK-ARCH-007-wire-observability.md}}

\subsection{Standardy migracyjne SK-MIG-001--013}\label{sec:standardy-mig}

Reguły \repopath{SK-MIG-} opisują pułapki przenoszenia systemu z~.NET Framework do
nowoczesnego .NET: System.Web, stan ambientowy, sync-over-async, konfigurację, przejście
NLS → ICU, wstrzykiwanie zależności, wychodzący HTTP, EF6 i~przede wszystkim zachowanie
zachowania systemu. Opis każdej reguły ma ten sam układ co w~sekcji~\ref{sec:standardy-arch}.

\subsubsection{SK-MIG-001 -- System.Web zastąpiony przez ASP.NET Core}\label{sec:standardy-sk-mig-001}

Tytuł źródłowy: \emph{Replace System.Web hosting types with ASP.NET Core}. Waga
\repopath{error}, kategoria \repopath{hosting}, \repopath{appliesTo}: \repopath{"*.cs"}.

{\raggedright
\paragraph{Reguła.} Zmigrowany kod nie zależy od modelu hostingu ASP.NET Framework: nie ma
referencji do assembly \repopath{System.Web} ani użycia typów \repopath{System.Web.*}, które
istnieją po to, by działać w~potoku ASP.NET w~IIS -- \repopath{HttpApplication},
\repopath{HttpContext} i~\repopath{HttpContextBase}, \repopath{HttpRequest},
\repopath{HttpResponse}, \repopath{IHttpModule}, \repopath{IHttpHandler},
\repopath{System.Web.Mvc}, \repopath{System.Web.Http}, \repopath{System.Web.Optimization},
\repopath{System.Web.Security}, \repopath{System.Web.SessionState}. Żądania obsługuje
ASP.NET Core: middleware, kontrolery albo endpointy oraz
\path{Microsoft.AspNetCore.Http.HttpContext}.\par}

\paragraph{Dlaczego.} \repopath{System.Web.dll} nie jest częścią nowoczesnego .NET, więc
projekt z~\repopath{net10.0} nie może go referencjonować. Kod, który nadal nazywa te typy,
albo przestaje się kompilować, albo kompiluje się z~warstwą zgodności
\path{Microsoft.AspNetCore.SystemWebAdapters}, która odtwarza część powierzchni
\repopath{System.Web} na ASP.NET Core. Adaptery są uprawnionym mostem przy migracji
przyrostowej, prowadzonej obok starego systemu, ale różnią się od oryginału w~udokumentowanych
szczegółach -- buforowanie żądania, buforowanie odpowiedzi i~stan sesji są włączane per
endpoint, a~\repopath{HttpContext.Current} nie jest już związany z~wątkiem -- więc każde
wywołanie przez nie jest miejscem, w~którym zachowanie legacy i~zmigrowane mogą się rozejść.
Migracja jest skończona dopiero wtedy, gdy model hostingu jest własnym modelem ASP.NET Core,
a~bramka raportuje to, co zostało.

Poza regułą jest jeden typ: \repopath{System.Web.HttpUtility} jest dostarczany
z~nowoczesnym .NET, a~jego enkodery dają inne bajty niż alternatywy. Przykłady z~reguły
(pierwszy niezgodny, drugi zgodny):

\begin{lstlisting}
// "Modernising" an encoder during the migration. HttpUtility produced "%2fcart%3fid%3d7";
// WebUtility produces "%2Fcart%3Fid%3D7" — different bytes in every redirect URL.
var returnUrl = WebUtility.UrlEncode("/cart?id=7");
\end{lstlisting}

\begin{lstlisting}
// System.Web.HttpUtility is part of modern .NET, so the output stays byte-identical
var returnUrl = System.Web.HttpUtility.UrlEncode("/cart?id=7"); // "%2fcart%3fid%3d7"
\end{lstlisting}

\paragraph{Jak spełnić.} Każde zdarzenie potoku (\path{Application_BeginRequest},
\path{Application_EndRequest}, handlery \repopath{IHttpModule}) staje się middleware,
rejestrowanym w~kolejności, w~jakiej zdarzenia się wykonywały; każdy \repopath{IHttpHandler}
i~plik \repopath{.ashx} -- endpointem. Kontrolery MVC i~Web API przechodzą na kontrolery
ASP.NET Core z~trasami i~nazwami akcji dokładnie takimi, jak były (SK-MIG-012). Jeśli
adaptery System.Web służą za most, każdy obszar od nich zależny jest zapisany jako flaga
zachowania, z~krokiem, który go usunie. Wywołania \repopath{HttpUtility} zostają bez zmian;
jeśli bramka zgłosi takie wywołanie, wiersz jest tłumiony z~uzasadnieniem cytującym regułę,
zamiast zmiany enkodera.

Przykład z~reguły dla zdarzenia potoku: niezgodny \path{Global.asax.cs} zawiera klasę
\texttt{ShopApplication : System.Web.HttpApplication}, której
\path{Application_BeginRequest} dodaje nagłówek \repopath{X-Shop-Node} przez
\repopath{HttpContext.Current}; zgodny ustawia ten sam nagłówek w~middleware, zanim zacznie
się odpowiedź:

\begin{lstlisting}
// Program.cs: the same header, set by middleware before the response starts
app.Use(async (context, next) =>
{
    context.Response.Headers["X-Shop-Node"] = Environment.MachineName;
    await next(context);
});
\end{lstlisting}

Portcullis domyślnie nie zgłasza \path{System.Web.HttpUtility} ani
\path{System.Web.IHtmlString} (lista \repopath{systemWebAllowedTypes},
sekcja~\ref{sec:portcullis-migracja}); ustawienie tej listy zastępuje wartość domyślną,
więc bramka zgłosi \repopath{HttpUtility} dopiero wtedy, gdy konfiguracja repozytorium poda
listę bez niego -- i~wtedy obowiązuje tłumienie z~uzasadnieniem cytującym regułę.

\paragraph{Zgodność i~flagi.} Bramka zgłasza \path{PORTCULLIS_MIG_SYSTEM_WEB} jako błąd.
Flagowane są: walidacja żądań -- ASP.NET Framework odrzucał żądania zawierające znaczniki
(„A potentially dangerous Request.Form value was detected”), ASP.NET Core nie ma
odpowiednika, więc zmigrowany system przyjmuje wejście, które legacy odrzucał; to się
zapisuje, bez doraźnego filtrowania imitującego dawną walidację -- oraz każde zachowanie
kodowania, kolejności potoku albo buforowania, którego zmigrowany kod nie odtwarza dokładnie.

\zrodla{\path{letsgolegacy.secondkey-standards/standards/SK-MIG-001-replace-system-web.md},
\path{letsgolegacy.portcullis/docs/CONFIGURATION.md},
\path{letsgolegacy.portcullis/docs/rules/MIGRATION.md}}

\subsubsection{SK-MIG-002 -- bez HttpContext.Current}\label{sec:standardy-sk-mig-002}

Tytuł źródłowy: \emph{Pass request data explicitly instead of reading HttpContext.Current}.
Waga \repopath{error}, kategoria \repopath{hosting}, \repopath{appliesTo}:
\repopath{"*.cs"}.

\paragraph{Reguła.} Żaden kod nie czyta ambientowego \repopath{HttpContext.Current}.
Kontrolery i~endpointy używają \repopath{HttpContext}, który dostają; usługi otrzymują
potrzebne wartości -- identyfikator użytkownika, host sklepu, kulturę -- jako parametry.
\repopath{IHttpContextAccessor} wstrzykuje się tylko do infrastruktury brzegowej, która nie
ma innej drogi do żądania, np. do enrichera logowania, i~nic nie trzyma referencji do
\repopath{HttpContext} dłużej niż żądanie, do którego kontekst należy.

\paragraph{Dlaczego.} \repopath{HttpContext.Current} jest statyczną, ambientową
zależnością: każda klasa, jakkolwiek głęboko, może sięgnąć do żądania, co ukrywa zależność
przed czytelnikiem i~uniemożliwia test bez serwera WWW. ASP.NET Core usuwa ten statyk.
Mechaniczny zamiennik -- \repopath{IHttpContextAccessor} wstrzyknięty wszędzie tam, gdzie
był \repopath{HttpContext.Current} -- zachowuje ukryte sprzężenie za interfejsem i~dodaje
tryb awarii: \repopath{HttpContext} nie jest bezpieczny wątkowo, a~odczyt poza żądaniem,
które go posiada (zadanie w~tle, kontynuacja wykonana po zakończeniu odpowiedzi, metoda
\repopath{async void}), zwraca dane innego żądania albo rzuca
\repopath{NullReferenceException}. Kod legacy miał te same zagrożenia na innych zasadach,
bo ASP.NET Framework wiązał kontekst z~wątkiem, więc migracja jest momentem, w~którym
przepływ danych ma stać się jawny.

\paragraph{Jak spełnić.} W~kontrolerach \repopath{HttpContext.Current} zastępują
właściwości kontrolera \repopath{HttpContext}, \repopath{Request}, \repopath{Response}
i~\repopath{User}. W~usługach dodaje się jako parametr wartość, której usługa faktycznie
używa, i~odczytuje ją raz, na brzegu; gdy wiele usług potrzebuje tej samej wartości,
preferowana jest wąska abstrakcja (\repopath{ICurrentCustomer}) zamiast
\repopath{IHttpContextAccessor}. \repopath{AddHttpContextAccessor()} rejestruje się tylko dla
infrastruktury, która tego potrzebuje. Przed startem pracy w~tle kopiuje się z~kontekstu
wartości, których ta praca potrzebuje.

Przykłady zgodne z~reguły: usługa dostaje host sklepu jako parametr, kontroler czyta
żądanie raz i~przekazuje w~dół zwykłe wartości, a~praca w~tle dostaje wartość skopiowaną,
zanim żądanie się skończy:

\begin{lstlisting}
public class OrderNumberService
{
    public string NextNumber(string storeHost) => $"{storeHost}-{_sequence.Next()}";
}

// The controller reads the request once and passes plain values down
public IActionResult PlaceOrder(OrderForm form)
{
    var number = _orderNumbers.NextNumber(Request.Host.Host);
    // ...
}
\end{lstlisting}

\begin{lstlisting}
// Copy what background work needs while the request is still in scope
public IActionResult Checkout()
{
    var userName = User.Identity?.Name;
    _queue.Enqueue(() => _audit.Write(userName, "checkout"));
    return RedirectToAction("Complete");
}
\end{lstlisting}

Przykłady niezgodne: \repopath{OrderNumberService} czytający
\path{HttpContext.Current.Request.Url.Host} w~usłudze domenowej oraz mechaniczny zamiennik
\repopath{AuditService}, który przez \repopath{IHttpContextAccessor} czyta użytkownika
w~\repopath{Task.Run}, już poza żądaniem.

\paragraph{Zgodność i~flagi.} Bramka zgłasza \path{PORTCULLIS_MIG_HTTPCONTEXT_CURRENT} jako
błąd. Flagowane są: wartość, której znaczenie zmienia się między modelami hostingu -- nazwa
anonimowego użytkownika była pustym napisem pod domyślnym modułem uwierzytelniania ASP.NET
Framework, a~pod ASP.NET Core jest \repopath{null}, więc kod polegający na pustym napisie
rzuca wyjątek albo rozgałęzia się inaczej -- oraz kod legacy, który czytał
\repopath{HttpContext.Current} z~wątku w~tle i~polegał na tym, że jest tam \repopath{null}.

\zrodla{\path{letsgolegacy.secondkey-standards/standards/SK-MIG-002-no-httpcontext-current.md}}

\subsubsection{SK-MIG-003 -- asynchroniczność do końca łańcucha}\label{sec:standardy-sk-mig-003}

Tytuł źródłowy: \emph{Keep asynchronous code asynchronous all the way; never block on a
task}. Waga \repopath{error}, kategoria \repopath{concurrency}, \repopath{appliesTo}:
\repopath{"*.cs"}.

\paragraph{Reguła.} Żaden kod nie blokuje się na pracy asynchronicznej przez
\repopath{.Result}, \repopath{.Wait()}, \repopath{.GetAwaiter().GetResult()},
\repopath{Task.WaitAll} ani \repopath{Task.WaitAny}. Metoda, która wywołuje asynchroniczne
API, sama jest asynchroniczna -- aż do akcji kontrolera albo endpointu, które rozpoczęły
żądanie.

\paragraph{Dlaczego.} Kontekst synchronizacji ASP.NET Framework robił z~blokowania na
zadaniu zakleszczenie, które tylko czekało na okazję. ASP.NET Core nie ma takiego kontekstu,
więc zakleszczenie zwykle znika -- i~właśnie dlatego sync-over-async przechodzi przez
migrację niezauważony. Nadal jednak trzyma wątek z~puli przez całe I/O, na które czeka, a~pod
obciążeniem pula się wyczerpuje: żądania czekają w~kolejce, opóźnienia rosną, health checki
przekraczają limit czasu -- bez żadnego błędu w~logu. Kestrel z~tego samego powodu domyślnie
odmawia synchronicznego I/O na ciele żądania i~odpowiedzi (\repopath{AllowSynchronousIO}
ma wartość \repopath{false}). Blokowanie zmienia też postać wyjątków: \repopath{.Result}
i~\repopath{.Wait()} opakowują błąd w~\repopath{AggregateException}, a~\repopath{await}
rzuca oryginalny wyjątek. Zamiana jednego na drugie zmienia, który blok \repopath{catch}
się wykona, więc jest zmianą zachowania do sprawdzenia, a~nie mechaniczną edycją.

\paragraph{Jak spełnić.} \repopath{async} i~\repopath{await} propaguje się w~górę łańcucha
wywołań aż do akcji. Nowe metody asynchroniczne dostają sufiks \repopath{Async}; zmiana nazwy
publicznego API używanego poza rozwiązaniem jest zmianą łamiącą i~jest flagowana. Tam,
gdzie łańcuch legacy jest synchroniczny od początku do końca, a~wywoływane API nadal ma
przeciążenie synchroniczne, pozostawienie wersji synchronicznej jest dozwolone; konwersja
jest preferowana, gdy migracja i~tak dotyka każdej metody łańcucha. Każdy
\texttt{catch (AggregateException)}, który rozpakowywał zablokowane zadanie, trzeba
przejrzeć: po konwersji wyjątek wewnętrzny przychodzi nieopakowany, więc handler przepisuje
się tak, by te same błędy trafiały do tej samej obsługi. Kodu synchronicznego nie opakowuje
się w~\repopath{Task.Run} wewnątrz żądania, żeby wyglądał na asynchroniczny, a~\repopath{async
void} służy tylko handlerom zdarzeń.

\paragraph{Zgodność i~flagi.} Bramka zgłasza \path{PORTCULLIS_MIG_SYNC_OVER_ASYNC} jako
błąd. Flagowane są: łańcuch, który nie może stać się asynchroniczny bez zmiany sygnatury
używanej poza rozwiązaniem (most zostaje, stłumiony w~tym wierszu z~uzasadnieniem cytującym
regułę); \texttt{AllowSynchronousIO = true} dodane po to, by działał synchroniczny
konsument strumienia; każda zmiana typu wyjątku, który dociera do wywołującego albo do
strony błędu.

\zrodla{\path{letsgolegacy.secondkey-standards/standards/SK-MIG-003-async-all-the-way.md}}

\subsubsection{SK-MIG-004 -- IOptions zamiast ConfigurationManager}\label{sec:standardy-sk-mig-004}

Tytuł źródłowy: \emph{Read settings through IOptions<T> or IConfiguration, never
ConfigurationManager}. Waga \repopath{error}, kategoria \repopath{configuration},
\repopath{appliesTo}: \repopath{"*.cs"}, zasada P5.

\paragraph{Reguła.} Ustawienia są wiązane z~klasami opcji (\repopath{IOptions<T>})
z~\repopath{IConfiguration}, które czyta \repopath{appsettings.json}, zmienne środowiskowe
i~magazyn sekretów platformy. Żaden kod nie czyta
\path{System.Configuration.ConfigurationManager.AppSettings} ani
\path{ConfigurationManager.ConnectionStrings}. Każdy klucz, wartość domyślna i~wartość, które
system legacy miał w~\repopath{web.config} lub \repopath{app.config}, ma jawny odpowiednik,
a~każda wartość znaczy to, co znaczyła wcześniej.

\paragraph{Dlaczego.} Pakiet \path{System.Configuration.ConfigurationManager} sprawia, że
wywołania legacy kompilują się na nowoczesnym .NET, i~właśnie to czyni go niebezpiecznym:
ASP.NET Core nie czyta \repopath{web.config}, więc
\path{ConfigurationManager.AppSettings["Shipping.Surcharge"]} po cichu zwraca
\repopath{null} w~czasie działania, a~kod przechodzi do swojej wartości domyślnej. Build jest
zielony, testy, które stubują konfigurację, przechodzą, a~system działa z~innymi
ustawieniami (w~przykładzie z~reguły dopłata wysyłkowa po cichu staje się zerem). Wiązanie
z~klasami opcji czyni kształt konfiguracji jawnym -- jedna klasa na sekcję, walidowana przy
starcie -- i~jest modelem konfiguracji z~konstytucji: hierarchiczne klucze wiązane z~klasami
opcji, dostarczane przez środowisko (P5).

\paragraph{Jak spełnić.} Inwentaryzuje się każdy klucz \repopath{<appSettings>} i~każdy
wpis \repopath{<connectionStrings>}, łącznie z~wartościami ustawianymi przez każdą
transformację \repopath{web.config} dla każdego środowiska, i~mapuje na klucze konfiguracji;
tabela mapowania trafia do notatek migracji, bo operatorzy ustawiają te wartości po nazwie.
Binder konfiguracji konwertuje wartości kulturą niezmienną (invariant), a~kod legacy często
parsował kulturą serwera (\repopath{decimal.Parse} bez providera): wartości liczbowe i~daty
przenosi się w~formacie niezmiennym i~sprawdza, że każda parsuje się do tej samej wartości.
Używa się \repopath{IOptions<T>}: jest czytany raz, tak jak \repopath{web.config}, którego
edycja restartowała aplikację; \repopath{IOptionsMonitor<T>} i~\repopath{IOptionsSnapshot<T>}
wprowadzają przeładowanie na żywo, czyli zmianę zachowania. Connection stringi czyta się
z~sekcji \repopath{ConnectionStrings} przez \repopath{GetConnectionString}; dane
uwierzytelniające nigdy nie trafiają do commitowanego pliku (SK-ARCH-003). Każdą wartość
domyślną legacy zachowuje się dokładnie; \repopath{ValidateOnStart} stosuje się tylko do
wartości, bez których system legacy nie mógł działać. Przykład zgodny z~reguły:

\begin{lstlisting}
public sealed class ShippingOptions
{
    public const string Section = "Shipping";

    public decimal Surcharge { get; init; }
}

// Registration, next to the service that uses it
builder.Services.AddOptions<ShippingOptions>()
    .Bind(builder.Configuration.GetSection(ShippingOptions.Section))
    .ValidateOnStart();

public class ShippingCalculator(IOptions<ShippingOptions> options)
{
    public decimal Surcharge => options.Value.Surcharge;
}
\end{lstlisting}

Wartość przechodzi z~\repopath{web.config}:

\begin{lstlisting}
<add key="Shipping.Surcharge" value="4.99" />
\end{lstlisting}

do \repopath{appsettings.json}:

\begin{lstlisting}
{
  "Shipping": {
    "Surcharge": 4.99
  }
}
\end{lstlisting}

\paragraph{Zgodność i~flagi.} Bramka zgłasza \path{PORTCULLIS_MIG_CONFIGURATION_MANAGER}
jako błąd. Flagowane są: klucz, który kod legacy czytał, a~którego nie definiował żaden
\repopath{web.config}, albo zdefiniowany, a~nigdy nieczytany; wartość, której wynik
parsowania różni się pod kulturą niezmienną; każde przejście na przeładowanie na żywo.

\zrodla{\path{letsgolegacy.secondkey-standards/standards/SK-MIG-004-options-not-configurationmanager.md}}

\subsubsection{SK-MIG-005 -- jawne StringComparison}\label{sec:standardy-sk-mig-005}

Tytuł źródłowy: \emph{State the StringComparison of every culture-sensitive string
comparison and sort}. Waga \repopath{warning}, kategoria \repopath{globalization},
\repopath{appliesTo}: \repopath{"*.cs"}, analyzer \repopath{CA1310}.

\paragraph{Reguła.} Każde porównanie, wyszukiwanie i~sortowanie napisów, którego domyślne
zachowanie zależy od kultury, deklaruje intencję: \repopath{StringComparison.Ordinal} albo
\repopath{OrdinalIgnoreCase} dla tekstu maszynowego -- identyfikatorów, kluczy, kodów, URL-i,
ścieżek plików, tokenów protokołów -- oraz \path{StringComparison.CurrentCulture} (albo
nazwaną kulturę) dla tekstu pokazywanego ludziom. Dotyczy to \repopath{string.Compare},
\repopath{CompareTo}, \repopath{StartsWith(string)}, \repopath{EndsWith(string)},
\repopath{IndexOf(string)}, \repopath{LastIndexOf(string)} oraz domyślnego comparera, do
którego wracają \repopath{OrderBy}, \repopath{List<string>.Sort}, \repopath{Array.Sort},
\repopath{SortedList} i~\repopath{SortedDictionary}.

\paragraph{Dlaczego.} .NET Framework wykonuje operacje na napisach zależne od kultury przez
NLS (National Language Support) systemu Windows. .NET 5 i~nowsze używają ICU -- na Windows
tak samo jak na Linuksie. Obie biblioteki się różnią:
\begin{itemize}
  \item ICU traktuje \repopath{\textbackslash 0} jako znak o~zerowej wadze, więc
        \path{"Hel\0lo".IndexOf("\0")} zwraca \repopath{0}, a~NLS zwracał \repopath{3};
        \path{"abc".EndsWith("\0")} daje \repopath{true}.
  \item Domyślne sortowanie ICU zachowuje się jak \path{CompareOptions.StringSort}, które
        sortuje interpunkcję przed literami: „bill's” sortuje się przed „bills”, a~lista nazw
        może wrócić w~innej kolejności.
  \item NLS traktuje ligatury jako równe ich literom („oeuf” i~„œuf”), ICU nie.
\end{itemize}
Nic z~tego nie jest błędem kompilacji, a~test jednostkowy zwykle przechodzi na tej
bibliotece, której akurat używa maszyna budująca. Jawne porównanie daje dwie rzeczy. Dla
tekstu maszynowego porównanie porządkowe (ordinal) zwraca tę samą odpowiedź na każdej
platformie i~usuwa zależność od obu bibliotek. Dla tekstu ludzkiego oznacza miejsce wywołania
jako lingwistyczne, więc można je wylistować i~sprawdzić jako miejsce, w~którym przejście
NLS → ICU może zmienić wynik. CA1310 znajduje wywołania, których domyślne zachowanie zależy
od kultury; comparerów, do których wracają API sortujące, nie zgłasza żaden analyzer, więc
wymagają wyszukania.

\paragraph{Jak spełnić.} Dla każdego miejsca wywołania rozstrzyga się, czy napisy są tekstem
maszynowym, czy ludzkim. Tekst maszynowy dostaje \repopath{Ordinal} albo
\repopath{OrdinalIgnoreCase}; tekst ludzki -- \repopath{CurrentCulture}, co zachowuje
semantykę legacy i~czyni ją jawną. Porównania porządkowe i~lingwistyczne zgadzają się dla
większości identyfikatorów ASCII, ale nie dla każdego wejścia: znaki o~zerowej wadze
i~ignorowalne (\repopath{\textbackslash 0}, miękki dywiz) oraz reguły wielkości liter się
różnią -- pod kulturą \repopath{tr-TR} „FILE” i~„file” nie są równe przy lingwistycznym
porównaniu bez rozróżniania wielkości liter, ale są równe przy \repopath{OrdinalIgnoreCase}.
Gdy wejście nie jest kontrolowane (kody wpisywane przez użytkowników, importowane dane) albo
serwer legacy działał z~taką kulturą, wybór zapisuje się jako flagę. Zapytań LINQ, które
Entity Framework tłumaczy na SQL, nie zmienia się: porównuje je collation bazy, a~nie .NET,
\repopath{StringComparer} nie da się przetłumaczyć, a~przejście NLS → ICU ich nie dotyczy.
\repopath{Dictionary<string, T>}, \repopath{HashSet<string>} i~\repopath{string.Equals(string)}
już porównują porządkowo -- zostają bez zmian. Przykład zgodny z~reguły:

\begin{lstlisting}
if (sku.StartsWith("GIFT-", StringComparison.Ordinal))          // a product code, not language
{
    ApplyGiftCardRules(order);
}

// Shown to people: still linguistic, now visibly so, and recorded as a behaviour flag
products.Sort((a, b) => string.Compare(a.Name, b.Name, StringComparison.CurrentCulture));
var byName = products.OrderBy(p => p.Name, StringComparer.CurrentCulture);
\end{lstlisting}

\paragraph{Zgodność i~flagi.} .NET SDK zgłasza \repopath{CA1310} z~wagą
\repopath{warning} (analyzer jest domyślnie wyłączony; włącza go wygenerowana konfiguracja).
Flagowane jest każde lingwistyczne sortowanie albo porównanie danych pokazywanych ludziom:
ICU może ułożyć lub dopasować je inaczej niż NLS, choć kod się nie zmienił. Flaga nazywa
miejsce wywołania i~dane, np. „product list sorted by name with CurrentCulture”.

\zrodla{\path{letsgolegacy.secondkey-standards/standards/SK-MIG-005-explicit-string-comparison.md},
\path{letsgolegacy.secondkey-standards/docs/adr/0001-standards-source-format.md}}

\subsubsection{SK-MIG-006 -- jawna kultura i~źródło kultury legacy}\label{sec:standardy-sk-mig-006}

Tytuł źródłowy: \emph{Name the culture for formatting, parsing and casing, and keep the
legacy culture source}. Waga \repopath{warning}, kategoria \repopath{globalization},
\repopath{appliesTo}: \repopath{"*.cs"}, analyzery \repopath{CA1304}, \repopath{CA1305},
\repopath{CA1311}.

\paragraph{Reguła.} Formatowanie i~parsowanie liczb, dat i~kwot oraz zamiana na wielkie
i~małe litery nazywają kulturę, której używają: \path{CultureInfo.InvariantCulture} dla
wartości przechowywanych, wymienianych lub porównywanych przez oprogramowanie -- zawartości
plików, kluczy cache, URL-i, tekstu w~bazie, ręcznie pisanego JSON -- oraz kulturę żądania
dla wartości pokazywanych ludziom. Każde żądanie działa pod kulturą, pod którą działało
w~systemie legacy, wziętą z~tego samego źródła.

\paragraph{Dlaczego.} W~migracji zmieniają się naraz dwie rzeczy. Po pierwsze dane kultury:
ICU i~NLS różnią się w~części z~nich -- dla kultury określonej tylko językiem, np.
\repopath{de}, \texttt{string.Format("\{0:C\}", 100)} daje „100,00~€” pod NLS
i~„100,00~¤” pod ICU, bo ICU traktuje walutę jako cechę kraju lub regionu, a~nie języka. Po
drugie, i~częściej, zmienia się \emph{źródło} kultury. ASP.NET Framework stosował
\texttt{<globalization culture="…" uiCulture="…" />} z~\repopath{web.config} do każdego
żądania. ASP.NET Core nie czyta \repopath{web.config}: działa z~kulturą procesu -- z~ustawień
regionalnych hosta na Windows, ze zmiennej środowiskowej albo bez niczego w~kontenerze --
chyba że skonfigurowano lokalizację żądań. Cena sformatowana kulturą ambientową może wrócić
z~innym separatorem dziesiętnym, a~data sparsowana tą kulturą może się nie sparsować albo
zamienić dzień z~miesiącem. Nazwanie kultury w~każdym miejscu wywołania czyni intencję
widoczną; lokalizacja żądań odtwarza resztę.

\paragraph{Jak spełnić.} Najpierw ustala się źródło kultury legacy: element
\repopath{<globalization>} w~\repopath{web.config} i~jego transformacjach,
\repopath{culture="auto"} (kultura z~nagłówka \repopath{Accept-Language} przeglądarki)
oraz ustawienia regionalne hosta IIS. Zapisuje się je na początku rejestru flag, bo od niego
zależy każda decyzja w~miejscach wywołań. Źródło odtwarza się przez
\repopath{UseRequestLocalization}: stała kultura staje się domyślną i~jedyną obsługiwaną;
\repopath{culture="auto"} -- dostawcą \repopath{Accept-Language} z~kulturami, które strona
faktycznie obsługiwała. W~każdym miejscu wywołania wybiera się kulturę, która odtwarza wynik
systemu legacy pod jego kulturą; tekst maszynowy zwykle przechodzi na
\repopath{InvariantCulture}, a~jeśli daje to inny wynik pod kulturą legacy (host
\repopath{tr-TR} zamieniający na wielką literę \repopath{i}), zapisuje się flagę. Parsowanie
wejścia użytkownika zachowuje kulturę żądania, którą miało; parsowanie tekstu
przechowywanego lub wymienianego przechodzi na \repopath{InvariantCulture} tylko wtedy, gdy
ten tekst był zapisywany niezmiennie. Przykłady z~reguły dla miejsc wywołań -- niezgodne:

\begin{lstlisting}[escapechar=!]
var cacheKey = "price:" + product.Id + ":" + price.ToString();  // separator depends on the host
var shipDate = DateTime.Parse(form["shipDate"]);                  // day/month order depends on the host
var code = couponCode.ToUpper();                                  // "file" becomes "F!\.{I}!LE" under tr-TR
\end{lstlisting}

i~zgodne:

\begin{lstlisting}
var cacheKey = string.Create(CultureInfo.InvariantCulture, $"price:{product.Id}:{price}");
var shipDate = DateTime.Parse(form["shipDate"], CultureInfo.CurrentCulture); // the request's culture, as before
var code = couponCode.ToUpperInvariant();
\end{lstlisting}

Odtworzenie źródła kultury z~\repopath{web.config}:

\begin{lstlisting}
// web.config had <globalization culture="pl-PL" uiCulture="pl-PL" />; reproduce it
var polish = new CultureInfo("pl-PL");
app.UseRequestLocalization(new RequestLocalizationOptions
{
    DefaultRequestCulture = new RequestCulture(polish),
    SupportedCultures = [polish],
    SupportedUICultures = [polish],
});
\end{lstlisting}

\paragraph{Zgodność i~flagi.} .NET SDK zgłasza \repopath{CA1304}, \repopath{CA1305}
i~\repopath{CA1311} z~wagą \repopath{warning}. Flagowane są: każda wartość pokazywana ludziom,
której format zależy od danych kultury (symbole walut, separatory, wzorce dat), bo dane ICU
mogą się różnić od danych NLS dla tej samej nazwy kultury, oraz każde miejsce, w~którym
jawna kultura daje inny wynik niż ambientowa kultura legacy.

\zrodla{\path{letsgolegacy.secondkey-standards/standards/SK-MIG-006-explicit-culture.md}}

\subsubsection{SK-MIG-007 -- tryb globalizacji bez zmian}\label{sec:standardy-sk-mig-007}

Tytuł źródłowy: \emph{Do not switch the globalization mode; flag every result NLS to ICU can
change}. Waga \repopath{error}, kategoria \repopath{globalization}, \repopath{appliesTo}:
\repopath{"*.csproj"}, \repopath{"*.props"}, \repopath{"*.json"}, \repopath{"Dockerfile"}.

\paragraph{Reguła.} Zmigrowany system działa z~ICU, domyślną biblioteką nowoczesnego .NET,
a~migracja nie zmienia tego za niczyimi plecami. Nie ustawia
\repopath{InvariantGlobalization}, \path{System.Globalization.UseNls} ani
\path{System.Globalization.AppLocalIcu} -- w~pliku projektu, w~\path{runtimeconfig.template.json}
ani w~zmiennej środowiskowej -- i~nie wybiera obrazu bazowego kontenera, który nie ma ICU.
Każde zachowanie, które może zmienić przejście z~NLS na ICU, jest zapisywane jako flaga
zachowania, a~nie łatane.

\paragraph{Dlaczego.} Biblioteka globalizacji jest jedyną różnicą migracyjną, której żaden
kompilator, analyzer ani test jednostkowy nie zgłasza niezawodnie: kod jest identyczny,
a~wyniki nie. Każdy przełącznik, który ją zmienia, jest uprawnionym narzędziem z~kosztem,
którego agent migracyjny nie jest w~stanie ocenić:
\begin{itemize}
  \item \path{InvariantGlobalization=true} usuwa dane kultury w~ogóle. Każda kultura
        zachowuje się jak kultura niezmienna, a~od .NET~6 utworzenie jakiejkolwiek innej
        kultury rzuca \repopath{CultureNotFoundException}, więc sortowanie, formatowanie
        i~wielkość liter zmieniają się naraz. Szablony projektów native AOT go włączają,
        a~niektóre obrazy kontenerów (np. Alpine) domyślnie działają w~trybie invariant.
  \item \repopath{UseNls=true} przywraca NLS, ale tylko na Windows: ten sam build zachowuje
        się inaczej na Linuksie, a~różnica pozostaje ukryta, dopóki platforma się nie zmieni.
  \item App-local ICU przypina jedną wersję ICU wszędzie. Spójnie, ale nadal jest to ICU,
        a~nie NLS.
\end{itemize}
Wybór jednego z~nich jest decyzją, która ma właściciela. Podjęty po cichu ukrywa dokładnie
ten rodzaj różnicy, który niezależna weryfikacja ma znaleźć.

\paragraph{Jak spełnić.} Jeśli test nie przechodzi z~powodu wyniku ICU, test zostaje taki,
jaki jest (SK-MIG-013), a~różnicę się zapisuje -- bez sięgania po przełącznik. Jeśli
platformą docelową jest kontener, sprawdza się, czy obraz bazowy ma ICU; jeśli nie ma, jest
to zapisywane jako decyzja do podjęcia. Jeśli człowiek zdecydował o~jednym z~przełączników,
wpisuje się go do projektu z~komentarzem wskazującym flagę, która zapisuje tę decyzję.
Zgodny plik projektu nie ma żadnego przełącznika (tylko \path{<TargetFramework>net10.0</TargetFramework>}),
a~sortowanie, które może się zmienić, trafia do rejestru, np.:

\begin{lstlisting}
| SK-MIG-007 | Catalog/ProductSorter.cs:41 | Products sorted by name with the current culture under NLS | Unchanged; ICU may order names with punctuation differently | Accept ICU order, or pin NLS on Windows |
\end{lstlisting}

\paragraph{Zgodność i~flagi.} Reguła nie ma diagnostyki -- kontrolą jest instrukcja agenta,
a~weryfikacją replay. Procedura P6 (sekcja~\ref{sec:standardy-modernize}) sprawdza ją
poleceniem \repopath{git grep} po nazwach przełączników. Flagowane są: każde lingwistyczne
sortowanie, wyszukiwanie lub porównanie tekstu ludzkiego (SK-MIG-005), każdy format zależny
od danych kultury pokazywany ludziom (SK-MIG-006) oraz nieprzechodzący test, którego
oczekiwana wartość zależy od zachowania NLS.

\zrodla{\path{letsgolegacy.secondkey-standards/standards/SK-MIG-007-keep-globalization-mode.md},
\path{letsgolegacy.secondkey-standards/docs/USING-WITH-MODERNIZE-DOTNET.md}}

\subsubsection{SK-MIG-008 -- wstrzykiwanie przez konstruktor}\label{sec:standardy-sk-mig-008}

Tytuł źródłowy: \emph{Resolve dependencies by constructor injection, not static singletons or
service locators}. Waga \repopath{warning}, kategoria \repopath{dependency-injection},
\repopath{appliesTo}: \repopath{"*.cs"}, zasada P10.

{\raggedright
\paragraph{Reguła.} Usługi są rejestrowane w~\path{Microsoft.Extensions.DependencyInjection}
(\repopath{IServiceCollection}) i~otrzymywane przez konstruktory. Migracja nie dodaje
statycznego singletona (statycznej właściwości \repopath{Instance}, statycznego pola
trzymającego usługę) ani wywołania service locatora
(\path{DependencyResolver.Current.GetService}, \path{IServiceProvider.GetService}
w~kodzie biznesowym, statycznego \repopath{Resolve} kontenera). Każda rejestracja zachowuje
czas życia, który miał obiekt legacy.\par}

\paragraph{Dlaczego.} Statyczny singleton albo service locator ukrywa zależności klasy
i~ustala je na cały czas życia procesu: czytelnik ich nie widzi, a~test nie może ich
podmienić. ASP.NET Core jest zbudowany na wbudowanym kontenerze. Każda możliwość opisana
przez konstytucję jest metodą rozszerzającą \repopath{IServiceCollection}, rozszerzenie to
interfejs i~jedna rejestracja (P10), a~konstytucja wyklucza własny kontener DI. Migracja
i~tak przepisuje okablowanie aplikacji, a~zapisanie go jako rejestracji kosztuje niewiele
więcej niż przenoszenie wywołań lokatora jedno po drugim. Niemechaniczną częścią jest czas
życia: obiekt per żądanie zarejestrowany jako singleton przenosi stan jednego klienta do
żądania innego, a~cache na cały proces zarejestrowany jako scoped po cichu traci zawartość
po każdym żądaniu. Czasy życia są zachowaniem.

\paragraph{Jak spełnić.} Każdą rejestrację legacy mapuje się na ten sam czas życia. Dla
Autofac: \repopath{SingleInstance} → \repopath{AddSingleton};
\repopath{InstancePerLifetimeScope} i~\repopath{InstancePerRequest} w~aplikacji webowej →
\repopath{AddScoped}; \repopath{InstancePerDependency} → \repopath{AddTransient}.
Wstrzykiwanie przez właściwości, dekoratory i~moduły nie mają odpowiednika jeden do jednego
we wbudowanym kontenerze; rejestracje z~kluczem mają (keyed services, .NET~8 i~nowsze).
Pozostawienie kontenera legacy na czas migracji przez jego integrację z~ASP.NET Core jest
decyzją do zapisania, a~nie wariantem domyślnym; nowe rejestracje i~tak przechodzą przez
\repopath{IServiceCollection}. Klasa statyczna bez stanu (czysta funkcja) może pozostać
statyczna. Tam, gdzie kod nie może przyjmować parametrów konstruktora (atrybut, statyczna
metoda rozszerzająca), przekazuje się mu to, czego potrzebuje; jeśli nie da się tego zrobić
bez zmiany zachowania -- flaga.

Przykład z~reguły: niezgodny kontroler pobiera \repopath{IOrderService} przez
\path{DependencyResolver.Current.GetService} i~cache cen przez statyczne
\repopath{PriceCache.Instance}; zgodny zachowuje czasy życia legacy -- cache cen był
procesowy, usługa zamówień była per żądanie:

\begin{lstlisting}
// The legacy lifetimes, kept: the price cache was process-wide, the order service per request
builder.Services.AddSingleton<IPriceCache, PriceCache>();
builder.Services.AddScoped<IOrderService, OrderService>();

public class CheckoutController(IOrderService orders, IPriceCache prices) : Controller
{
    public IActionResult Confirm()
    {
        // ...
    }
}
\end{lstlisting}

\paragraph{Zgodność i~flagi.} Reguła nie ma diagnostyki; jest tylko instrukcją dla agenta.
Flagowane są: czas życia, którego wbudowany kontener nie odtworzy, albo defekt legacy, np.
obiekt per żądanie przechwycony przez singleton (zachowanie się zachowuje i~zapisuje), oraz
zastąpienie kontenera legacy, jeśli migracja to robi.

\zrodla{\path{letsgolegacy.secondkey-standards/standards/SK-MIG-008-constructor-injection.md}}

\subsubsection{SK-MIG-009 -- klienci HTTP przez IHttpClientFactory}\label{sec:standardy-sk-mig-009}

Tytuł źródłowy: \emph{Create outbound HTTP clients through IHttpClientFactory, never per
call}. Waga \repopath{warning}, kategoria \repopath{outbound-http}, \repopath{appliesTo}:
\repopath{"*.cs"}, zasada P2.

\paragraph{Reguła.} Wychodzące wywołania HTTP używają klientów tworzonych przez
\repopath{IHttpClientFactory} -- typowanych, rejestrowanych przez
\path{AddHttpClient<TClient>()}, albo nazwanych -- skonfigurowanych raz, przy rejestracji,
adresem bazowym, nagłówkami i~limitem czasu, których używał klient legacy. Żaden kod nie
tworzy \repopath{HttpClient}, \repopath{WebClient} ani \repopath{HttpWebRequest} na każde
wywołanie. Migracja zachowuje zachowanie wychodzące: nie dodaje ponowień, circuit breakera
ani zmienionego limitu czasu.

\paragraph{Dlaczego.} \repopath{HttpClient} posiada pulę połączeń. Tworzenie go i~zwalnianie
przy każdym wywołaniu wyrzuca pulę za każdym razem, a~pod obciążeniem zamknięte połączenia
pozostają w~stanie \repopath{TIME\_WAIT}, dopóki hostowi nie zabraknie portów -- awaria
widoczna jako sporadyczne błędy gniazd daleko od kodu, który ją spowodował.
\repopath{IHttpClientFactory} pooluje i~wymienia handlery (co przy okazji uwzględnia zmiany
DNS) i~zamienia konfigurację każdego klienta w~rejestrację, którą czytelnik może znaleźć.
\repopath{WebClient} i~\repopath{HttpWebRequest} są w~nowoczesnym .NET przestarzałe
(\repopath{SYSLIB0014}). Pułapką po drugiej stronie są ponowienia. Stanem docelowym
konstytucji jest standardowy handler odporności na każdym kliencie (P2), ale ponowienie
wysyła żądanie jeszcze raz: \repopath{POST} do dostawcy płatności albo wysyłki, który
przekroczył limit czasu, może zostać wykonany dwukrotnie. To zmienia kontrakt wychodzący,
który według reguły sprawdza Second Key (w~rdzeniu \texttt{sk} nagrywa tylko ruch HTTP
wchodzący, więc wywołań wychodzących jeszcze nie sprawdza), więc należy do osobnej zmiany
z~własnym przeglądem, po migracji.
Ustawienia \repopath{ServicePointManager} (\repopath{SecurityProtocol},
\repopath{DefaultConnectionLimit}, \path{ServerCertificateValidationCallback}) konfigurowały
stos sieciowy .NET Framework; w~nowoczesnym .NET nie dotyczą \repopath{HttpClient}, więc
każde z~nich albo przechodzi jawnie do handlera klienta, albo przestaje działać.

\paragraph{Jak spełnić.} Każdy system zewnętrzny dostaje jednego klienta typowanego,
z~przeniesionym adresem bazowym, domyślnymi nagłówkami, obsługą poświadczeń i~limitem czasu
(\repopath{HttpWebRequest.Timeout}, podklasy \repopath{WebClient}, które go ustawiały). Każde
ustawienie \repopath{ServicePointManager}, na którym polegał kod legacy, przechodzi do
handlera (\path{ConfigurePrimaryHttpMessageHandler}) albo jest zapisane jako już
nieobowiązujące. Ciała żądań i~odpowiedzi pozostają bajtowo zgodne: ustawienia serializera,
kodowanie wartości w~query i~wielkość liter nagłówków takie, jakie produkował klient legacy.
Liczba wywołań zostaje ta sama: żadnej polityki ponowień, hedgingu ani circuit breakera
w~zmianie migracyjnej. W~przykładzie niezgodnym
\path{AddHttpClient<ShippingLabelClient>().AddStandardResilienceHandler()} sprawia, że zakup
etykiety, który przekroczył limit czasu, może zostać wykonany dwa razy.

Przykład zgodny z~reguły -- klient typowany, skonfigurowany raz przy rejestracji, z~limitem
czasu przeniesionym z~klienta legacy:

\begin{lstlisting}
builder.Services.AddHttpClient<RatesClient>(client =>
{
    client.BaseAddress = new Uri("https://rates.example.com/");
    client.Timeout = TimeSpan.FromSeconds(30); // the legacy client's timeout, carried over
});

public class RatesClient(HttpClient client)
{
    public Task<Rate[]?> GetRatesAsync(string zip) =>
        client.GetFromJsonAsync<Rate[]>($"v1/rates/{zip}");
}
\end{lstlisting}

Pierwszy przykład niezgodny tworzy klienta przy każdym wywołaniu
(\texttt{using var client = new HttpClient \{ BaseAddress = ... \};}), czyli nową pulę
połączeń na każde wywołanie.

\paragraph{Zgodność i~flagi.} Reguła nie ma diagnostyki; jest tylko instrukcją dla agenta,
a~zmiany wywołań wychodzących nie wykrywa w~tej wersji nic: \texttt{sk} nagrywa tylko ruch HTTP
wchodzący, a~nagrywanie wywołań wychodzących należy do fazy~02. Flagowane są: callback walidacji
certyfikatu albo ustawienie TLS, które przestało działać; limit czasu, przekierowanie albo
dekompresja, których wartości domyślne różnią się między klientem legacy a~\repopath{HttpClient},
gdy migracja nie może ich odtworzyć dokładnie; dodanie odporności -- zapisane jako zmiana
uzupełniająca, o~którą prosi konstytucja.

\zrodla{\path{letsgolegacy.secondkey-standards/standards/SK-MIG-009-httpclient-factory.md},
\path{letsgolegacy.secondkey/schemas/skcap.schema.json},
\path{letsgolegacy.secondkey/docs/formats/contract.md}}

\subsubsection{SK-MIG-010 -- zachowanie EF6 przy dostępie do danych}\label{sec:standardy-sk-mig-010}

Tytuł źródłowy: \emph{Preserve query, loading and validation behaviour when data access
leaves EF6}. Waga \repopath{error}, kategoria \repopath{data-access}, \repopath{appliesTo}:
\repopath{"*.cs"}, zasada P4.

\paragraph{Reguła.} Przeniesienie platformy zachowuje dostęp do danych dokładnie takim, jaki
był. EF6 (6.4 i~nowsze) działa na nowoczesnym .NET, więc domyślnie najpierw przenosi się
system na .NET~10 z~EF6, a~port do EF Core jest osobnym krokiem. Gdy migracja jednak portuje
do EF Core, odtwarza każde zachowanie EF6, na którym polegał kod: dane powiązane czytane przez
lazy loading, nazwy tabel i~typy kolumn istniejącego schematu, walidację wykonywaną przez EF6
przy zapisie oraz stany encji nadawane przez \repopath{Attach} i~\repopath{Update} w~grafie
obiektów.

\paragraph{Dlaczego.} Istotne różnice EF6 → EF Core to te, które się kompilują:
\begin{itemize}
  \item EF6 domyślnie ładuje leniwie wirtualne właściwości nawigacyjne. EF Core nie, chyba że
        dodano i~włączono pakiet proxy, więc nawigacja, na której polegał kod, jest pusta albo
        \repopath{null} -- zamówienie bez pozycji, wyrenderowane bez błędu.
  \item EF6 wyprowadza nazwy tabel przez pluralizację nazwy klasy; EF Core używa nazwy
        właściwości \repopath{DbSet} albo nazwy klasy.
  \item EF6 waliduje encje względem adnotacji danych przy \repopath{SaveChanges}. EF Core nie
        waliduje, więc dane odrzucane przez EF6 docierają teraz do bazy.
  \item EF Core decyduje, czy encja znaleziona w~grafie jest dodana, czy niezmieniona, na
        podstawie tego, czy ma ustawiony klucz generowany przez bazę -- inaczej niż EF6.
  \item EF Core mapuje domyślnie \repopath{DateTime} na \repopath{datetime2}, a~schemat
        utworzony przez EF6 używa \repopath{datetime}, z~inną precyzją.
\end{itemize}
Wytyczne Microsoft dotyczące portowania wymieniają te różnice jako zmiany zachowania, które
nie pojawiają się jako błędy kompilacji, i~zalecają przejście na EF6 w~nowoczesnym .NET przed
portem do EF Core.

\paragraph{Jak spełnić.} Dla przeniesienia platformy preferowany jest EF6 6.4 lub nowszy na
\repopath{net10.0}, a~port do EF Core -- w~późniejszej zmianie. Przy portowaniu
inwentaryzuje się każdą właściwość nawigacyjną czytaną bez \repopath{Include} i~ładuje ją
jawnie albo włącza proxy lazy loading (\repopath{UseLazyLoadingProxies}) jako zapisaną
decyzję. Każdą encję mapuje się jawnie na istniejącą tabelę i~typy kolumn; pierwsza migracja
EF Core opisuje istniejący schemat, a~nie go zmienia (SK-ARCH-001). Tam, gdzie kod legacy
polegał na tym, że EF6 odrzuci niepoprawną encję przy zapisie, walidacja odbywa się przed
\repopath{SaveChanges}, żeby to samo wejście nadal było odrzucane. EF Core rzuca wyjątek
w~czasie działania dla zapytania, którego nie potrafi przetłumaczyć; takie zapytania
przepisuje się tak, by zwracały te same wiersze, i~porównuje wygenerowany SQL tam, gdzie
wyniki mogłyby się różnić. Przykłady zgodne z~reguły:

\begin{lstlisting}
// Load exactly what the legacy code read through lazy loading
var order = await db.Orders
    .Include(o => o.Lines)
    .SingleAsync(o => o.Id == id);
var total = order.Lines.Sum(l => l.Quantity * l.UnitPrice);
\end{lstlisting}

\begin{lstlisting}
// Map to the existing schema instead of letting conventions choose
modelBuilder.Entity<Order>(order =>
{
    order.ToTable("Orders");
    order.Property(o => o.CreatedOnUtc).HasColumnType("datetime");
});
\end{lstlisting}

\paragraph{Zgodność i~flagi.} Reguła nie ma diagnostyki; kontrolą jest instrukcja agenta,
a~weryfikacją replay. Flagowane są: walidacja wykonywana przez EF6, której zmigrowany kod nie
odtwarza dokładnie; zapytanie bez \repopath{OrderBy}, którego wiersze wracają pod nowym SQL
w~innej kolejności (kolejność nigdy nie była gwarantowana, więc zapisuje się różnicę zamiast
dodawać sortowanie, którego kod legacy nie miał); każda zmiana typu kolumny, nazwy tabeli albo
klucza, które baza już ma.

\zrodla{\path{letsgolegacy.secondkey-standards/standards/SK-MIG-010-preserve-ef-behaviour.md}}

\subsubsection{SK-MIG-011 -- zachowaj zachowanie, flaguj zamiast naprawiać}\label{sec:standardy-sk-mig-011}

Tytuł źródłowy: \emph{Preserve behaviour; flag a suspected defect instead of fixing it}. Waga
\repopath{error}, kategoria \repopath{behaviour}, \repopath{appliesTo}: \repopath{"*"}
(wszystkie pliki), zasada P14.

\paragraph{Reguła.} Migracja zmienia platformę, a~nie biznes. Każde obserwowalne zachowanie
systemu legacy -- wyniki, zaokrąglenia, kolejność, wyniki walidacji, odpowiedzi błędów oraz
skutki uboczne, takie jak e-maile, rekordy audytu i~wywołania wychodzące -- zostaje
zachowane, także zachowanie, które wygląda na błędne. Gdy migracja znajduje podejrzewany
defekt, martwy kod albo niespójność, zostawia logikę taką, jaka była, oznacza miejsce
komentarzem \repopath{SECONDKEY-FLAG} i~zapisuje je w~rejestrze flag zachowania. Decyzja
o~naprawie należy do człowieka, w~osobnej zmianie.

\paragraph{Dlaczego.} Migracja ma sens tylko wtedy, gdy nowy system robi to, co stary,
a~jedynym sposobem, żeby to pokazać, jest porównanie obu. \finding{„Naprawa” wykonana
w~trakcie migracji jest w~tym porównaniu nieodróżnialna od regresji: obie są różnicą, o~którą
nikt nie prosił.} Taka naprawa łączy też dwa przeglądy -- „czy przeniesienie platformy się
udało?” i~„czy ta zmiana biznesowa jest słuszna?” -- w~jeden diff, który na żadne z~tych
pytań nie odpowiada dobrze. Zapisanie znaleziska nie traci niczego, co agent zauważył,
a~poprawkę można przejrzeć osobno, gdy zmigrowany system jest już znany jako równoważny.
Zapisanie decyzji w~repozytorium, z~powodem, jest regułą konstytucji dla decyzji (P14).
W~przykładzie niezgodnym kod legacy stosuje rabaty procentowe w~kolejności listy (więc dwa
rabaty się składają), a~migracja „naprawia” to do zastosowania tylko największego -- cicha
zmiana każdej ceny.

\paragraph{Jak spełnić.} Sekcja \emph{Migration} tej reguły definiuje protokół flag, do
którego odwołuje się sekcja \emph{Flag instead of fixing} każdej innej reguły; opisuje go
sekcja~\ref{sec:standardy-flagi}. Jako zachowanie do zachowania liczy się:
\begin{itemize}
  \item Dodanie sprawdzenia \repopath{null}, wartości domyślnej albo bloku
        \repopath{try}/\repopath{catch}, których kod legacy nie miał, zmienia zachowanie:
        \repopath{NullReferenceException}, który dawał HTTP 500, nadal go daje.
  \item Zaokrąglanie zostaje takie, jak było. \repopath{Math.Round} domyślnie używa
        zaokrąglania bankierskiego (\repopath{MidpointRounding.ToEven}) na obu platformach;
        trybu się nie zmienia.
  \item Skutki uboczne zachowują wyzwalacze i~kolejność: ten sam e-mail dla tego samego
        zdarzenia, ten sam rekord audytu, to samo wywołanie wychodzące.
  \item Ostrzeżenia kompilatora o~nieosiągalnym kodzie albo nieużytych wartościach są
        znaleziskami do zapisania, a~nie kodem do usunięcia, jeśli usunięcie zmieniłoby to,
        co się wykonuje.
\end{itemize}

\paragraph{Zgodność i~flagi.} Reguła nie ma diagnostyki; kontrolą jest instrukcja agenta,
a~weryfikacją -- replay i~porównanie Second Key. Flagowane są: podejrzewany defekt w~logice
biznesowej, nawet oczywisty; słabość bezpieczeństwa w~kodzie legacy -- flagowana na początku
rejestru, żeby człowiek zobaczył ją pierwszą, a~naprawiana w~migracji tylko wtedy, gdy sama
migracja inaczej uczyniłaby ją możliwą do wykorzystania (co flaga mówi wprost); martwy kod,
zduplikowane reguły, które się ze sobą nie zgadzają, i~walidacja niespójna między dwiema
ścieżkami.

\zrodla{\path{letsgolegacy.secondkey-standards/standards/SK-MIG-011-preserve-behaviour.md}}

\subsubsection{SK-MIG-012 -- kontrakt HTTP bez zmian}\label{sec:standardy-sk-mig-012}

Tytuł źródłowy: \emph{Keep the HTTP contract --- routes, status codes, headers, cookies and
JSON shape}. Waga \repopath{error}, kategoria \repopath{behaviour}, \repopath{appliesTo}:
\repopath{"*.cs"}, \repopath{"*.cshtml"}.

\paragraph{Reguła.} Na każde żądanie, na które odpowiadał system legacy, po migracji pada ta
sama odpowiedź: ten sam URL i~ta sama metoda HTTP trafiają do tej samej akcji, a~odpowiedź
ma ten sam kod statusu, cel przekierowania, content type, nagłówki, nazwy i~atrybuty
ciasteczek oraz -- dla JSON -- te same nazwy właściwości i~wielkość liter, formaty dat
i~liczb oraz traktowanie wartości \repopath{null} i~pól.

\paragraph{Dlaczego.} Przeglądarki, zakładki, wyszukiwarki, callbacki płatności i~inne
systemy zależą od formatu na łączu, a~nie od kodu. ASP.NET Core zmienia naraz kilka ustawień
domyślnych:
\begin{itemize}
  \item JSON jest domyślnie w~camelCase, podczas gdy \repopath{JsonResult} ASP.NET MVC
        i~Web API zapisywały nazwy właściwości tak, jak je zadeklarowano.
  \item \repopath{System.Text.Json} ignoruje publiczne pola, które \repopath{Newtonsoft.Json}
        i~serializer MVC zapisywały, a~\repopath{Controller.Json} w~MVC~5 zapisywał daty jako
        \repopath{\textbackslash/Date(…)\textbackslash/} zamiast ISO~8601.
  \item Nazwy ciasteczek się różnią: forms authentication używało \repopath{.ASPXAUTH},
        a~stan sesji \repopath{ASP.NET\_SessionId}; domyślne nazwy ASP.NET Core to
        \repopath{.AspNetCore.Cookies} i~\repopath{.AspNetCore.Session}.
\end{itemize}
Każda z~tych zmian jest niewidoczna w~przeglądzie kodu i~oczywista dla klienta HTTP. Second
Key odtwarza nagrany ruch na obu systemach i~porównuje odpowiedzi pole po polu, więc każda
z~nich jest też zgłaszaną różnicą.

\paragraph{Jak spełnić.} Trasy przenosi się z~tymi samymi szablonami, wartościami domyślnymi
i~ograniczeniami, w~tej samej kolejności; trasy konwencyjne pozostają konwencyjne. Dla
każdego endpointu JSON porównuje się odpowiedź legacy ze zmigrowaną -- nazwy, wielkość liter,
daty, liczby, wartości \repopath{null}, publiczne pola, reprezentację enumów -- i~konfiguruje
serializer albo dodaje konwerter, aż się zgadzają. Kody statusu i~przekierowania zostają,
także w~odpowiedziach błędów: nieobsłużony wyjątek, który dawał 500 z~własną stroną błędu,
nadal ją daje. Nazwy nagłówków i~ciasteczek czytane przez klientów zostają; tam, gdzie nazwa
musi się zmienić (ciasteczek uwierzytelniania nie da się współdzielić między oboma
formatami), zapisuje się flagę. Przykłady zgodne z~reguły:

\begin{lstlisting}
builder.Services.AddControllersWithViews()
    .AddJsonOptions(options =>
    {
        // Keep property names as declared, as MVC 5 and Web API did
        options.JsonSerializerOptions.PropertyNamingPolicy = null;
    });
\end{lstlisting}

\begin{lstlisting}
// Web API endpoints built on Newtonsoft.Json can keep it, with its settings, during the migration
builder.Services.AddControllers().AddNewtonsoftJson();
\end{lstlisting}

\paragraph{Zgodność i~flagi.} Reguła nie ma diagnostyki; różnice wykrywa replay
i~porównanie pole po polu. Flagowane są: różnica ciasteczka, nagłówka albo odpowiedzi,
której migracja nie może uniknąć, np. ciasteczka uwierzytelniania, przez które przy
przełączeniu każdy użytkownik zostaje wylogowany, oraz odpowiedź legacy, która wygląda na
błędną (błąd zwracany ze statusem 200, literówka w~nazwie właściwości) -- zachowuje się ją
i~zapisuje (SK-MIG-011).

\zrodla{\path{letsgolegacy.secondkey-standards/standards/SK-MIG-012-keep-http-contract.md}}

\subsubsection{SK-MIG-013 -- testy są dowodem}\label{sec:standardy-sk-mig-013}

Tytuł źródłowy: \emph{Keep existing tests and their expectations; a failing test is a
finding}. Waga \repopath{error}, kategoria \repopath{behaviour}, \repopath{appliesTo}:
\repopath{"*.cs"}, zasada P13.

\paragraph{Reguła.} Istniejące testy migrują razem z~kodem i~zachowują asercje. Test może
się zmienić tylko po to, by nadążyć za przeniesieniem API -- przestrzeń nazw, host testowy,
sygnatura \repopath{async} -- nigdy w~oczekiwanych wartościach, tolerancjach ani wejściach,
które sprawdza. Gdy test nie przechodzi po migracji, zmienia się zmigrowany kod, aż test
przejdzie w~obecnej postaci. Jeśli różnicy nie da się usunąć, test zostaje taki, jaki jest,
a~porażka jest zapisana jako flaga zachowania; testu nigdy się nie usuwa, nie pomija ani nie
przepisuje tak, by przeszedł. Testy charakteryzujące, które przypinają zachowanie legacy,
pisze się przeciw kodowi legacy, przed przeniesieniem.

\paragraph{Dlaczego.} Testy są dowodem migracji na to, że zachowanie przetrwało.
Oczekiwanie poprawione tak, by pasowało do nowego wyniku, zamienia ten dowód w~jego
przeciwieństwo: jedyny sygnał, który złapał różnicę, zostaje uciszony. Konstytucja wymaga
testów charakteryzujących przed przeniesieniem, a~nie po nim, z~tego samego powodu (P13):
test napisany po migracji opisuje nowe zachowanie, jakiekolwiek się ono okazało. Przykłady
niezgodne z~reguły: oczekiwana wartość VAT zmieniona z~\repopath{24.39m} na
\repopath{24.40m}, „żeby build był zielony”, oraz test wyłączony atrybutem
\texttt{[Fact(Skip = ...)]} z~powodem „Fails after the upgrade”.

\paragraph{Jak spełnić.} Projekty testowe przenosi się razem z~kodem: ten sam framework
testowy albo jego bieżąca wersja, te same asercje. Nieprzechodzący test najpierw oznacza, że
zmigrowany kod jest błędny: trzeba znaleźć zachowanie, które się zmieniło, i~je przywrócić
(w~przykładzie zgodnym test sumy zamówienia przestał przechodzić, bo EF Core nie ładuje
leniwie pozycji; zapytanie dołącza pozycje zgodnie z~SK-MIG-010, a~asercja jest nietknięta).
Tylko wtedy, gdy różnica jest nieodłączną cechą platformy -- kolejność sortowania ICU,
usunięte API bez odpowiednika -- test zostaje nieprzechodzący, dodaje się flagę i~pisze o~tym
w~podsumowaniu. Nowe testy pisane w~trakcie migracji przypinają zachowanie legacy,
zweryfikowane na systemie legacy; nie asertują tego, co akurat robi system zmigrowany.

\paragraph{Zgodność i~flagi.} Reguła nie ma diagnostyki; kontrolą są instrukcja agenta
i~uruchomienie istniejących testów po każdym zadaniu. Flagowane są: każdy test, który na
koniec migracji nadal nie przechodzi, z~przyczyną, oraz test, który już przed migracją nie
przechodził albo był pomijany -- zostaje taki, jaki był, i~jest zapisany.

\zrodla{\path{letsgolegacy.secondkey-standards/standards/SK-MIG-013-tests-are-evidence.md}}

\subsubsection{Źródła pierwotne cytowane przez reguły}\label{sec:standardy-referencje}

Każda reguła poza SK-MIG-011 i~SK-MIG-013 kończy się sekcją \emph{References}, która podaje
źródła pierwotne jej twierdzeń technicznych; generator przenosi ją bez zmian do
\path{references/<id>.md} skilla.

{\small
\begin{xltabular}{\linewidth}{@{}>{\raggedright\arraybackslash}p{2.3cm} >{\raggedright\arraybackslash}X@{}}
\toprule
Reguła & Źródła (tytuły jak w~regule) \\
\midrule
\endfirsthead
\toprule
Reguła & Źródła (tytuły jak w~regule) \\
\midrule
\endhead
SK-ARCH-001 & Konstytucja, P4 „Persistence is provider-portable and its schema is migrated, never ensured”; EF Core, „Porting from EF6 to EF Core” -- inicjalizacja bazy w~Code First \\
SK-ARCH-002 & Konstytucja, P4 -- „Migrations describe schema; reference data is seeded separately” \\
SK-ARCH-003 & Konstytucja, P5 „Configuration through the environment; secrets through the platform” \\
SK-ARCH-004 & Konstytucja, P9 „Program.cs is a manifest; wiring lives in extension methods” \\
SK-ARCH-005 & Konstytucja, P9 -- warstwy kontrolerów, orkiestratorów, usług domenowych i~repozytoriów \\
SK-ARCH-006 & Konstytucja, P10 „Extensibility through interface + registration, not inheritance” \\
SK-ARCH-007 & Konstytucja, P15 „Observability is a build-time decision, not an afterthought” i~P2a „every service calls AddServiceDefaults()” \\
SK-MIG-001 & ASP.NET Core, „Migrate from ASP.NET Framework to ASP.NET Core” -- sekcja \path{migration/fx-to-core}, w~tym adaptery System.Web i~obszar \repopath{HttpContext} (buforowanie, powiązanie z~wątkiem) \\
SK-MIG-002 & ASP.NET Core, „Access HttpContext in ASP.NET Core” -- bezpieczeństwo wątkowe i~praca w~tle \\
SK-MIG-003 & ASP.NET Core, „Kestrel options” -- \repopath{AllowSynchronousIO} domyślnie \repopath{false} \\
SK-MIG-004 & ASP.NET Core, „Migrate configuration to ASP.NET Core” (\path{migration/fx-to-core}); .NET, „Options pattern in .NET” \\
SK-MIG-005 & .NET, „Globalization and ICU” -- różnice NLS i~ICU; .NET, „Best practices for comparing strings in .NET”; reguła CA1310 „Specify StringComparison for correctness” \\
SK-MIG-006 & .NET, „Globalization APIs use ICU libraries on Windows” (zmiana łamiąca w~.NET~5) -- przykład z~walutą; reguły CA1304, CA1305 i~CA1311 \\
SK-MIG-007 & .NET, „Globalization and ICU” -- przełączniki NLS i~app-local ICU; .NET, „Runtime configuration options for globalization”; zmiana łamiąca .NET~6 „Culture creation and case mapping in globalization-invariant mode” \\
SK-MIG-008 & Konstytucja, P10 „Extensibility through interface + registration, not inheritance” i~§4 „Deliberate non-goals” (bez własnego kontenera DI) \\
SK-MIG-009 & .NET, „HttpClient guidelines for .NET” -- pule połączeń, DNS i~\repopath{IHttpClientFactory}; konstytucja, P2 -- handler odporności jako stan docelowy \\
SK-MIG-010 & EF Core, „Porting from EF6 to EF Core” -- „Behavior changes between EF6 and EF Core” i~zalecana kolejność kroków; EF Core, „Lazy loading of related data” \\
SK-MIG-012 & ASP.NET Core, „Format response data in ASP.NET Core Web API” -- camelCase domyślnie, \texttt{PropertyNamingPolicy = null} dla nazw zadeklarowanych; .NET, „Migrate from Newtonsoft.Json to System.Text.Json” -- pola, escapowanie i~wielkość liter \\
\bottomrule
\end{xltabular}
}

\zrodla{\path{letsgolegacy.secondkey-standards/standards/} (sekcje \emph{References} reguł),
\path{letsgolegacy.secondkey-standards/AGENTS.md}}

\subsection{Protokół flag zachowania}\label{sec:standardy-flagi}

Flaga zachowania jest sposobem, w~jaki agent zgłasza różnicę, której nie wolno mu rozstrzygnąć
samodzielnie. Protokół jest zdefiniowany w~regule SK-MIG-011, bo „flag instead of fixing” to
centralne zachowanie, o~które proszą standardy; jest więc wersjonowany razem z~regułami
i~dociera do agenta tym samym kanałem. Cztery kroki:

\begin{enumerate}
  \item \textbf{W~miejscu kodu.} Komentarz w~wierszu nad zachowanym kodem:
        \texttt{// SECONDKEY-FLAG <rule-id>: <...>}, gdzie po dwukropku jest jedno zdanie
        o~tym, co się różni albo wygląda źle. W~Razor:
        \repopath{@* SECONDKEY-FLAG … *@}; w~XML i~plikach konfiguracyjnych:
        \repopath{<!-{}- SECONDKEY-FLAG … -{}->}.
  \item \textbf{W~rejestrze.} Jeden wiersz w~\path{docs/migration/behaviour-flags.md}
        zmigrowanego repozytorium (plik tworzy się, jeśli nie istnieje), z~kolumnami:
        identyfikator reguły, \repopath{path:line}, zachowanie legacy, co zrobiła migracja,
        decyzja, którą musi podjąć człowiek.
  \item \textbf{Przy tłumieniu diagnostyki bramki}, bo zachowanie musi zostać -- tłumienie
        nazywa flagę (przykład poniżej).
  \item \textbf{W~podsumowaniu.} Raport końcowy migracji i~opis jej pull requestu kończą się
        sekcją „Behaviour flags”: liczba flag pogrupowana według reguł i~link do rejestru.
        Według skilla znaleziska bezpieczeństwa idą na początek.
\end{enumerate}

Formaty przepisane ze skilla i~z~reguły SK-MIG-011 (przykład w~miejscu kodu, nagłówek rejestru
z~przykładowym wierszem, ogólna postać tłumienia oraz przykład tłumienia z~reguły):

\begin{lstlisting}
// SECONDKEY-FLAG SK-MIG-011: percentage discounts compound in list order; looks unintended.
\end{lstlisting}

\begin{lstlisting}
# Behaviour flags

Legacy culture source: (the web.config globalization element, or the host's regional settings)

| Rule | Location | Legacy behaviour | What the migration did | Decision needed |
|---|---|---|---|---|
| SK-MIG-011 | Pricing/DiscountCalculator.cs:57 | Percentage discounts compound in list order | Preserved unchanged | Is compounding intended? |
\end{lstlisting}

\begin{lstlisting}[escapechar=!]
#pragma warning disable <DIAGNOSTIC> // SECONDKEY-FLAG <rule-id>: <why>
#pragma warning disable PORTCULLIS_MIG_SYNC_OVER_ASYNC // SECONDKEY-FLAG SK-MIG-003: !\ldots!
\end{lstlisting}

Protokół występuje w~dwóch miejscach: definiuje go reguła (reguły są źródłem), a~skill
powtarza format, bo agent potrzebuje go w~tej części skilla, którą ma zawsze załadowaną. Test
\repopath{SkillContentTests} sprawdza, że oba teksty zawierają ten sam znacznik
(\repopath{SECONDKEY-FLAG <rule-id>:}), tę samą ścieżkę rejestru i~ten sam nagłówek tabeli
rejestru (ADR 0004, decyzja 6). Ścieżkę rejestru generator wstawia do skilla ze stałej
\path{SkillEmitter.FlagRegister}.

\zrodla{\path{letsgolegacy.secondkey-standards/standards/SK-MIG-011-preserve-behaviour.md},
\path{letsgolegacy.secondkey-standards/standards/README.md},
\path{letsgolegacy.secondkey-standards/catalog/skill.template.md},
\path{letsgolegacy.secondkey-standards/docs/adr/0004-skill-for-github-copilot-upgrade.md},
\path{letsgolegacy.secondkey-standards/tests/SecondKey.Standards.Generator.Tests/Generation/SkillContentTests.cs},
\path{letsgolegacy.secondkey-standards/src/SecondKey.Standards.Generator/Generation/SkillEmitter.cs}}

\subsection{Generator secondkey-standards (ADR 0002)}\label{sec:standardy-generator}

Standardy potrzebują programu, który je waliduje (C4-a), kompiluje do skilla, konfiguracji
analyzerów i~pakietu NuGet (C4-b) oraz sprawdza wersję przypiętą przez konsumenta (C4-c).
Drugi generator w~estate (\path{architecture-standards/scripts/build-marketplace.mjs}) jest
skryptem Node; zasada produktu nr 3 Second Key brzmi natomiast: tylko technologie, które
odbiorca już utrzymuje -- .NET, SQL Server, Entra ID, GitHub albo Azure DevOps.

\subsubsection{Decyzje ADR 0002}

\begin{enumerate}
  \item \textbf{Aplikacja konsolowa .NET~10, pakowana jako dotnet tool}
        (\repopath{secondkey-standards}). Osoby, które ją uruchamiają, mają już .NET SDK;
        nic więcej nie jest potrzebne. Ta sama binarka działa w~CI tego repozytorium i~--
        jako drift check -- w~CI konsumenta.
  \item \textbf{Dwie zależności, obie zadeklarowane} (w~\path{Directory.Packages.props}
        i~w~tabeli zależności README). \repopath{YamlDotNet} czyta front matter: własny
        podzbiór YAML byłby mały, ale przypadki brzegowe YAML (cytowanie, aliasy, zduplikowane
        klucze) są dokładnie tym miejscem, w~którym ręcznie pisany parser po cichu źle
        odczytuje regułę; biblioteka to jedno assembly bez zależności, a~jej model węzłów daje
        numery wierszy potrzebne w~komunikatach. \repopath{System.CommandLine} (2.0, stabilna)
        daje pomoc, parsowanie i~błędy dla trzech poleceń bez ręcznej obsługi argumentów.
        Sekcje Markdown są parsowane ręcznie: gramatyka reguły to nagłówki i~bloki kodu,
        a~jedyną subtelność -- \repopath{\#\# heading} wewnątrz bloku kodu nie jest nagłówkiem
        -- obsługuje jawne śledzenie ogrodzeń. Biblioteka Markdown byłaby większą zależnością
        dającą mniej kontroli.
  \item \textbf{Walidator powstaje w~C4-a, razem ze standardami.} Kryterium akceptacji
        C4-a -- każda reguła ma identyfikator, wagę, uzasadnienie, przykład zgodny i~niezgodny
        -- egzekwuje CI w~tym samym pull requeście, który wprowadza reguły, a~nie przegląd aż
        do C4-b. C4-b dodaje do tego samego narzędzia emisję i~\repopath{-{}-check}.
  \item \textbf{Problemy są zbierane, a~nie rzucane.} Jedno uruchomienie zgłasza każdą
        zepsutą regułę z~plikiem i~wierszem; w~GitHub Actions każdy problem jest też adnotacją
        na diffie.
  \item \textbf{Kody wyjścia wspólne dla wszystkich poleceń}: 0 -- sukces, 1 -- sprawdzenie
        się wykonało i~wykazało problem, 2 -- polecenie nie mogło się wykonać. Pipeline
        odróżnia nieudane sprawdzenie od zepsutego wywołania.
  \item \textbf{Testy na xUnit v3 i~Microsoft.Testing.Platform}, włączonej przez
        \repopath{global.json}, czego xUnit v3 wymaga od \repopath{dotnet test} w~.NET~10 SDK.
  \item \textbf{Repozytorium stosuje do siebie własne reguły globalizacji.}
        \repopath{.editorconfig} podnosi CA1304, CA1305, CA1307, CA1309, CA1310 i~CA1311 do
        ostrzeżeń, a~\path{TreatWarningsAsErrors} robi z~nich błędy builda: obsługa
        napisów w~generatorze jest porządkowa i~jawna co do kultury. Ten sam plik podnosi
        CA1849 (blokujące wywołanie tam, gdzie istnieje wersja z~\repopath{await}) -- według
        komentarza sync-over-async nie ma analyzera w~SDK, a~CA1849 jest najbliższy.
\end{enumerate}

Konsekwencje: \texttt{dotnet run -{}-project src/SecondKey.Standards.Generator -{}- <polecenie>}
działa z~klonu, w~którym jest tylko SDK, a~\repopath{dotnet pack} tworzy pakiet narzędzia.
Każde polecenie przyjmuje strumienie, katalog roboczy i~środowisko jako parametry
(\repopath{CliContext}), więc testy uruchamiają prawdziwą powierzchnię poleceń w~procesie,
na jednorazowych repozytoriach.

\subsubsection{Projekt i~budowanie}

\begin{tabularx}{\linewidth}{@{}>{\raggedright\arraybackslash}p{4.6cm} >{\raggedright\arraybackslash}X@{}}
\toprule
Plik & Ustawienia \\
\midrule
\path{SecondKey.Standards.Generator.csproj} & \repopath{OutputType} \repopath{Exe}, \repopath{PackAsTool} \repopath{true}, \repopath{ToolCommandName} \repopath{secondkey-standards}, \repopath{PackageId} \path{SecondKey.Standards.Generator}, \repopath{Version} \repopath{0.1.0}; \repopath{InternalsVisibleTo} dla projektu testów \\
\path{Directory.Build.props} & \repopath{net10.0}, \repopath{LangVersion} \repopath{latest}, \repopath{Nullable} \repopath{enable}, \repopath{ImplicitUsings} \repopath{enable}, \repopath{TreatWarningsAsErrors} \repopath{true}, \repopath{AnalysisLevel} \repopath{latest}, \repopath{Deterministic} \repopath{true} \\
\path{global.json} & SDK \repopath{10.0.100}, \repopath{rollForward} \repopath{latestFeature}; \repopath{test.runner} \path{Microsoft.Testing.Platform} \\
\path{SecondKey.Standards.slnx} & Dwa projekty: generator (\path{src/}) i~testy (\path{tests/}) \\
\path{.gitattributes} & \repopath{eol=lf} dla wszystkich plików tekstowych: \repopath{-{}-check} porównuje bajty, więc checkout z~CRLF na Windows pokazałby każdy wygenerowany plik jako nieaktualny; \path{generated/**} oznaczone \repopath{linguist-generated} \\
\bottomrule
\end{tabularx}

Polecenia pracy w~repozytorium (README); wymagania: .NET~10 SDK oraz gitleaks albo Docker
dla skanu sekretów w~hooku pre-commit:

\begin{lstlisting}
./scripts/setup.sh                                                          # prerequisites + pre-commit hook
dotnet test                                                                 # the tests
dotnet run --project src/SecondKey.Standards.Generator -- validate          # check every rule
dotnet run --project src/SecondKey.Standards.Generator -- generate          # regenerate generated/
dotnet run --project src/SecondKey.Standards.Generator -- generate --check  # what CI runs
./scripts/verify-package.sh                                                 # pack the NuGet package and prove a build picks it up
./scripts/scan-secrets.sh                                                   # mirror the CI secret scan
\end{lstlisting}

Zmiana reguły to: edycja pliku, \repopath{generate}, commit obu. CI odrzuca pull request,
którego \path{generated/} nie odpowiada źródłom, i~wymienia każdy nieaktualny plik.

\subsubsection{Polecenia}

\begin{xltabular}{\linewidth}{@{}>{\raggedright\arraybackslash}p{3.3cm} >{\raggedright\arraybackslash}p{3.4cm} >{\raggedright\arraybackslash}X@{}}
\toprule
Polecenie & Opcje & Działanie i~wynik \\
\midrule
\endfirsthead
\toprule
Polecenie & Opcje & Działanie i~wynik \\
\midrule
\endhead
\repopath{validate} & \repopath{-{}-root} & Sprawdza każdą regułę w~\path{standards/}: front matter, słowniki, nazwy plików, uzasadnienie i~oba przykłady. Sukces: „\texttt{<n> rules are valid (<m> mapped to gate diagnostics).}”, kod 0. Problemy: każdy jako \texttt{ścieżka:wiersz: komunikat} na stderr i~podsumowanie „\texttt{<k> problem(s) in standards/.}”, kod 1 \\
\repopath{generate} & \repopath{-{}-root}, \repopath{-{}-check} & Kompiluje \path{standards/} do \path{generated/}: skill, \repopath{.editorconfig} i~\repopath{.globalconfig}, projekt pakietu NuGet, manifest. Odmawia zmiany treści wydanej wersji. Sukces: „\texttt{Generated standards <wersja>: <n> files (<w> written, <r> removed).}”. Problemy źródeł: „\texttt{<k> problem(s); nothing was generated.}”, kod 1 \\
\repopath{generate -{}-check} & & Niczego nie zapisuje; kończy się błędem, gdy \path{generated/} różni się od tego, co dają źródła (to uruchamia CI). Zgodność: „\texttt{generated/ is up to date (<n> files, standards <wersja>).}”. Różnice: każdy plik jako \texttt{missing}, \texttt{stale} albo \texttt{orphaned} i~„\texttt{generated/ is out of date (<k> file(s)). Run: secondkey-standards generate}”, kod 1 \\
\repopath{drift-check} & \repopath{-{}-repo}, \repopath{-{}-source}, \repopath{-{}-package-id}, \repopath{-{}-skill-name}, \repopath{-{}-pinned-version}, \repopath{-{}-latest-version} & Sprawdza repozytorium konsumenta (sekcja~\ref{sec:standardy-drift}) \\
\bottomrule
\end{xltabular}

Opcja \repopath{-{}-root} wskazuje repozytorium standardów; domyślnie jest to najbliższy
katalog w~bieżącym lub wyższym, który zawiera \path{catalog/pack.json}. Gdy takiego nie ma,
polecenie kończy się komunikatem „\texttt{no catalog/pack.json found at or above <katalog>;
pass -{}-root}” i~kodem 2. Błąd parsowania argumentów daje „\texttt{error: <komunikat>}”
i~„\texttt{Run with -{}-help for usage.}” oraz kod 2.

Według \path{ExitCodes.cs} oba kody niezerowe kończą krok CI błędem, ale ponawiać warto
tylko kod 2 (polecenie nie mogło się wykonać: złe argumenty, brakujący plik, nieosiągalne
źródło). \repopath{drift-check} z~\repopath{-{}-repo} wskazującym nieistniejący katalog
kończy się komunikatem „\texttt{the repository to check, <ścieżka>, does not exist}”
i~kodem 2.

W~GitHub Actions (\repopath{GITHUB\_ACTIONS=true}) każdy problem jest dodatkowo wypisywany
jako komenda workflow \texttt{::error file=<ścieżka>,line=<wiersz>::<komunikat>}, a~błąd
ogólny jako \texttt{::error::<komunikat>}; dane komunikatu mają zakodowane \repopath{\%},
CR i~LF, a~wartości właściwości dodatkowo \repopath{:} i~\repopath{,}.

\subsubsection{Walidacja catalog/pack.json}

\path{catalog/pack.json} jest pisany ręcznie i~niesie to, co dotyczy pakietu jako całości:
nazwę (\repopath{secondkey-standards}), wersję (\repopath{0.1.0}), adres repozytorium, adres
konstytucji, konfigurację skilla i~pakietu oraz zamknięte słowniki. Jest czytany ściśle:
nieznana właściwość, komentarz albo przecinek na końcu listy są błędem. Plik z~defektem
kończy \repopath{validate} i~\repopath{generate} kodem 2 (\repopath{drift-check} go nie czyta). Reguły walidacji:

\begin{itemize}
  \item \repopath{name} wymagane; \repopath{version} w~postaci \repopath{MAJOR.MINOR.PATCH}
        bez sufiksu pre-release; \repopath{repository} i~\repopath{constitution} to adresy
        \repopath{https}.
  \item \repopath{skill.name}: małe litery, cyfry i~pojedyncze myślniki, najwyżej 64 znaki,
        bez słów zastrzeżonych przez specyfikację skilli; \repopath{skill.description}: od 1
        do 1024 znaków, bez \repopath{<} i~\repopath{>} (opis skilla nie może zawierać
        znaczników XML); \repopath{skill.discovery}: \repopath{preload} albo
        \repopath{lazy}; \repopath{skill.traits} wymagane (np. \repopath{.NET|CSharp}).
  \item \repopath{package.id} zgodny ze wzorcem identyfikatora NuGet;
        \repopath{package.description} wymagany.
  \item \repopath{categories}: co najmniej jedna, bez duplikatów; każdy wpis
        \repopath{portcullisRules} ma postać \path{PORTCULLIS_<SLUG>}.
\end{itemize}

Lista \repopath{portcullisRules} w~wersji 0.1.0 zawiera trzynaście identyfikatorów: cztery
migracyjne (\path{PORTCULLIS_MIG_SYSTEM_WEB}, \path{PORTCULLIS_MIG_HTTPCONTEXT_CURRENT},
\path{PORTCULLIS_MIG_SYNC_OVER_ASYNC}, \path{PORTCULLIS_MIG_CONFIGURATION_MANAGER}) i~dziewięć
architektonicznych (\path{PORTCULLIS_P2_KERNEL_LOC_CEILING},
\path{PORTCULLIS_P2_KERNEL_ENTITY_REFERENCE}, \path{PORTCULLIS_P4_ENSURE_CREATED_OUTSIDE_TEST},
\path{PORTCULLIS_P4_SEED_DATA_IN_MODEL}, \path{PORTCULLIS_P9_CONTROLLER_NO_DBCONTEXT},
\path{PORTCULLIS_P9_ORPHAN_ENTITY}, \path{PORTCULLIS_P10_CUSTOM_BASE_CLASS},
\path{PORTCULLIS_P11_VENDOR_SDK_OUTSIDE_ADAPTER}, \path{PORTCULLIS_P15_MISSING_SERVICE_DEFAULTS}).
Reguły korzystają z~dziewięciu z~nich; na cztery pozostałe -- oba identyfikatory P2,
\path{PORTCULLIS_P9_ORPHAN_ENTITY} i~\path{PORTCULLIS_P11_VENDOR_SDK_OUTSIDE_ADAPTER} -- nie
mapuje się żadna reguła. \path{standards/README.md} mówi to wprost: lista jest słownikiem
wszystkich reguł, które bramka dostarcza, zmapowanych albo nie; diagnostyka, której nie
mapuje żadna reguła, nie trafia do generowanej konfiguracji, więc obowiązuje dla niej
domyślna severity Portcullis.

\subsubsection{Przebieg generowania}

\repopath{generate} i~\repopath{generate -{}-check} dzielą jedną ścieżkę kodu
(\repopath{PackGenerator.Build}) i~różnią się tylko tym, co robią z~wynikiem, więc nie mogą
się nie zgadzać co do tego, jaki wynik jest poprawny:

\begin{enumerate}
  \item Wczytanie reguł: wszystkie pliki \path{standards/*.md} najwyższego poziomu poza
        \path{README.md}, w~porządku porządkowym nazw; parsowanie każdego pliku; sprawdzenia
        obejmujące cały zbiór; sortowanie po identyfikatorze.
  \item Wczytanie szablonu \path{catalog/skill.template.md}: nieznany placeholder jest
        błędem z~numerem wiersza (agent przeczytałby go dosłownie), brak wymaganych
        \repopath{\{\{rule-table\}\}} i~\repopath{\{\{gate-table\}\}} także.
  \item Wczytanie i~parsowanie \path{CHANGELOG.md}. Brak \path{catalog/skill.template.md}
        albo \path{CHANGELOG.md} jest problemem („\texttt{this source file is missing}”).
  \item Jeśli jest jakikolwiek problem -- koniec, nic nie jest generowane.
  \item Zbudowanie całego drzewa w~pamięci: skill, konfiguracje analyzerów, pakiet; następnie
        skrót treści i~manifest.
  \item Bramka wersji względem zatwierdzonego \path{generated/manifest.json}
        (sekcja~\ref{sec:standardy-wersje}).
  \item \repopath{generate} zapisuje pliki, które się różnią, usuwa osierocone i~puste
        katalogi; \repopath{-{}-check} tylko porównuje.
\end{enumerate}

Szczegóły emisji:

\begin{itemize}
  \item \textbf{Skill} (\repopath{SkillEmitter}). \repopath{SKILL.md} to front matter
        (\repopath{name}; \repopath{description} jako jedna linia w~podwójnym cudzysłowie YAML;
        \repopath{metadata} z~\repopath{discovery}, \repopath{traits}, \repopath{version},
        \repopath{source}), wyrenderowany szablon i~stopka „\texttt{Second Key standards
        <wersja>, generated from standards/ … Do not edit this copy: take a newer version from
        the source repository instead.}” z~linkiem do katalogu \path{standards/} pod tagiem
        \repopath{v<wersja>}. Placeholdery szablonu: \repopath{\{\{version\}\}},
        \repopath{\{\{rule-count\}\}}, \repopath{\{\{rule-table\}\}} (kolumny \emph{Rule},
        \emph{Severity}, \emph{What migrated code must do}, \emph{Gate diagnostic}; identyfikator
        linkuje \path{references/<id>.md}), \repopath{\{\{gate-table\}\}} (kolumny
        \emph{Diagnostic}, \emph{Severity}, \emph{Rule}), \repopath{\{\{flag-register\}\}}
        (\path{docs/migration/behaviour-flags.md}) i~\repopath{\{\{repository\}\}}. Każdy plik
        \path{references/<id>.md} ma nagłówek „\texttt{\# <id> — <tytuł>}”, tabelę metadanych
        (\emph{Severity}, \emph{Category}, \emph{Applies to}, \emph{Gate diagnostic},
        \emph{Principle} z~linkiem do kotwicy zasady w~konstytucji), wiersz
        „\texttt{Contents:}” z~nagłówkami sekcji reguły, treść reguły bez zmian i~stopkę
        z~linkiem do pliku źródłowego pod tagiem wersji; taki link działa dopiero wtedy,
        gdy wersja zostanie otagowana.
  \item \textbf{Konfiguracja analyzerów} (\repopath{AnalyzerConfigEmitter}).
        \repopath{.editorconfig} ma dla każdego wzorca \repopath{appliesTo} reguł z~diagnostyką
        sekcję \repopath{[wzorzec]} z~wierszami \texttt{\# <id> <tytuł>}
        i~\texttt{dotnet\_diagnostic.<id>.severity = <waga>}. \repopath{.globalconfig} ma
        \texttt{is\_global = true}, \texttt{global\_level = 0} i~te same wiersze dla wszystkich
        reguł z~diagnostyką. W~wersji 0.1.0 wszystkie reguły z~diagnostyką mają
        \repopath{appliesTo} \repopath{"*.cs"}, więc \repopath{.editorconfig} ma jedną sekcję
        \repopath{[*.cs]} z~trzynastoma wierszami.
  \item \textbf{Pakiet NuGet} (\repopath{PackageEmitter}): projekt
        \path{generated/nuget/SecondKey.Standards.csproj}, plik
        \path{build/SecondKey.Standards.props} i~\path{README.md} pakietu (instalacja, tabela
        wag, nadpisywanie) -- szczegóły w~sekcji~\ref{sec:standardy-wersje}.
  \item \textbf{Manifest} (\repopath{Manifest}): \repopath{name}, \repopath{version},
        \repopath{contentHash}, \repopath{skill}, \repopath{package} i~lista \repopath{rules}
        (\repopath{id}, \repopath{title}, \repopath{severity}, \repopath{category},
        \repopath{diagnostics}, \repopath{principle}, \repopath{hash}). Skrót to
        \repopath{sha256:} i~pierwsze 16 znaków szesnastkowych SHA-256. Skrót reguły liczony
        jest z~identyfikatora, tytułu, wagi, kategorii, \repopath{appliesTo}, diagnostyk,
        zasady i~treści, więc każda edycja reguły go zmienia. \repopath{contentHash} obejmuje
        każdy wygenerowany plik poza manifestem, wyrenderowany z~wersją zastąpioną stałą
        \repopath{0.0.0}: pliki wymieniają wersję (\repopath{<Version>} pakietu, metadane
        skilla), więc bez tego każde podniesienie wersji wyglądałoby jak zmiana treści.
\end{itemize}

Drzewo w~pamięci przyjmuje wyłącznie ścieżki pod \path{generated/}, każdą raz; treść jest
normalizowana do LF i~jednego końcowego znaku nowej linii, a~pliki są zapisywane w~UTF-8 bez
BOM. Porównanie jest bajtowe: kopia z~CRLF albo z~BOM jest \texttt{stale}. W~wersji 0.1.0
drzewo liczy 27 plików: \repopath{SKILL.md}, 20 plików \path{references/}, dwa pliki
konfiguracji, trzy pliki pakietu i~manifest.

\zrodla{\path{letsgolegacy.secondkey-standards/docs/adr/0002-generator-as-a-dotnet-tool.md},
\path{letsgolegacy.secondkey-standards/README.md},
\path{letsgolegacy.secondkey-standards/src/SecondKey.Standards.Generator/SecondKey.Standards.Generator.csproj},
\path{letsgolegacy.secondkey-standards/Directory.Build.props},
\path{letsgolegacy.secondkey-standards/global.json},
\path{letsgolegacy.secondkey-standards/.editorconfig},
\path{letsgolegacy.secondkey-standards/.gitattributes},
\path{letsgolegacy.secondkey-standards/catalog/pack.json},
\path{letsgolegacy.secondkey-standards/src/SecondKey.Standards.Generator/Cli/CommandLineApp.cs},
\path{letsgolegacy.secondkey-standards/src/SecondKey.Standards.Generator/Cli/Reporter.cs},
\path{letsgolegacy.secondkey-standards/src/SecondKey.Standards.Generator/Cli/ExitCodes.cs},
\path{letsgolegacy.secondkey-standards/src/SecondKey.Standards.Generator/Model/PackConfig.cs},
\path{letsgolegacy.secondkey-standards/src/SecondKey.Standards.Generator/Generation/PackGenerator.cs},
\path{letsgolegacy.secondkey-standards/src/SecondKey.Standards.Generator/Generation/SkillEmitter.cs},
\path{letsgolegacy.secondkey-standards/src/SecondKey.Standards.Generator/Generation/SkillTemplate.cs},
\path{letsgolegacy.secondkey-standards/src/SecondKey.Standards.Generator/Generation/AnalyzerConfigEmitter.cs},
\path{letsgolegacy.secondkey-standards/src/SecondKey.Standards.Generator/Generation/Manifest.cs},
\path{letsgolegacy.secondkey-standards/src/SecondKey.Standards.Generator/Generation/GeneratedTree.cs},
\path{letsgolegacy.secondkey-standards/src/SecondKey.Standards.Generator/Generation/TreeSync.cs},
\path{letsgolegacy.secondkey-standards/generated/}}

\subsection{Drzewo generowane, wersjonowanie i~pakiet NuGet (ADR 0003)}\label{sec:standardy-wersje}

Reguły są kompilowane do wszystkiego, co musi się co do nich zgadzać. Skompilowany wynik ma
być widoczny w~przeglądzie pull requestu, niemożliwy do rozjechania ze źródłem
i~wersjonowany tak, żeby konsument mógł go przypiąć (C4-c).

\subsubsection{Decyzje ADR 0003}

\begin{enumerate}
  \item \textbf{Wynik jest commitowany w~\path{generated/} i~nigdy nie jest edytowany
        ręcznie.} Recenzent widzi w~tym samym diffie, co zmiana reguły robi ze skillem
        i~z~konfiguracją bramki; konsument może skopiować skill z~tagu bez uruchamiania
        czegokolwiek. Tak samo działa warstwa pakowania w~\path{architecture-standards}.
\end{enumerate}

\begin{tabularx}{\linewidth}{@{}>{\raggedright\arraybackslash}p{7.2cm} >{\raggedright\arraybackslash}X@{}}
\toprule
Ścieżka & Zawartość \\
\midrule
\path{generated/skills/<skill>/SKILL.md}, \path{references/<id>.md} & Katalog Agent Skills \\
\path{generated/config/.editorconfig} & \repopath{dotnet\_diagnostic.<id>.severity} dla każdej reguły z~diagnostyką, w~sekcjach zawężonych do jej \repopath{appliesTo} \\
\path{generated/config/.globalconfig} & Te same wagi jako globalna konfiguracja analyzerów \\
\path{generated/nuget/} & Projekt pakietu \repopath{SecondKey.Standards}, jego \path{build/} props i~README \\
\path{generated/manifest.json} & Wersja, skrót treści i~każda reguła z~własnym skrótem \\
\bottomrule
\end{tabularx}

\begin{enumerate}
  \setcounter{enumi}{1}
  \item \textbf{\repopath{generate -{}-check} buduje całe drzewo w~pamięci i~porównuje je
        bajt po bajcie z~dyskiem}; brakujące, nieaktualne i~osierocone pliki są nazywane
        osobno. Generator jest właścicielem całego \path{generated/} -- plik, którego już nie
        produkuje, jest usuwany -- z~wyjątkiem \path{bin/} i~\path{obj/}, czyli ignorowanego
        przez git wyniku pakowania projektu pakietu. CI uruchamia \repopath{-{}-check}, a~test
        jednostkowy wykonuje to samo porównanie, więc \repopath{dotnet test} wykrywa
        nieaktualne drzewo także lokalnie.
  \item \textbf{Wersję ustawia człowiek, a~build nie pozwala o~niej zapomnieć} -- zasada
        z~\path{architecture-standards}, dostosowana do tego, że wydaniami są tagi git, a~nie
        każdy merge. Skrót treści w~manifeście jest liczony z~wersją zastąpioną stałą, więc
        zmienia się wtedy i~tylko wtedy, gdy zmienia się treść. Jeśli zatwierdzony manifest ma
        tę samą wersję, inny skrót, a~\path{CHANGELOG.md} datuje tę wersję, uruchomienie kończy
        się błędem -- zarówno zwykłe, jak i~\repopath{-{}-check} -- zanim cokolwiek zostanie
        zapisane, więc regeneracja nie może „wyprać” zmiany. Wersja, której wpis mówi
        \repopath{Unreleased}, nie została otagowana, więc nikt nie mógł jej przypiąć, a~jej
        treść może się jeszcze zmieniać. Wersje nigdy nie maleją, a~każda wersja potrzebuje
        wpisu w~changelogu, bo to ten wpis drift check pokazuje konsumentowi, który jest
        w~tyle.
  \item \textbf{Znaczenie podniesienia wersji} (README, „Versioning”) -- tabela niżej. Numeru
        nie da się wyprowadzić z~diffu: tylko człowiek może powiedzieć, czy przeredagowany
        akapit odwraca regułę.
  \item \textbf{Pakiet jest wyłącznie konfiguracją i~ładuje się przez props w~\path{build/}.}
        \path{SecondKey.Standards} nie niesie assembly (\path{IncludeBuildOutput=false}),
        jest zależnością deweloperską (nie trafia do zależności runtime konsumenta), a~jego
        \path{build/SecondKey.Standards.props} dodaje globalną konfigurację do elementu
        \repopath{EditorConfigFiles} kompilatora -- tego samego, którego .NET SDK używa dla
        własnych konfiguracji poziomu analizy. Konsument dostaje wszystkie wagi, dodając jedną
        referencję pakietu, i~może go wyłączyć przez
        \path{SecondKeyStandardsGlobalConfig=false}.
  \item \textbf{\repopath{global\_level = 0}.} Konfiguracje poziomu analizy SDK używają
        poziomów ujemnych, własny \repopath{.globalconfig} repozytorium ma domyślnie 100,
        a~każdy wpis \repopath{.editorconfig} wygrywa z~każdą konfiguracją globalną. Poziom 0
        stawia standardy nad domyślnymi ustawieniami SDK i~pod wszystkim, co ustawi konsument.
        Sprawdzone na buildzie konsumenta: \repopath{CA1310 = warning} z~pakietu obowiązuje;
        \repopath{.globalconfig} konsumenta z~\repopath{none} wygrywa bez ostrzeżenia
        o~konflikcie; \repopath{.editorconfig} konsumenta z~\repopath{error} wygrywa
        z~obydwoma.
  \item \textbf{Pakiet jest pakowany i~weryfikowany w~CI, a~nie publikowany.} Job
        \repopath{package} pakuje go, buduje na nim jednorazowego konsumenta (bez pakietu CA1310
        milczy; z~pakietem jest ostrzeżeniem; z~wyłączeniem znów milczy) i~wysyła pliki
        \repopath{.nupkg} jako artefakt workflow. \path{scripts/verify-package.sh} to ten sam
        job do uruchomienia lokalnie.
  \item \textbf{Tag musi być wydaniem.} Na tagu \repopath{v*} CI sprawdza, że tag jest równy
        \repopath{v} + wersja pakietu, że \path{CHANGELOG.md} datuje tę wersję i~że
        zatwierdzony manifest ma tę wersję. Drift check ufa tagom, więc tagi są pilnowane.
\end{enumerate}

\begin{tabularx}{\linewidth}{@{}>{\raggedright\arraybackslash}p{1.6cm} >{\raggedright\arraybackslash}X@{}}
\toprule
Zmiana & Kiedy (pytanie konsumenta: ile kosztuje mnie aktualizacja?) \\
\midrule
\textbf{major} & Reguła jest wycofana lub odwrócona albo waga obniżona: praca wykonana pod starą regułę może wymagać powtórzenia \\
\textbf{minor} & Reguła dodana albo waga podniesiona: nic już zbudowanego nie jest błędne, jest więcej do spełnienia \\
\textbf{patch} & Nic wymaganego się nie zmienia: sformułowania, przykłady, linki \\
\bottomrule
\end{tabularx}

\subsubsection{Bramka wersji i~format CHANGELOG}

Bramka wersji (\repopath{VersionGate}) zgłasza problem, gdy:
\begin{itemize}
  \item \path{CHANGELOG.md} nie ma wpisu \texttt{\#\# [<wersja>]} dla wersji z~\path{catalog/pack.json};
  \item wersja w~\path{catalog/pack.json} jest niższa niż w~zatwierdzonym manifeście
        („versions only go up”);
  \item wersja jest równa zatwierdzonej, skrót treści się różni, a~wpis w~changelogu ma datę
        wydania -- komunikat każe podnieść wersję i~dodać dla niej wpis \repopath{Unreleased}.
\end{itemize}
Gdy zatwierdzonego \path{generated/manifest.json} nie ma (pierwsze generowanie) albo nie da
się go odczytać, bramka sprawdza tylko wpis w~changelogu. Podniesienie wersji bez zmiany
treści jest dozwolone (zwykłe ponowne wydanie).

Changelog ma jeden wpis na wersję, najnowszy pierwszy, z~nagłówkiem
\texttt{\#\# [MAJOR.MINOR.PATCH] — YYYY-MM-DD} albo \texttt{\#\# [MAJOR.MINOR.PATCH] —
Unreleased}; parser przyjmuje jako separator pauzę, półpauzę albo dywiz. Nagłówek w~innej
postaci, status, który nie jest ani datą, ani \repopath{Unreleased}, oraz wersja
powtórzona lub niepoprawnie uporządkowana są błędami.

\subsubsection{Wydawanie}

Według README: w~\path{CHANGELOG.md} zastępuje się \repopath{Unreleased} datą, uruchamia
\repopath{generate}, scala zmianę i~taguje \repopath{v<wersja>} na \repopath{main}. CI na
tagu sprawdza, że tag odpowiada wersji pakietu i~changelogowi (job \repopath{release-guard},
sekcja~\ref{sec:standardy-ci}). Po wydaniu każda zmiana treści wymaga nowej wersji. Wersja
0.1.0 ma w~changelogu status \repopath{Unreleased}; żaden tag nie istnieje, więc bramka
wersji nie zamraża jeszcze niczego.

\subsubsection{Pakiet SecondKey.Standards u~konsumenta}

Pakiet jest pakowany przez CI i~dołączany do każdego uruchomienia jako artefakt workflow
\repopath{secondkey-standards-packages}; nie jest publikowany w~nuget.org. Aby go użyć,
umieszcza się plik \repopath{.nupkg} w~lokalnym lub wewnętrznym feedzie (albo pakuje
samodzielnie przez \path{scripts/verify-package.sh}), a~następnie:

\begin{lstlisting}
<!-- Directory.Packages.props, with central package management: every project gets it -->
<GlobalPackageReference Include="SecondKey.Standards" Version="0.1.0" />
\end{lstlisting}

albo, w~każdym projekcie lub w~\path{Directory.Build.props}:

\begin{lstlisting}
<PackageReference Include="SecondKey.Standards" Version="0.1.0" PrivateAssets="all" />
\end{lstlisting}

Diagnostyki zgłaszają analyzery, do których należą: Portcullis identyfikatory
\repopath{PORTCULLIS\_}, .NET SDK identyfikatory \repopath{CA}; pakiet ustawia ich wagę, ale
nie dodaje analyzerów. Każdą wagę można nadpisać w~\repopath{.editorconfig} konsumenta, który
zawsze wygrywa, np.:

\begin{lstlisting}
[*.cs]
dotnet_diagnostic.CA1305.severity = suggestion
\end{lstlisting}

\path{<SecondKeyStandardsGlobalConfig>false</SecondKeyStandardsGlobalConfig>} wyłącza pakiet
w~projekcie, a~props ustawia też właściwość MSBuild \repopath{SecondKeyStandardsVersion}
z~wersją standardów. Bez pakietu są dwie drogi. \path{generated/config/.globalconfig}
kopiuje się do katalogu głównego repozytorium jako \path{.globalconfig}; jako globalna
konfiguracja analyzerów obowiązuje dla każdego pliku każdego projektu, który ją ładuje, bez
zawężenia do \repopath{appliesTo}. Sekcje \path{generated/config/.editorconfig} -- zawężone
do wzorców \repopath{appliesTo} reguł -- scala się z~własnym \path{.editorconfig}
repozytorium; plik nie ma \texttt{root = true}, a~wpisy \path{.editorconfig} mają
pierwszeństwo przed każdą globalną konfiguracją analyzerów, także tą z~pakietu.
Diagnostyki Portcullis są skonfigurowane, zanim ich
analyzery pojawią się w~buildzie konsumenta; waga \repopath{.editorconfig} dla nieznanego
identyfikatora diagnostyki nie ma żadnego efektu, więc konfigurację można przyjąć przed bramką.

Projekt pakietu: \repopath{netstandard2.0}, \repopath{IncludeBuildOutput} \repopath{false},
\repopath{DevelopmentDependency} \repopath{true}, \path{SuppressDependenciesWhenPacking}
\repopath{true}, \repopath{NoDefaultExcludes} \repopath{true} (źródło globalnej konfiguracji
to plik z~kropką, który NuGet inaczej pomija). Do pakietu trafiają
\path{build/SecondKey.Standards.props}, \path{build/SecondKey.Standards.globalconfig} (kopia
\path{generated/config/.globalconfig}), \path{README.md} i~\path{LICENSE}. Pakowanie ręczne:
\texttt{dotnet pack generated/nuget/SecondKey.Standards.csproj -o artifacts/packages}.
Metadane pakietu: \repopath{PackageId} i~\repopath{Version} z~\path{catalog/pack.json}
(wersja pakietu jest wersją standardów), \repopath{Description}
z~\repopath{package.description}, \repopath{PackageTags}
\path{secondkey;standards;migration;analyzers;globalconfig;editorconfig},
\repopath{PackageProjectUrl} -- adres repozytorium, \repopath{PackageReadmeFile}
\repopath{README.md}, \repopath{PackageLicenseFile} \repopath{LICENSE},
\repopath{IsPackable} \repopath{true}.

Skrypt \path{scripts/verify-package.sh} (opcjonalny argument: katalog wyjściowy, domyślnie
\path{artifacts/packages}) pakuje pakiet, sprawdza jego zawartość (obecne
\path{build/<id>.props}, \path{build/<id>.globalconfig}, \path{README.md}, \path{LICENSE};
brak katalogu \path{lib/}), tworzy w~katalogu tymczasowym poza repozytorium projekt
konsumenta z~jednym plikiem wywołującym \repopath{string.Compare(left, right)} i~sprawdza
trzy przypadki: bez pakietu CA1310 nie jest zgłaszany; z~pakietem jest zgłaszany jako
\repopath{warning}; z~\path{-p:SecondKeyStandardsGlobalConfig=false} znów nie jest zgłaszany.
CA1310 zastępuje każdą przypisaną diagnostykę: analyzery Portcullis są osobnym pakietem, ale
wagi docierają do kompilatora tą samą konfiguracją globalną. CI sprawdza więc tylko CA1310;
współpracy z~pakietem analyzerów Portcullis nie sprawdza.

Skrypt działa tylko w~klonie git repozytorium (\texttt{git rev-parse -{}-show-toplevel}),
czyta identyfikator i~wersję pakietu z~\path{catalog/pack.json}, a~bez
\path{generated/nuget/SecondKey.Standards.csproj} kończy się komunikatem, żeby najpierw
uruchomić generator. Zawartość pakietu sprawdza poleceniem \repopath{unzip}, a~konsument
tymczasowy ma \path{nuget.config} z~dwoma źródłami: katalogiem wyjściowym i~nuget.org.
Wynikiem jest plik \path{<katalog>/SecondKey.Standards.<wersja>.nupkg} -- dla wersji 0.1.0
domyślnie \path{artifacts/packages/SecondKey.Standards.0.1.0.nupkg}, ten sam, który krok~7
procedury (sekcja~\ref{sec:standardy-modernize}) kopiuje do \path{.packages/}. Każda
nieudana kontrola kończy skrypt komunikatem „\texttt{verify-package: ...}” i~kodem 1.

\subsubsection{Odrzucone alternatywy}

\begin{itemize}
  \item \textbf{Generowanie tylko w~CI, bez commitowania} -- nic nie mogłoby się zestarzeć, ale
        nie byłoby diffu pokazującego, co robi zmiana reguły, ani możliwości skopiowania
        skilla z~tagu bez budowania.
  \item \textbf{Wyprowadzanie wersji z~diffu} -- narzędzie, które zgaduje, w~końcu wyda wersję
        major jako patch, co jest gorsze niż brak wersjonowania.
  \item \textbf{Dostarczanie konfiguracji jako \repopath{.editorconfig} w~pakiecie} -- NuGet
        nie umieści pliku \repopath{.editorconfig} w~drzewie konsumenta; globalna konfiguracja
        dodana przez props builda jest mechanizmem, który MSBuild obsługuje.
\end{itemize}

\zrodla{\path{letsgolegacy.secondkey-standards/docs/adr/0003-generated-tree-and-versioning.md},
\path{letsgolegacy.secondkey-standards/README.md},
\path{letsgolegacy.secondkey-standards/CHANGELOG.md},
\path{letsgolegacy.secondkey-standards/src/SecondKey.Standards.Generator/Versioning/VersionGate.cs},
\path{letsgolegacy.secondkey-standards/src/SecondKey.Standards.Generator/Versioning/Changelog.cs},
\path{letsgolegacy.secondkey-standards/src/SecondKey.Standards.Generator/Generation/PackGenerator.cs},
\path{letsgolegacy.secondkey-standards/src/SecondKey.Standards.Generator/Generation/PackageEmitter.cs},
\path{letsgolegacy.secondkey-standards/src/SecondKey.Standards.Generator/Generation/AnalyzerConfigEmitter.cs},
\path{letsgolegacy.secondkey-standards/generated/nuget/README.md},
\path{letsgolegacy.secondkey-standards/generated/nuget/SecondKey.Standards.csproj},
\path{letsgolegacy.secondkey-standards/generated/nuget/build/SecondKey.Standards.props},
\path{letsgolegacy.secondkey-standards/generated/config/.globalconfig},
\path{letsgolegacy.secondkey-standards/generated/config/.editorconfig},
\path{letsgolegacy.secondkey-standards/scripts/verify-package.sh},
\path{letsgolegacy.secondkey-standards/docs/USING-WITH-MODERNIZE-DOTNET.md}}

\subsection{Procedura: zmiana standardu od edycji do wydania}\label{sec:standardy-procedura-zmiany}

Kroki zebrane z~README, \path{AGENTS.md}, \path{standards/README.md}, ADR 0003, szablonu
pull requestu i~joba \repopath{release-guard}:

\begin{enumerate}
  \item Przygotować klon: \texttt{./scripts/setup.sh} (.NET~10 SDK, skaner sekretów, hook
        pre-commit).
  \item Edytować plik reguły albo dodać nowy plik \path{standards/<id>-<slug>.md}
        z~nowym, unikalnym identyfikatorem. Dane pakietu jako całości zmienia się
        w~\path{catalog/pack.json}, treść skilla -- w~\path{catalog/skill.template.md}.
        Diagnostykę Portcullis dopisuje się do \repopath{portcullisRules}
        w~\path{catalog/pack.json} dopiero wtedy, gdy bramka ją dostarcza.
  \item \texttt{dotnet run -{}-project src/SecondKey.Standards.Generator -{}- validate} musi
        przejść.
  \item Podjąć decyzję o~wersji według tabeli podniesień (major, minor, patch;
        sekcja~\ref{sec:standardy-wersje}). Jeśli bieżąca wersja ma w~\path{CHANGELOG.md}
        datę wydania, każda zmiana wygenerowanej treści wymaga nowej wersji: podnieść
        \repopath{version} w~\path{catalog/pack.json}
        i~dodać na początku changeloga wpis \texttt{\#\# [<nowa wersja>] — Unreleased}
        z~opisem zmiany. Wersja, której wpis mówi \repopath{Unreleased}, może się jeszcze
        zmieniać bez podnoszenia.
  \item Uruchomić \repopath{generate} i~zacommitować źródło razem z~\path{generated/}.
  \item Sprawdzić lokalnie to, co sprawdza CI: \repopath{dotnet test},
        \texttt{generate -{}-check}, \path{./scripts/verify-package.sh},
        \path{./scripts/scan-secrets.sh}.
  \item Otworzyć pull request według szablonu: ticket i~jego „done when”, zmiana z~powodem
        i~dowodem, na którym reguła się opiera (dokumentacja, kod albo zaobserwowane
        zachowanie), decyzja wersji (major, minor, patch albo „No standards content
        changed”) i~lista weryfikacji.
  \item Wydanie: w~\path{CHANGELOG.md} zastąpić \repopath{Unreleased} datą
        \texttt{RRRR-MM-DD}, uruchomić \repopath{generate}, scalić i~otagować
        \repopath{v<wersja>} na \repopath{main}. Job \repopath{release-guard} odrzuca tag,
        który nie równa się \repopath{v} i~wersji z~\path{catalog/pack.json} („\texttt{Tag …
        does not match catalog/pack.json version …}”), wersję bez daty w~changelogu
        („\texttt{replace 'Unreleased' with the release date before tagging}”) oraz manifest
        w~innej wersji („\texttt{regenerate before tagging}”).
\end{enumerate}

\begin{checklista}
  \item Reguła zmieniona albo dodana; \repopath{validate} przechodzi.
  \item Decyzja wersji podjęta; \path{CHANGELOG.md} ma wpis dla wersji z~\path{catalog/pack.json}.
  \item \path{generated/} zregenerowane i~zacommitowane razem ze źródłem.
  \item \repopath{dotnet test}, \repopath{generate -{}-check} i~skan sekretów przechodzą.
  \item Przy wydaniu: data w~changelogu, tag \repopath{v<wersja>} na \repopath{main}.
\end{checklista}

\zrodla{\path{letsgolegacy.secondkey-standards/README.md},
\path{letsgolegacy.secondkey-standards/AGENTS.md},
\path{letsgolegacy.secondkey-standards/standards/README.md},
\path{letsgolegacy.secondkey-standards/docs/adr/0003-generated-tree-and-versioning.md},
\path{letsgolegacy.secondkey-standards/.github/pull_request_template.md},
\path{letsgolegacy.secondkey-standards/.github/workflows/ci.yml},
\path{letsgolegacy.secondkey-standards/src/SecondKey.Standards.Generator/Versioning/VersionGate.cs},
\path{letsgolegacy.secondkey-standards/src/SecondKey.Standards.Generator/Versioning/Changelog.cs}}

\subsection{Wagi standardów a~kanały Portcullis}\label{sec:standardy-wagi-kanaly}

Wagi standardów trafiają do kompilatora: pakiet \repopath{SecondKey.Standards} dodaje
globalną konfigurację do elementu \repopath{EditorConfigFiles} buildu, a~sam nie dodaje
analyzerów -- diagnostyki \repopath{PORTCULLIS\_} zgłasza w~buildzie pakiet analyzerów
Portcullis (sekcja~\ref{sec:portcullis-dystrybucja}). Skill mówi agentowi, że bramka
Portcullis i~analyzery .NET SDK zgłaszają te diagnostyki na pull requeście migracji „at
these severities”. Źródła Portcullis mówią co innego o~kanale CLI, akcji i~kontenera:
\repopath{portcullis scan} buduje własną kompilację, raportuje każdą diagnostykę na jej
poziomie domyślnym i~nie stosuje nadpisań \texttt{dotnet\_diagnostic.<id>.severity}
(sekcja~\ref{sec:portcullis-konfiguracja}); to samo dotyczy SARIF z~tego skanu
i~\texttt{sk gate}, który czyta ten SARIF. Cztery diagnostyki \repopath{CA} w~ogóle nie
należą do Portcullis: skan uruchamia tylko jego analyzery, więc ich nie ma ani w~SARIF, ani
w~\texttt{sk gate}; zgłasza je wyłącznie kompilator, i~to dopiero po włączeniu (domyślnie
są wyłączone, a~włącza je m.in. konfiguracja standardów). Dla pięciu z~dziewięciu
zmapowanych diagnostyk Portcullis poziom domyślny różni się od wagi standardu:

\begin{xltabular}{\linewidth}{@{}>{\raggedright\arraybackslash}X >{\raggedright\arraybackslash}p{2.75cm} >{\raggedright\arraybackslash}p{2.25cm} >{\raggedright\arraybackslash}p{2.25cm}@{}}
\toprule
Diagnostyka & Reguła & Waga standardu & Poziom domyślny Portcullis \\
\midrule
\endfirsthead
\toprule
Diagnostyka & Reguła & Waga standardu & Poziom domyślny Portcullis \\
\midrule
\endhead
\path{PORTCULLIS_P4_ENSURE_CREATED_OUTSIDE_TEST} & SK-ARCH-001 & \repopath{error} & \repopath{error} \\
\path{PORTCULLIS_P4_SEED_DATA_IN_MODEL} & SK-ARCH-002 & \repopath{warning} & \repopath{warning} \\
\path{PORTCULLIS_P9_CONTROLLER_NO_DBCONTEXT} & SK-ARCH-005 & \repopath{warning} & \repopath{error} \\
\path{PORTCULLIS_P10_CUSTOM_BASE_CLASS} & SK-ARCH-006 & \repopath{suggestion} & \repopath{warning} \\
\path{PORTCULLIS_P15_MISSING_SERVICE_DEFAULTS} & SK-ARCH-007 & \repopath{suggestion} & \repopath{warning} \\
\path{PORTCULLIS_MIG_SYSTEM_WEB} & SK-MIG-001 & \repopath{error} & \repopath{warning} \\
\path{PORTCULLIS_MIG_HTTPCONTEXT_CURRENT} & SK-MIG-002 & \repopath{error} & \repopath{error} \\
\path{PORTCULLIS_MIG_SYNC_OVER_ASYNC} & SK-MIG-003 & \repopath{error} & \repopath{warning} \\
\path{PORTCULLIS_MIG_CONFIGURATION_MANAGER} & SK-MIG-004 & \repopath{error} & \repopath{error} \\
\bottomrule
\end{xltabular}

Bramka CLI blokuje tylko finding o~poziomie \repopath{error} -- także ostrzeżenie
podniesione do \repopath{error} na linii przypisanej AI -- w~zakresie bramki
i~nieprzyjęty przez baseline (sekcja~\ref{sec:portcullis-spec}). W~skanie CLI
\path{PORTCULLIS_P9_CONTROLLER_NO_DBCONTEXT} blokuje więc w~zakresie bramki, choć
SK-ARCH-005 każe zapisać przeniesiony \repopath{DbContext} jako dług;
\path{PORTCULLIS_P10_CUSTOM_BASE_CLASS} i~\path{PORTCULLIS_P15_MISSING_SERVICE_DEFAULTS}
mogą zablokować po eskalacji, choć sugestia według standardów nigdy nie blokuje;
\path{PORTCULLIS_MIG_SYSTEM_WEB} i~\path{PORTCULLIS_MIG_SYNC_OVER_ASYNC} blokują tylko po
eskalacji, choć standard daje im \repopath{error}. Wagi z~tabel skilla i~pakietu oraz
z~opisów reguł (sekcje~\ref{sec:standardy-arch} i~\ref{sec:standardy-mig}) obowiązują
w~buildzie, który ma pakiet analyzerów Portcullis i~konfigurację standardów. Źródła
standardów tej różnicy nie opisują.

\zrodla{\path{letsgolegacy.secondkey-standards/catalog/skill.template.md},
\path{letsgolegacy.secondkey-standards/generated/nuget/README.md},
\path{letsgolegacy.secondkey-standards/generated/nuget/build/SecondKey.Standards.props},
\path{letsgolegacy.secondkey-standards/standards/README.md},
\path{letsgolegacy.portcullis/README.md},
\path{letsgolegacy.portcullis/docs/CONFIGURATION.md},
\path{letsgolegacy.portcullis/docs/SPEC.md},
\path{letsgolegacy.portcullis/src/Portcullis.Engine/Scanner.cs},
\path{letsgolegacy.secondkey/src/SecondKey.Evidence/Sarif/SarifLog.cs}}

\subsection{Skill dla GitHub Copilot upgrade (ADR 0004)}\label{sec:standardy-skill}

Agentem migracyjnym w~demonstracji Second Key jest agent modernizacyjny Microsoft dla .NET --
w~backlogu \emph{modernize-dotnet}, obecnie dokumentowany jako \textbf{GitHub Copilot
upgrade}. Visual Studio nadal instaluje go jako komponent „GitHub Copilot app modernization”
i~wywołuje jako \repopath{@Modernize}; w~pozostałych środowiskach jest to agent
\repopath{@upgrade}. Agent przyjmuje własne skille, które kodują standardy zespołu. Ticket R4
pakuje standardy jako taki skill; ticket P6 benchu (planowany) ma uruchomić agenta z~tym
skillem (rozdział~\ref{sec:bench}).

\subsubsection{Fakty sprawdzone w~źródłach}

Mechanizm sprawdzono w~bieżących źródłach zamiast go zakładać (odczyt 2026-09-26):

\begin{xltabular}{\linewidth}{@{}>{\raggedright\arraybackslash}X >{\raggedright\arraybackslash}p{5.6cm}@{}}
\toprule
Fakt & Źródło \\
\midrule
\endfirsthead
\toprule
Fakt & Źródło \\
\midrule
\endhead
Skill to plik Markdown z~front matter \repopath{name}, \repopath{description} i~opcjonalnymi \repopath{metadata.discovery} (\repopath{preload}, \repopath{lazy}, \repopath{scenario}) oraz \repopath{metadata.traits} & \repopath{dotnet/docs}: \path{docs/core/porting/github-copilot-upgrade/customization.md} (2026-09-21) \\
Lokalizacje i~priorytet: \texttt{\%UserProfile\%/.copilot/skills/} > \path{.github/upgrades/skills/} > \path{.github/skills/} > wbudowane & jw. \\
Instrukcje sformułowane „From now on…” są zapisywane do \path{.github/upgrades/<scenarioId>/scenario-instructions.md} i~ładowane do każdej decyzji & jw. \\
Skille dostarczane przez Microsoft to katalogi \path{<name>/SKILL.md} z~plikami referencyjnymi; nazwy zaczynają się od gerundium; opisy do 1024 znaków, w~trzeciej osobie, bez znaczników XML; \repopath{SKILL.md} poniżej 500 wierszy & \path{microsoft/upgrade-agent-plugins} (jego skill \repopath{creating-skills} i~reguły walidacji) \\
Skille projektu mieszkają w~\path{.github/skills}, \path{.claude/skills} albo \path{.agents/skills}, po jednym katalogu na skill, z~plikiem \repopath{SKILL.md}; działają w~agencie chmurowym, CLI, aplikacji Copilot i~VS Code & \repopath{github/docs}: \emph{About agent skills}, \emph{Adding agent skills} \\
Format katalogu i~ograniczenia nazwy oraz opisu & specyfikacja Agent Skills (\path{agentskills/agentskills}) \\
\bottomrule
\end{xltabular}

\subsubsection{Decyzje ADR 0004}

\begin{enumerate}
  \item \textbf{Katalog Agent Skills, generowany, kopiowalny bez zmian.}
        \path{generated/skills/<name>/} zawiera \repopath{SKILL.md}
        i~\path{references/<rule-id>.md}; każdy link w~środku jest względny wobec katalogu albo
        bezwzględny, co sprawdza test. Instalacja to skopiowanie jednego katalogu do
        \path{.github/skills/} repozytorium docelowego. Skill \textbf{nie} jest umieszczony we
        własnym \path{.github/skills/} tego repozytorium: to repozytorium nie jest migrowane,
        a~skill w~tym miejscu byłby ładowany przez każdego agenta pracującego nad samymi
        standardami.
  \item \textbf{Nazwa \path{applying-dotnet-migration-standards}.} Przykładem w~tickecie
        była \path{dotnet-migration-standards}; reguły autorskie Microsoft dla skilli
        upgrade wymagają gerundium na początku, a~nazwa jest jednocześnie katalogiem,
        poleceniem \repopath{/skill} i~kluczem routingu, więc idzie za konwencją agenta.
  \item \textbf{\texttt{metadata.discovery: preload}.} Standardy są przekrojowe: każde
        zadanie migracji dotyka zachowania, a~skill ładowany leniwie jest ładowany tylko
        wtedy, gdy sformułowanie zadania pasuje do jego opisu. Microsoft zaleca
        \repopath{lazy} dla większości własnych skilli, żeby oszczędzać kontekst; ten skill
        utrzymuje zawsze ładowaną część na poziomie około 200 wierszy, a~szczegóły każdej
        reguły trzyma w~pliku referencyjnym ładowanym na żądanie -- to jest stopniowe
        ujawnianie, które chroni tamto zalecenie.
  \item \textbf{Skill niesie swoją wersję} (\repopath{metadata.version},
        \repopath{metadata.source}). Skopiowany skill jest inaczej anonimowy; z~wersją drift
        check (C4-c) może powiedzieć konsumentowi, czy instrukcje jego agenta są aktualne
        i~czy zgadzają się z~konfiguracją bramki, z~którą buduje.
  \item \textbf{Treść jest pisana ręcznie, tabele są generowane.}
        \path{catalog/skill.template.md} zawiera instrukcje -- dwie reguły nadrzędne, przebieg
        pracy nałożony na etapy agenta (baseline przed pierwszym zadaniem, reguły w~każdym
        zadaniu, build, bramka i~testy po każdym, flagi na końcu), to, czego nigdy nie zmienia
        się po cichu, format flag, kryteria sukcesu i~obsługę błędów -- czyli strukturę,
        której wymagają reguły autorskie Microsoft. Tabele reguł i~bramki są generowane
        z~reguł, więc nie mogą się rozjechać. Szablon musi zawierać obie i~może używać tylko
        znanych placeholderów.
  \item \textbf{Protokół flag występuje dwa razy, a~test pilnuje, żeby był identyczny}
        (sekcja~\ref{sec:standardy-flagi}).
  \item \textbf{Podwójne zabezpieczenie w~P6.} Procedura
        (\path{docs/USING-WITH-MODERNIZE-DOTNET.md}) każe nazwać skill w~pierwszym prompcie
        i~uczynić go trwałą instrukcją, więc przebieg nie zależy wyłącznie od automatycznego
        wykrycia.
  \item \textbf{Opis to jedna linia w~podwójnym cudzysłowie, z~klauzulą „Use when…”.}
        Parsery YAML czytają blok złożony równie dobrze, ale liniowe sprawdzenia w~narzędziach
        dla skilli widzą tylko pierwszą linię skalara blokowego. Każdy plik referencyjny
        zaczyna się od wiersza \repopath{Contents:} z~sekcjami reguły, co spełnia regułę
        autorską Microsoft (spis treści w~plikach dłuższych niż 100 wierszy); wiersz jest
        w~każdym pliku, żeby nie pojawiał się i~nie znikał, gdy reguła przekroczy próg.
        Tak brzmi decyzja po poprawce ADR z~26 IX 2026; wcześniej mówiła o~wierszu
        \repopath{Contents:} tylko w~plikach dłuższych niż 100 wierszy, a~poprawka uznała za
        zamierzone to, co robi generator: wiersz w~każdym pliku referencyjnym.
  \item \textbf{Jednorazowe sprawdzenie walidatorem skilli Microsoft, poza CI.} Skrypt
        \path{validate_skill.sh} z~\path{microsoft/upgrade-agent-plugins} zgłasza 0 błędów
        i~dwa ostrzeżenia, obydwa heurystyczne i~nieadekwatne: jego lista prefiksów gerundium
        nie zawiera \repopath{applying-}, a~\repopath{\textbackslash 0} w~napisie C\#
        w~SK-MIG-005 wygląda dla niego jak ścieżka Windows. Nie jest podłączony do CI, bo to
        repozytorium jest dystrybuowane na przedpremierowych warunkach licencyjnych Microsoft,
        które ograniczają użycie do produktów Microsoft; istotne ograniczenia, które sprawdza
        (nazwa, długość i~treść opisu, \repopath{SKILL.md} poniżej 500 wierszy, referencje
        jeden poziom w~głąb), egzekwują własny generator i~testy tego repozytorium.
\end{enumerate}

\subsubsection{Front matter wygenerowanego skilla}

Wygenerowany \path{generated/skills/applying-dotnet-migration-standards/SKILL.md} (192 wiersze
w~wersji 0.1.0) zaczyna się tak:

\begin{lstlisting}
---
name: applying-dotnet-migration-standards
description: "Applies the Second Key migration standards while a .NET Framework application is upgraded to .NET 10 or another modern .NET version: preserves observable behaviour and records a behaviour flag instead of fixing a suspected defect; replaces System.Web, HttpContext.Current, ConfigurationManager and blocking on tasks; makes string comparison and culture explicit because .NET Framework uses NLS and modern .NET uses ICU; keeps EF6 query, loading and validation behaviour; keeps routes, JSON shape and cookies; names the diagnostics the Portcullis gate checks. Use when upgrading or migrating a .NET Framework solution to .NET 10, moving ASP.NET MVC or Web API to ASP.NET Core, porting EF6 to EF Core, or moving web.config settings to appsettings.json. Also use when asked to upgrade to .NET 10, migrate from .NET Framework, modernize a solution, or follow the migration standards."
metadata:
  discovery: "preload"
  traits: ".NET|CSharp"
  version: "0.1.0"
  source: "https://github.com/konradcinkusz/letsgolegacy.secondkey-standards"
---
\end{lstlisting}

\subsubsection{Treść skilla}

Treść pochodzi z~\path{catalog/skill.template.md}. Wstęp mówi agentowi, że migracja będzie
weryfikowana niezależnie: Second Key odtwarza nagrany ruch na systemie legacy i~zmigrowanym
i~porównuje każdą odpowiedź, zmianę w~bazie i~wywołanie wychodzące z~pisanym kontraktem,
a~różnica wprowadzona celowo jest zgłaszana dokładnie tak jak regresja. To opis ze skilla,
szerszy niż to, co robi rdzeń w~opisywanej wersji: \texttt{sk} nagrywa tylko ruch HTTP
wchodzący, porównuje odpowiedzi oraz zmiany wierszy wskazanych tabel SQL Server, które
zapisuje sonda przy replayu, a~wywołań wychodzących nie nagrywa ani nie porównuje
(faza~02, rozdział~\ref{sec:rdzen}). Sekcje (każdą sprawdza test
\repopath{SkillContentTests}):

\begin{description}
  \item[Two rules above all others.] (1) \emph{Preserve behaviour} -- zmiana platformy, nie
        biznesu: wyniki, zaokrąglenia, kolejność, walidacja, odpowiedzi błędów, kontrakt HTTP
        i~skutki uboczne zostają, także gdy wyglądają na błędne (SK-MIG-011, SK-MIG-012).
        (2) \emph{Flag instead of fixing} -- gdy jedyną drogą do spełnienia reguły lub
        ukończenia zadania jest zmiana zachowania albo gdy agent zauważy podejrzewany defekt,
        zachowuje zachowanie legacy i~zapisuje flagę; nigdy nie zmienia po cichu, decyduje
        człowiek, w~osobnej zmianie.
  \item[Workflow.] Lista postępu, którą agent kopiuje do notatek postępu upgrade'u i~aktualizuje
        (zestawienie niżej). \emph{Before the first task}: agent tworzy rejestr flag
        i~zapisuje, zanim zmieni się jakikolwiek kod, fakty, które znikną, gdy migracja
        przepisze pliki: źródło kultury legacy (SK-MIG-006); każdy klucz
        \repopath{<appSettings>} i~wpis \repopath{<connectionStrings>} dla każdej transformacji
        \repopath{web.config}, z~oznaczeniem tych, które zawierają dane uwierzytelniające
        (SK-MIG-004, SK-ARCH-003); stos dostępu do danych -- wersję EF6, nawigacje czytane bez
        \repopath{Include}, inicjalizatory bazy i~metody \repopath{Seed} (SK-MIG-010,
        SK-ARCH-001, SK-ARCH-002); każdy endpoint JSON i~serializer za nim (SK-MIG-012);
        projekty testowe i~każdy test, który już nie przechodzi albo jest pomijany
        (SK-MIG-013). Po napisaniu oceny (assessment) lub planu dodaje je jako kontekst, żeby
        widziało je każde zadanie. \emph{In every task}: przed edycją pliku znaleźć reguły,
        które go dotyczą, i~przeczytać ich pliki \path{references/<rule-id>.md}; wprowadzić
        zmianę pokazaną w~przykładzie \emph{Compliant} albo flagę w~sytuacjach z~sekcji
        \emph{Flag instead of fixing}; najmniejsza zmiana, która się kompiluje i~zachowuje
        zachowanie -- bez refaktoryzacji, zmiany nazw publicznych API, przestawiania logiki
        i~porządkowania kodu, bo diff migracji, który przy okazji sprząta, nie daje się
        zweryfikować. \emph{After every task}: build -- diagnostyka z~tabeli bramki jest
        znaleziskiem, naprawianym zgodnie z~regułą albo tłumionym z~uzasadnieniem nazywającym
        flagę; uruchomienie istniejących testów -- nieprzechodzący test oznacza zmianę
        zachowania do przywrócenia, nigdy edycję oczekiwań, pominięcie ani usunięcie testu
        (SK-MIG-013). \emph{At the end}: każda flaga ma komentarz w~kodzie i~wiersz w~rejestrze;
        podsumowanie i~opis pull requestu kończą się sekcją „Behaviour flags” z~liczbą flag
        na regułę i~linkiem do rejestru, znaleziska bezpieczeństwa na początku.
  \item[Never change silently.] Tryb globalizacji (SK-MIG-007); porównania napisów i~kultura,
        z~flagą dla lingwistycznych sortowań i~porównań tekstu ludzkiego (SK-MIG-005,
        SK-MIG-006); wywołania wychodzące -- bez ponowień, circuit breakerów, hedgingu i~zmian
        limitów czasu (SK-MIG-009); dostęp do danych -- lazy loading, mapowanie tabel
        i~kolumn, walidacja przy zapisie (SK-MIG-010), schemat i~brak \repopath{EnsureCreated}
        poza testami (SK-ARCH-001); konfiguracja -- każdy klucz i~wartość przeniesione,
        czytane raz przez \repopath{IOptions<T>} (SK-MIG-004), bez danych uwierzytelniających
        w~commitowanym pliku (SK-ARCH-003); testy -- oczekiwane wartości, tolerancje i~wejścia
        (SK-MIG-013).
  \item[Rules.] Generowana tabela 20 reguł (sekcja~\ref{sec:standardy-przeglad}) i~znaczenie
        wag dla agenta.
  \item[What the gate checks.] Generowana tabela 13 diagnostyk z~wagami. Bramka patrzy na
        wiersze zmienione przez pull request, więc kod przepisany przez migrację jest
        sprawdzany nawet tam, gdzie kod legacy miał ten sam wzorzec. W~skanie CLI Portcullis
        obowiązują poziomy domyślne reguł, nie wagi z~tej tabeli
        (sekcja~\ref{sec:standardy-wagi-kanaly}).
  \item[Flag format.] Format komentarza i~nagłówek rejestru (sekcja~\ref{sec:standardy-flagi}).
  \item[Success criteria.] Lista kryteriów sukcesu (zestawienie niżej).
  \item[When something goes wrong.] Reguła i~zadanie się nie zgadzają (krok dodałby
        ponowienia, włączyłby invariant globalization albo przeniósł dane seedujące do
        \repopath{HasData}) -- wygrywa reguła: zadanie bez tej części i~flaga opisująca
        konflikt. Build wymaga zmiany zachowania -- najmniejsza zmiana zachowująca zachowanie;
        jeśli jej nie ma, flaga, stłumienie tej jednej diagnostyki z~flagą i~dalsza praca.
        Test nie przechodzi z~powodu nieodłącznego od platformy (sortowanie ICU, API bez
        odpowiednika) -- test zostaje nieprzechodzący, flaga i~wzmianka w~podsumowaniu. Brak
        pewności, czy zmiana zmienia zachowanie -- założyć, że zmienia, i~oflagować.
\end{description}

Lista postępu z~sekcji \emph{Workflow} (przepisana z~wygenerowanego skilla):

\begin{checklista}
  \item Before the first task: legacy baseline recorded in \path{docs/migration/behaviour-flags.md}
  \item Every task: rules applied to the code the task touches
  \item After every task: build clean of gate errors, existing tests pass unchanged
  \item At the end: flags register complete, summary ends with „Behaviour flags”
\end{checklista}

Kryteria sukcesu z~sekcji \emph{Success criteria}:

\begin{checklista}
  \item The solution builds on .NET 10, and every test that passed before still passes, unchanged.
  \item No error-severity diagnostic from the gate table remains unsuppressed, and every
        suppression names a flag.
  \item No globalization switch, no new retry policy, no \repopath{EnsureCreated} outside tests,
        and no credential in a committed file.
  \item \path{docs/migration/behaviour-flags.md} lists every flag, and the summary ends with the
        Behaviour flags section.
\end{checklista}

\subsubsection{Konsekwencje i~otwarte punkty}

Kryterium akceptacji R4 -- „agent podejmuje skill w~tickecie P6 benchu” -- nie może zostać
sprawdzone w~tym repozytorium: nic tutaj nie potrafi uruchomić GitHub Copilot. Procedura
wymienia kontrole, które wykonuje człowiek, i~to, co ma zapisać
(sekcja~\ref{sec:standardy-modernize}). Do czasu P6 otwarte są dwa punkty: jak
\repopath{preload} zachowuje się dla skilla repozytorium w~każdym kliencie oraz rozbieżność
na stronie Microsoft, której przykładowa ścieżka jest pojedynczym plikiem
(\path{.github/skills/my-skill.md}), podczas gdy dostarczane skille Microsoft i~dokumentacja
GitHub używają formy katalogu \path{<name>/SKILL.md}, użytej tutaj.

\zrodla{\path{letsgolegacy.secondkey-standards/docs/adr/0004-skill-for-github-copilot-upgrade.md},
\path{letsgolegacy.secondkey-standards/catalog/skill.template.md},
\path{letsgolegacy.secondkey-standards/catalog/pack.json},
\path{letsgolegacy.secondkey-standards/generated/skills/applying-dotnet-migration-standards/SKILL.md},
\path{letsgolegacy.secondkey-standards/src/SecondKey.Standards.Generator/Generation/SkillEmitter.cs},
\path{letsgolegacy.secondkey-standards/tests/SecondKey.Standards.Generator.Tests/Generation/SkillContentTests.cs},
\path{letsgolegacy.secondkey/schemas/skcap.schema.json},
\path{letsgolegacy.secondkey/src/SecondKey.Replay/Probes/SqlServerTableProbe.cs}}

\subsection{Drift check (ADR 0005)}\label{sec:standardy-drift}

Repozytorium konsumenta przyjmuje wersję standardów, a~potem zostaje w~tyle, gdy standardy
się zmieniają. Repozytorium \repopath{context-pin} rozwiązywało to usługą: API manifestu,
plik lock synchronizowany z~nim i~GitHub Action porównująca jedno z~drugim. C4-c zachowuje
ideę -- wersję przypiętą w~każdym repozytorium i~CI, które zawodzi, gdy pin jest w~tyle --
bez usługi: pakiet NuGet, tag git i~jeden krok CI.

\subsubsection{Decyzje ADR 0005}

\begin{enumerate}
  \item \textbf{Pin jest czytany tam, gdzie konsument już go deklaruje.} Konsument ma z~tego
        repozytorium dwa artefakty i~oba niosą wersję: referencję pakietu
        \repopath{SecondKey.Standards} (konfiguracja bramki -- \repopath{PackageVersion},
        \repopath{GlobalPackageReference} albo \repopath{PackageReference} w~dowolnym pliku
        MSBuild, albo \repopath{packages.config}) oraz \repopath{metadata.version}
        zainstalowanego skilla (instrukcje agenta). Nie ma trzeciego pliku do utrzymywania
        w~zgodzie. Przykładowa lokalizacja z~ticketu, \path{.secondkey/standards.lock}, została
        odrzucona z~konkretnego powodu: \path{.secondkey/} jest ignorowany przez git
        w~repozytoriach estate (to katalog wyników przebiegów Second Key), także w~benchu
        nopCommerce, więc pin w~tym miejscu nigdy nie zostałby zacommitowany.
        \repopath{-{}-pinned-version} obsługuje konsumenta, który przypina inaczej.
  \item \textbf{Niezgodne piny zawodzą, nawet gdy jeden z~nich jest aktualny.} Jeśli skill
        mówi 0.2.0, a~pakiet 0.1.0, agent migracyjny i~bramka pracują na różnych
        standardach -- to dokładnie ten dryf, któremu komponent ma zapobiegać.
  \item \textbf{Najnowsze wydanie to najwyższy tag \repopath{v<MAJOR.MINOR.PATCH>}}
        repozytorium standardów, listowany przez \repopath{git ls-remote}. Tagi pre-release
        i~tagi niebędące wersjami są pomijane. Bramka wersji generatora i~strażnik wydań w~CI
        (ADR 0003) czynią tag wiarygodnym: treść otagowanej wersji nigdy się nie zmienia.
        \repopath{-{}-latest-version} obsługuje konfigurację offline lub z~mirrorem;
        \repopath{-{}-source} przyjmuje dowolny URL git albo ścieżkę lokalną, i~tak też
        uruchamiają się offline testy i~fixture CI.
  \item \textbf{„Co się zmieniło” jest drukowane w~dwóch połowach.} Mechaniczna porównuje
        \path{generated/manifest.json} obu wydań: reguły dodane, usunięte i~zmienione -- ze
        zmianą wagi wypisaną wprost, bo to ona decyduje, czy aktualizacja kosztuje pracę.
        Ludzka drukuje wpisy \path{CHANGELOG.md} między pinem a~najnowszym wydaniem. Obie są
        czytane z~tagów płytkim \repopath{git fetch}, więc są dokładnie tym, co te wydania
        dostarczyły.
  \item \textbf{Jedna implementacja, dwa punkty wejścia.} Logika to polecenie
        \repopath{drift-check} narzędzia. Composite action
        (\path{actions/drift-check/action.yml}) instaluje .NET~10 SDK tylko wtedy, gdy runner
        go nie ma, i~uruchamia polecenie z~własnego checkoutu (więc SDK wybiera
        \repopath{global.json} tego repozytorium, a~nie konsumenta). Wejścia docierają do
        skryptu przez zmienne środowiskowe, nigdy przez interpolację wyrażeń w~treść skryptu.
        W~GitHub Actions wynik jest też adnotacją błędu, podsumowaniem kroku i~wyjściami kroku
        (\repopath{status}, \repopath{pinned-version}, \repopath{latest-version}).
  \item \textbf{Kody wyjścia oddzielają znalezisko od niemożności sprawdzenia.} W~tyle,
        przed wydaniami, niezgodność albo brak pinu -- kod 1: sprawdzenie się wykonało
        i~znalazło problem. Nieczytelny pin, źródło bez wydania albo nieosiągalne źródło --
        kod 2. Oba kończą krok błędem: sprawdzenie, które nie mogło się wykonać, nie może
        wyglądać jak sprawdzenie, które przeszło.
  \item \textbf{Fixture, które nie mogą się zestarzeć.} CI sprawdza zatwierdzone fixture
        z~pinami \repopath{0.0.1}/\repopath{0.0.2} -- poniżej każdej prawdziwej wersji, więc
        zawsze w~tyle -- oraz \emph{aktualnego} konsumenta złożonego z~commita, który jest
        testowany, dokładnie tak, jak konsument robi aktualizację. Źródło wydań jest budowane
        z~tego samego commita i~tagowane wersją pakietu. Nic w~fixture nie musi się zmieniać,
        gdy zmienia się wersja.
\end{enumerate}

\subsubsection{Odczyt pinu}

Drift check przegląda rekurencyjnie repozytorium konsumenta, pomijając katalogi
\path{.git}, \path{.vs}, \path{bin}, \path{obj}, \path{node_modules}, \path{packages},
\path{artifacts} i~\path{TestResults} (bez rozróżniania wielkości liter):

\begin{itemize}
  \item \textbf{Pliki MSBuild} (\repopath{.csproj}, \repopath{.vbproj}, \repopath{.fsproj},
        \repopath{.props}, \repopath{.targets}): elementy \repopath{PackageReference},
        \repopath{PackageVersion} i~\repopath{GlobalPackageReference}, których
        \repopath{Include} albo \repopath{Update} jest identyfikatorem pakietu (bez
        rozróżniania wielkości liter). Wersja pochodzi z~atrybutu \repopath{Version},
        atrybutu \repopath{VersionOverride} albo z~elementu podrzędnego \repopath{Version}
        lub \repopath{VersionOverride}. \repopath{PackageReference} bez wersji (central
        package management) jest pomijany -- \repopath{PackageVersion} z~wersją znajduje się
        osobno. Pliki, które nie są poprawnym XML, są pomijane; przetwarzanie DTD jest
        zabronione.
  \item \textbf{\repopath{packages.config}}: element \repopath{package} z~\repopath{id}
        pakietu i~atrybutem \repopath{version}.
  \item \textbf{Skill}: \texttt{<katalog skilli>/<skill-name>/SKILL.md} w~jednym z~czterech
        katalogów skilli, sprawdzanych w~tej kolejności: \path{.github/skills},
        \path{.github/upgrades/skills}, \path{.claude/skills} i~\path{.agents/skills}; wersja
        to \repopath{metadata.version} z~front matter. Skill bez czytelnej wersji jest
        problemem z~komunikatem, żeby skopiować skill ponownie z~wydania zamiast go edytować.
\end{itemize}

Tekst wersji \repopath{[1.2.3]} (dokładna wersja w~notacji NuGet) jest przyjmowany jako pin;
każdy inny zakres albo wersja pływająca -- nie. Wersja ustawiana właściwością MSBuild
(\repopath{\$(...)}) jest problemem z~podpowiedzią, żeby użyć \repopath{-{}-pinned-version};
check nie wylicza MSBuild. Każdy znaleziony pin jest raportowany, żeby niezgodności były
widoczne.

Pin liczy się osobno w~każdym miejscu: referencja w~każdym pliku MSBuild
i~\repopath{packages.config} oraz kopia skilla w~każdym z~czterech katalogów. Status
\repopath{disagree} daje każda para pinów o~różnych wersjach, nie tylko skill i~pakiet --
także dwie kopie skilla, np. po skopiowaniu katalogu również do
\path{.github/upgrades/skills/}, jak radzi tabela rozwiązywania problemów
(sekcja~\ref{sec:standardy-modernize}); przy aktualizacji trzeba więc zaktualizować każdą
kopię. Bez \repopath{-{}-pinned-version} nieczytelny pin w~jednym pliku daje status
\repopath{unreadable} (kod 2), nawet gdy pozostałe piny są poprawne.

\subsubsection{Wyniki i~kody wyjścia}

\begin{xltabular}{\linewidth}{@{}>{\raggedright\arraybackslash}p{3.3cm} >{\raggedright\arraybackslash}p{0.9cm} >{\raggedright\arraybackslash}X@{}}
\toprule
Status & Kod & Znaczenie \\
\midrule
\endfirsthead
\toprule
Status & Kod & Znaczenie \\
\midrule
\endhead
\path{current} & 0 & Pin jest najnowszym wydaniem: „\texttt{Up to date: this repository pins standards <v>, the latest release.}” \\
\path{behind} & 1 & Istnieje nowsze wydanie; wypisuje piny, reguły zmienione między wersjami, wpisy changeloga i~instrukcję aktualizacji \\
\path{ahead} & 1 & Pin jest wyższy niż każde wydanie: przypina coś, co nigdy nie zostało wydane \\
\path{disagree} & 1 & Piny mają różne wersje: skill agenta i~pakiet bramki albo dowolne dwa inne piny (dwa pliki projektu, dwie kopie skilla) \\
\path{unpinned} & 1 & Nic w~repozytorium nie przypina standardów (oczekiwana referencja pakietu albo skill w~\path{.github/skills/<skill>/SKILL.md}) \\
\path{unreadable} & 2 & Pin istnieje, ale nie da się go odczytać jako dokładnej wersji \\
\path{no-release} & 2 & Źródło nie ma wydania (tagu \repopath{v<MAJOR.MINOR.PATCH>}), z~którym można porównać \\
\path{source-unavailable} & 2 & Źródła nie dało się odczytać \\
\bottomrule
\end{xltabular}

Dla statusu \repopath{behind} część mechaniczna ma postać listy
\texttt{Rules changed between <pin> and <latest>:} z~wierszami \texttt{added},
\texttt{removed} albo \texttt{changed}, identyfikatorem, szczegółem i~tytułem; szczegół
\texttt{changed} to zmiana wagi (\texttt{severity warning → error}), zmiana diagnostyk albo
\texttt{text}, gdy zmienił się tylko skrót reguły. Gdy żadna reguła się nie zmieniła, lista
mówi, że różnica jest w~pakowaniu. Część ludzka to wpisy changeloga po pinie (wyłącznie) do
najnowszego wydania (włącznie). Jeśli przypięta wersja nie jest wydaniem, wynik to dodatkowo
zaznacza. Na końcu jest instrukcja: zmienić referencję pakietu na najnowszą wersję i~ponownie
skopiować skill z~tagu \repopath{v<latest>} (z~\path{generated/skills/<skill>/} do
\path{.github/skills/<skill>/}).

Podsumowania pozostałych statusów (na stderr z~prefiksem \texttt{error:}, w~GitHub Actions
także jako adnotacja): \repopath{behind} -- „\texttt{Behind: this repository pins standards
<pin>; the latest release is <latest>.}”; \repopath{ahead} -- „\texttt{This repository pins
standards <pin>, but the latest release is <latest>: <pin> was never released.}” ze
wskazówką „\texttt{Pin a released version (the latest is <latest>).}”; \repopath{disagree}
-- „\texttt{The pins disagree (<a> and <b>): the migration agent and the gate are working
from different standards.}”; \repopath{unpinned} -- „\texttt{This repository does not pin
the Second Key standards.}”; \repopath{unreadable} -- „\texttt{A standards pin could not be
read.}” z~listą problemów; \repopath{no-release} -- „\texttt{No release (a tag
v<MAJOR.MINOR.PATCH>) was found at <źródło>; there is nothing to compare the pin with.}”;
\repopath{source-unavailable} -- „\texttt{The standards releases could not be read:
<powód>}”. Dla wierszy \texttt{added} i~\texttt{removed} szczegółem jest waga reguły. Gdy
plików wydań nie da się pobrać, zamiast list wynik mówi „\texttt{What changed could not be
read: <powód>}”.

W~GitHub Actions wynik inny niż \repopath{current} jest adnotacją \texttt{::error::}, do
\repopath{\$GITHUB\_STEP\_SUMMARY} trafia sekcja „\texttt{\#\#\# Second Key standards drift
check: <status>}” ze źródłem wydań i~szczegółami, a~do \repopath{\$GITHUB\_OUTPUT} --
\repopath{status}, \repopath{pinned-version} i~\repopath{latest-version}.

\subsubsection{Odczyt wydań przez git}

Lista wydań pochodzi z~\texttt{git ls-remote -{}-tags -{}-refs <źródło>}: brane są referencje
\path{refs/tags/v<wersja>}, wersje pre-release są pomijane, a~lista jest sortowana malejąco.
Pliki wydania (\path{generated/manifest.json}, \path{CHANGELOG.md}) są czytane w~tymczasowym
repozytorium \path{secondkey-drift-*}: \texttt{git init -{}-quiet}, potem \texttt{git fetch
-{}-quiet -{}-depth 1 <źródło> refs/tags/v<x>:refs/tags/v<x>} i~\texttt{git show v<x>:<ścieżka>};
brak tagu („couldn't find remote ref”) oznacza brak pliku. Każde polecenie git ma limit
60~sekund, a~\repopath{GIT\_TERMINAL\_PROMPT=0} sprawia, że źródło wymagające poświadczeń jest
raportowane jako nieosiągalne zamiast czekać na prompt. Brak git w~systemie oznacza status
\repopath{source-unavailable}. \repopath{-{}-source} będący URL-em ze schematem
\repopath{https}, \repopath{http}, \repopath{ssh}, \repopath{git} lub \repopath{file} albo
adresem w~stylu scp (\repopath{user@host:ścieżka}) jest przekazywany do git bez zmian; każda
inna wartość jest ścieżką lokalną względem katalogu bieżącego.

\subsubsection{Użycie}

Krok CI w~repozytorium konsumenta (README):

\begin{lstlisting}
# .github/workflows/standards.yml in the consuming repository
name: Standards drift
on: [pull_request, push]
jobs:
  drift:
    runs-on: ubuntu-latest
    steps:
      - uses: actions/checkout@v7
      - uses: konradcinkusz/letsgolegacy.secondkey-standards/actions/drift-check@main
\end{lstlisting}

To samo sprawdzenie bez akcji:

\begin{lstlisting}
dotnet run --project src/SecondKey.Standards.Generator -- drift-check --repo /path/to/consumer
\end{lstlisting}

\begin{xltabular}{\linewidth}{@{}>{\raggedright\arraybackslash}p{3.6cm} >{\raggedright\arraybackslash}X@{}}
\toprule
Wejście akcji / opcja & Znaczenie i~wartość domyślna \\
\midrule
\endfirsthead
\toprule
Wejście akcji / opcja & Znaczenie i~wartość domyślna \\
\midrule
\endhead
\path{working-directory} / \path{--repo} & Sprawdzane repozytorium konsumenta. Akcja: względem workspace (albo bezwzględna), domyślnie \repopath{.}; polecenie: domyślnie katalog bieżący \\
\path{source} / \path{--source} & Skąd czytane są wydania: URL git repozytorium standardów albo ścieżka lokalna do klonu lub mirrora; domyślnie \path{https://github.com/konradcinkusz/letsgolegacy.secondkey-standards}. W~akcji ścieżka względna jest względna wobec workspace \\
\path{package-id} / \path{--package-id} & Pakiet NuGet z~konfiguracją bramki; domyślnie \repopath{SecondKey.Standards} \\
\path{skill-name} / \path{--skill-name} & Nazwa katalogu zainstalowanego skilla; domyślnie \path{applying-dotnet-migration-standards} \\
\path{pinned-version} / \path{--pinned-version} & Sprawdza tę wersję zamiast czytać referencję pakietu i~skill; domyślnie puste \\
\path{latest-version} / \path{--latest-version} & Traktuje tę wersję jako najnowsze wydanie zamiast listować tagi (konfiguracja offline lub z~mirrorem); domyślnie puste \\
\bottomrule
\end{xltabular}

Wersje w~\repopath{-{}-pinned-version} i~\repopath{-{}-latest-version}, tak jak piny czytane
z~plików, mają postać \repopath{MAJOR.MINOR.PATCH} bez zer wiodących i~bez prefiksu
\repopath{v}, z~opcjonalnym sufiksem pre-release (\repopath{-rc.1}) i~metadanymi builda
(\repopath{+build.7}); opcja w~innej postaci kończy się kodem 2 („\texttt{-{}-pinned-version
and -{}-latest-version take a MAJOR.MINOR.PATCH version}”). Metadane builda nie są częścią
wersji (\repopath{1.2.3+build.7} równa się \repopath{1.2.3}), a~pre-release sortuje się
przed swoim wydaniem. Wyjścia akcji: \repopath{status} (jedna z~wartości
z~tabeli wyników), \repopath{pinned-version}, \repopath{latest-version}.

Kroki composite action: (1) sprawdzenie, czy na runnerze jest .NET~10 SDK (\texttt{dotnet
-{}-list-sdks}); (2) \repopath{actions/setup-dotnet@v6} z~\repopath{10.0.x} tylko wtedy, gdy
go nie ma, żeby nie ruszać konfiguracji SDK konsumenta; (3) uruchomienie
\texttt{dotnet run -{}-project src/SecondKey.Standards.Generator -{}-configuration Release -{}-
drift-check …} w~katalogu \texttt{\$\{\{ github.action\_path \}\}/../..}, z~wejściami
przekazanymi jako zmienne \repopath{INPUT\_*} oraz \repopath{DOTNET\_NOLOGO},
\path{DOTNET_CLI_TELEMETRY_OPTOUT} i~\path{DOTNET_SKIP_FIRST_TIME_EXPERIENCE}.

\subsubsection{Konsekwencje i~odrzucone alternatywy}

Dopóki nie ma pierwszego otagowanego wydania, sprawdzenie względem prawdziwego repozytorium
zwraca „no release” (kod 2) -- według ADR poprawnie, bo nie ma jeszcze czego przypinać.
Konsument, który odwołuje się do pakietu przez właściwość MSBuild
(\texttt{Version="\$(StandardsVersion)"}), musi podać \repopath{-{}-pinned-version}. Akcja
buduje narzędzie przy każdym uruchomieniu (dziesiątki sekund); konsument, który uruchamia ją
często, może zainstalować spakowane narzędzie
(\texttt{dotnet tool install SecondKey.Standards.Generator}) z~własnego feedu i~uruchamiać
\texttt{secondkey-standards drift-check}.

Odrzucone warianty:
\begin{itemize}
  \item \textbf{Usługa i~plik lock z~\repopath{context-pin}} -- serwer do utrzymania
        i~zabezpieczenia oraz lock do synchronizowania: dwa nowe ruchome elementy dla pytania,
        na które odpowiadają już tagi git.
  \item \textbf{Porównanie z~najnowszą wersją w~feedzie NuGet} -- pakiet nie jest publikowany
        w~nuget.org, a~wewnętrzny feed konsumenta może być opóźniony; tagi są jedynym
        miejscem, w~którym istnieje każde wydanie.
  \item \textbf{Implementacja w~bashu wewnątrz akcji} -- dwie implementacje tego samego
        sprawdzenia rozjechałyby się; akcja i~polecenie dzielą jedną.
\end{itemize}

\zrodla{\path{letsgolegacy.secondkey-standards/docs/adr/0005-drift-check.md},
\path{letsgolegacy.secondkey-standards/README.md},
\path{letsgolegacy.secondkey-standards/actions/drift-check/action.yml},
\path{letsgolegacy.secondkey-standards/src/SecondKey.Standards.Generator/Drift/DriftChecker.cs},
\path{letsgolegacy.secondkey-standards/src/SecondKey.Standards.Generator/Drift/PinReader.cs},
\path{letsgolegacy.secondkey-standards/src/SecondKey.Standards.Generator/Drift/Pin.cs},
\path{letsgolegacy.secondkey-standards/src/SecondKey.Standards.Generator/Drift/GitReleaseSource.cs},
\path{letsgolegacy.secondkey-standards/src/SecondKey.Standards.Generator/Drift/ManifestDiff.cs},
\path{letsgolegacy.secondkey-standards/src/SecondKey.Standards.Generator/Drift/DriftReport.cs},
\path{letsgolegacy.secondkey-standards/src/SecondKey.Standards.Generator/Model/SemanticVersion.cs},
\path{letsgolegacy.secondkey-standards/src/SecondKey.Standards.Generator/Cli/CommandLineApp.cs},
\path{letsgolegacy.secondkey-standards/docs/USING-WITH-MODERNIZE-DOTNET.md}}

\subsection{Użycie z~modernize-dotnet}\label{sec:standardy-modernize}

Ta sekcja przenosi \path{docs/USING-WITH-MODERNIZE-DOTNET.md}: jak osoba z~dostępem do GitHub
Copilot instaluje standardy Second Key w~repozytorium docelowym jako własny skill
i~uruchamia agenta modernizacyjnego Microsoft dla .NET tak, żeby ten skill podjął. Jest to
procedura dla ticketu P6 benchu (przebieg modernize-dotnet na nopCommerce ze skillem
standardów), napisana tak, by dało się ją zastosować do każdego repozytorium .NET Framework;
nikt jej dotąd nie przeprowadził, bo P6 jest planowane (WORKPLAN benchu).

\textbf{Nazewnictwo.} Microsoft dokumentuje obecnie agenta jako \textbf{GitHub Copilot
upgrade} (część .NET nazywała się wcześniej „GitHub Copilot app modernization for .NET”,
a~backlog nazywa wtyczkę \emph{modernize-dotnet}). W~Visual Studio jest nadal instalowany jako
komponent \textbf{GitHub Copilot app modernization} i~wywoływany jako \repopath{@Modernize};
w~pozostałych środowiskach to agent \textbf{Upgrade} (\repopath{@upgrade}). To ten sam agent
i~ten sam mechanizm skilli.

\subsubsection{Krok 1: co czyta agent}

Skill jest generowany w~tym repozytorium w~\path{generated/skills/applying-dotnet-migration-standards/}:
\repopath{SKILL.md} (front matter i~instrukcje) oraz po jednym pliku
\path{references/<rule-id>.md} na regułę. To format Agent Skills -- katalog nazwany jak skill,
zawierający \repopath{SKILL.md} -- z~dwoma polami metadanych, które dodaje GitHub Copilot
upgrade:

\begin{lstlisting}
name: applying-dotnet-migration-standards
description: "Applies the Second Key migration standards while a .NET Framework application is upgraded ..."
metadata:
  discovery: "preload"      # GitHub Copilot upgrade: always available, not only on a description match
  traits: ".NET|CSharp"     # GitHub Copilot upgrade: the technologies it applies to
  version: "0.1.0"          # the standards version this copy came from (the drift check reads it)
  source: "https://github.com/konradcinkusz/letsgolegacy.secondkey-standards"
\end{lstlisting}

GitHub Copilot upgrade czyta własne skille z~tych miejsc, od najwyższego priorytetu:

\begin{tabularx}{\linewidth}{@{}>{\raggedright\arraybackslash}p{6.2cm} >{\raggedright\arraybackslash}X@{}}
\toprule
Lokalizacja & Zasięg \\
\midrule
\texttt{\%UserProfile\%/.copilot/skills/} & Profil osoby, każde repozytorium \\
\path{.github/upgrades/skills/} & Repozytorium, tylko dla upgrade \\
\path{.github/skills/} & Repozytorium, współdzielone z~zespołem -- \textbf{używane w~tej procedurze} \\
\bottomrule
\end{tabularx}

Wszystko, czego skill potrzebuje, jest w~jego katalogu, więc skopiowanie katalogu jest całą
instalacją.

\subsubsection{Krok 2: wymagania wstępne}

\begin{itemize}
  \item Subskrypcja GitHub Copilot (płatna albo darmowa) na koncie GitHub, którym osoba się
        loguje. Agent chmurowy (ostatnia opcja w~kroku 4) wymaga Copilot Business albo
        Enterprise z~włączonymi coding agents.
  \item Repozytorium docelowe sklonowane lokalnie, na gałęzi przeznaczonej na migrację. Dla P6
        to źródła nopCommerce 3.x przypięte przez bench (bench je pobiera; najpierw należy je
        zacommitować do repozytorium roboczego, żeby zmiany agenta dało się przejrzeć jako
        diff).
  \item Git i~.NET~10 SDK (rozszerzenie Visual Studio Code instaluje SDK, jeśli go brakuje).
\end{itemize}

\subsubsection{Krok 3: instalacja skilla w~repozytorium docelowym}

Wybiera się wersję standardów: tag wydania \repopath{v<wersja>}, gdy już istnieje, albo
\repopath{main} przed pierwszym wydaniem (wtedy trzeba zapisać commit, z~którego pochodzi
kopia).

Bash:

\begin{lstlisting}
STANDARDS_REF=main            # or a release tag, e.g. v0.1.0
SKILL=applying-dotnet-migration-standards

git clone --depth 1 --branch "$STANDARDS_REF" \
  https://github.com/konradcinkusz/letsgolegacy.secondkey-standards.git /tmp/secondkey-standards

cd /path/to/target-repository
mkdir -p .github/skills
rm -rf ".github/skills/$SKILL"
cp -R "/tmp/secondkey-standards/generated/skills/$SKILL" .github/skills/
git add ".github/skills/$SKILL"
git commit -m "Add Second Key migration standards skill ($STANDARDS_REF)"
\end{lstlisting}

PowerShell:

\begin{lstlisting}
$StandardsRef = 'main'        # or a release tag, e.g. 'v0.1.0'
$Skill = 'applying-dotnet-migration-standards'
$Source = Join-Path $env:TEMP 'secondkey-standards'

git clone --depth 1 --branch $StandardsRef `
  https://github.com/konradcinkusz/letsgolegacy.secondkey-standards.git $Source

Set-Location C:\path\to\target-repository
New-Item -ItemType Directory -Force .github\skills | Out-Null
Remove-Item -Recurse -Force ".github\skills\$Skill" -ErrorAction SilentlyContinue
Copy-Item -Recurse "$Source\generated\skills\$Skill" .github\skills\
git add ".github/skills/$Skill"
git commit -m "Add Second Key migration standards skill ($StandardsRef)"
\end{lstlisting}

Sprawdzenie wyniku: istnieje
\path{.github/skills/applying-dotnet-migration-standards/SKILL.md}, a~\repopath{metadata.version}
w~jego front matter to wersja, którą zamierzano wziąć. Kopii się nie edytuje; nowszą wersję
bierze się z~repozytorium standardów.

\subsubsection{Krok 4: instalacja agenta}

Jedna z~poniższych opcji. Kroki pochodzą od Microsoft; źródła wymienia koniec tej sekcji.

\begin{description}
  \item[Visual Studio (Windows).] Visual Studio 2026 albo Visual Studio 2022 17.14.17 lub
        nowsze. W~Visual Studio Installer, w~workloadzie \textbf{.NET desktop development},
        włącza się opcjonalne komponenty \textbf{GitHub Copilot} i~\textbf{GitHub Copilot app
        modernization}. Logowanie do Visual Studio kontem GitHub z~Copilot. Sprawdzenie: po
        kliknięciu projektu prawym przyciskiem w~Solution Explorer widać \textbf{Modernize}.
  \item[Visual Studio Code.] Instalacja rozszerzenia \textbf{GitHub Copilot}, a~potem
        rozszerzenia \textbf{GitHub Copilot upgrade} z~widoku Extensions. Sprawdzenie:
        w~Copilot Chat odpowiada \repopath{@upgrade} albo lista agentów zawiera
        \textbf{Upgrade}.
  \item[GitHub Copilot CLI.] W~sesji Copilot CLI:
\begin{lstlisting}
/plugin marketplace add microsoft/upgrade-agent-plugins
/plugin install upgrade-agent@upgrade-agent-plugins
\end{lstlisting}
        Sprawdzenie: \repopath{/agent} wymienia \repopath{upgrade-agent}.
  \item[Aplikacja GitHub Copilot.] Dodanie marketplace \path{microsoft/upgrade-agent-plugins}
        w~\textbf{Settings → Plugins} i~instalacja wtyczki \textbf{upgrade-agent}.
        Sprawdzenie: lista agentów zawiera \textbf{Upgrade}.
  \item[GitHub.com (agent chmurowy).] Skopiowanie \path{cloud-agent/upgrade.agent.md}
        z~\path{microsoft/upgrade-agent-plugins} do \path{.github/agents/} repozytorium
        docelowego oraz \path{cloud-agent/windows/copilot-setup-steps.yml} (dla .NET
        Framework) do \path{.github/workflows/copilot-setup-steps.yml}. Konfiguracja Windows
        wymaga też wyłączenia firewalla coding agenta (\textbf{Settings → Copilot → Coding
        agent}), co usuwa jego ograniczenia sieciowe -- trzeba postępować zgodnie z~README
        wtyczki i~rozważyć to przed użyciem. Skille z~\path{.github/skills/} są dostępne także
        dla agenta chmurowego.
\end{description}

\subsubsection{Krok 5: uruchomienie upgrade'u tak, żeby użył skilla}

\begin{enumerate}
  \item Otworzyć repozytorium docelowe (w~Visual Studio -- rozwiązanie).
  \item Uruchomić agenta:
        \begin{itemize}
          \item Visual Studio: prawy przycisk na rozwiązaniu → \textbf{Modernize} albo
                \repopath{@Modernize} w~Copilot Chat.
          \item Visual Studio Code: \repopath{@upgrade} w~Copilot Chat albo \textbf{Upgrade}
                na liście agentów.
          \item Copilot CLI: \repopath{/agent} i~wybór \textbf{Upgrade} (albo prompt zaczynający
                się od \repopath{@upgrade}). Najpierw \repopath{/skills list}; jeśli skilla nie
                ma -- \repopath{/skills reload} (krok 6).
          \item Aplikacja Copilot: \textbf{Upgrade} na liście agentów.
        \end{itemize}
  \item Jako pierwszy prompt wysłać dokładnie:
\begin{lstlisting}
Upgrade this solution to .NET 10. Apply the applying-dotnet-migration-standards skill from
.github/skills to every task: preserve behaviour, and record every behaviour flag in
docs/migration/behaviour-flags.md instead of fixing it.
\end{lstlisting}
  \item Uczynić to trwałym dla całego upgrade'u, żeby widziało to każde zadanie i~każda
        późniejsza sesja:
\begin{lstlisting}
From now on, for all tasks in this upgrade, follow the applying-dotnet-migration-standards skill.
\end{lstlisting}
        Agent zapisuje instrukcje sformułowane w~ten sposób do
        \path{.github/upgrades/<scenarioId>/scenario-instructions.md}, który ładuje do każdej
        decyzji.
  \item Przejść przez etapy agenta jak zwykle -- \textbf{assessment}
        (\repopath{assessment.md}), \textbf{planning} (\repopath{plan.md}),
        \textbf{execution} (\repopath{tasks.md}) -- przeglądając każdy dokument, zanim agent
        dostanie polecenie kontynuacji. Przy ocenie (assessment) sprawdzić, czy jest w~niej
        baseline legacy ze skilla (źródło kultury, klucze konfiguracji, lazy loading EF6,
        endpointy JSON, nieprzechodzące testy); jeśli nie ma, napisać:
\begin{lstlisting}
Record the legacy baseline the applying-dotnet-migration-standards skill asks for before planning
\end{lstlisting}
  \item Gdy agent skończy, akceptacja P6 jest kryterium Microsoft: kandydat się buduje,
        a~istniejące testy przechodzą. Flagi zachowania są tym, czego agent \emph{nie}
        rozstrzygnął; przekazuje się je razem z~kandydatem do P7 (replay i~porównanie) i~P8
        (bramka).
\end{enumerate}

\subsubsection{Krok 6: potwierdzenie, że agent podjął skill}

To jest dowód, który zapisuje P6. Musi być spełnione co najmniej jedno z~poniższych,
a~powinny -- pierwsze dwa:

\begin{xltabular}{\linewidth}{@{}>{\raggedright\arraybackslash}p{3.1cm} >{\raggedright\arraybackslash}X >{\raggedright\arraybackslash}p{3.6cm}@{}}
\toprule
Kontrola & Jak & Oczekiwany wynik \\
\midrule
\endfirsthead
\toprule
Kontrola & Jak & Oczekiwany wynik \\
\midrule
\endhead
Agent zna skill & Zapytać: \texttt{Which custom skills from .github/skills are you applying?} (Copilot CLI: \texttt{/skills info applying-dotnet-migration-standards}) & Wymienia \path{applying-dotnet-migration-standards} \\
Instrukcja została utrwalona & Otworzyć \path{.github/upgrades/<scenarioId>/scenario-instructions.md} & Zawiera instrukcję z~kroku 5.4 \\
Baseline został zapisany & Otworzyć \path{docs/migration/behaviour-flags.md} & Jest nagłówek i~źródło kultury legacy \\
Reguły zostały zastosowane & Pierwsze polecenie poniżej; identyfikatory reguł (\repopath{SK-MIG-}) w~ocenie, planie albo notatkach zadań & Flagi w~miejscach, które wymienia rejestr \\
Brak cichego przełączenia & Drugie polecenie poniżej & Nic nowego albo każde trafienie oflagowane \\
\bottomrule
\end{xltabular}

\begin{lstlisting}
git grep -n "SECONDKEY-FLAG"
git grep -n -E "InvariantGlobalization|UseNls|AppLocalIcu|EnsureCreated"
\end{lstlisting}

Do notatek P6 zapisuje się: narzędzie i~jego wersję, wersję standardów
z~\repopath{metadata.version} skilla, czy skill został podjęty samodzielnie, czy dopiero po
jawnym prompcie, oraz liczbę flag na regułę.

\subsubsection{Krok 7: konfiguracja bramki w~zmigrowanym rozwiązaniu}

Nie jest potrzebna, żeby agent użył skilla. Dotyczy buildu zmigrowanego rozwiązania, w~którym
wagi standardów mają obowiązywać w~kompilatorze. Pakiet NuGet \repopath{SecondKey.Standards}
ustawia wagę każdej diagnostyki, która sprawdza standard, ale tylko tam: P8 benchu ma uruchomić
na pull requeście kandydata skan Portcullis z~wynikiem SARIF, a~ten nie czyta wag z~pakietu
(sekcja~\ref{sec:standardy-wagi-kanaly}). Pakietu nie ma w~nuget.org: pobiera się artefakt
\repopath{secondkey-standards-packages} z~uruchomienia CI tego repozytorium dla wziętej wersji
albo pakuje samodzielnie przez \path{./scripts/verify-package.sh} w~klonie repozytorium.
Następnie w~zmigrowanym rozwiązaniu:

\begin{lstlisting}
mkdir -p .packages && cp /path/to/SecondKey.Standards.0.1.0.nupkg .packages/
\end{lstlisting}

\begin{lstlisting}
<!-- nuget.config at the repository root: a feed relative to this file -->
<configuration>
  <packageSources>
    <add key="secondkey-local" value=".packages" />
  </packageSources>
</configuration>
\end{lstlisting}

\begin{lstlisting}
<!-- Directory.Packages.props (central package management) -->
<GlobalPackageReference Include="SecondKey.Standards" Version="0.1.0" />

<!-- or, without central package management, in Directory.Build.props -->
<ItemGroup>
  <PackageReference Include="SecondKey.Standards" Version="0.1.0" PrivateAssets="all" />
</ItemGroup>
\end{lstlisting}

Wersja musi być ta sama co wersja skilla: agent i~bramka muszą dostać te same standardy.

Pakiet ustawia tylko wagi: diagnostyki \repopath{CA} zgłasza .NET SDK, a~diagnostyki
\repopath{PORTCULLIS\_} pojawią się w~buildzie dopiero z~pakietem analyzerów Portcullis
(sekcja~\ref{sec:portcullis-dystrybucja}). Sprawdzenie, że konfiguracja dotarła do buildu
-- tak, jak robi to \path{scripts/verify-package.sh}: wywołanie zależne od kultury, np.
\texttt{string.Compare(left, right)}, daje w~\repopath{dotnet build}
\texttt{warning CA1310}; z~\path{-p:SecondKeyStandardsGlobalConfig=false} (albo bez
pakietu) CA1310 nie jest zgłaszany.

\subsubsection{Krok 8: utrzymanie aktualnego pinu}

Kopia skilla i~referencja pakietu są pinem repozytorium docelowego: obie niosą wersję
standardów. Jeden krok w~jego CI kończy się błędem, gdy istnieje nowsze wydanie (i~wypisuje
brakujące reguły oraz wpisy changeloga) albo gdy skill i~pakiet się nie zgadzają:

\begin{lstlisting}
- uses: konradcinkusz/letsgolegacy.secondkey-standards/actions/drift-check@main
\end{lstlisting}

Aktualizacja: wziąć nowszy tag, ponownie skopiować skill (krok 3) i~ustawić zgodną wersję
pakietu (krok 7). Dopóki pierwsze wydanie standardów nie jest otagowane, krok zgłasza „no
release” i~kończy się błędem -- nie ma jeszcze czego przypinać; dla P6 przed wydaniem krok
się pomija.

\subsubsection{Rozwiązywanie problemów}

\begin{xltabular}{\linewidth}{@{}>{\raggedright\arraybackslash}p{4.6cm} >{\raggedright\arraybackslash}p{3.6cm} >{\raggedright\arraybackslash}X@{}}
\toprule
Objaw & Przyczyna & Co zrobić \\
\midrule
\endfirsthead
\toprule
Objaw & Przyczyna & Co zrobić \\
\midrule
\endhead
Agent nie wspomina o~skillu & Skill nie został wykryty albo nie został dopasowany do żądania & Sprawdzić, że ścieżka to \path{.github/skills/applying-dotnet-migration-standards/SKILL.md}; poprosić o~skill z~nazwy (krok 5.3); zrestartować Visual Studio albo VS Code po dodaniu skilla; w~Copilot CLI uruchomić \repopath{/skills reload} \\
Nadal nie podjęty & Inna wersja klienta czyta inną lokalizację & Skopiować ten sam katalog także do \path{.github/upgrades/skills/} (lokalizacja upgrade, wyższy priorytet), zacommitować i~zrestartować \\
\repopath{/skills info} nie pokazuje opisu albo pokazuje błąd & Front matter został edytowany albo uszkodzony przy kopiowaniu (końce linii, kodowanie) & Skopiować katalog ponownie z~tagu; nie edytować go \\
Agent włącza \path{InvariantGlobalization} albo zmienia oczekiwaną wartość testu & Poszedł za ogólną poprawką zamiast za regułą & Wskazać mu regułę: \texttt{This violates SK-MIG-007 (or SK-MIG-013). Revert it and record a behaviour flag instead.} i~dodać regułę do trwałej instrukcji \\
Build kończy się błędem \repopath{PORTCULLIS\_} albo \repopath{CA} & Konfiguracja bramki jest zainstalowana (krok 7), a~kod narusza regułę o~wadze \repopath{error} & Naprawić zgodnie z~regułą albo stłumić wiersz z~uzasadnieniem \repopath{SECONDKEY-FLAG}, jak opisuje skill \\
\bottomrule
\end{xltabular}

W~wersji 0.1.0 wszystkie zmapowane diagnostyki \repopath{CA} mają wagę \repopath{warning}
(sekcja~\ref{sec:standardy-przeglad}); wagę \repopath{error} mają w~tej wersji tylko
diagnostyki \repopath{PORTCULLIS\_}.

\subsubsection{Co jest zweryfikowane, a~co nie}

Zweryfikowane z~bieżącą dokumentacją Microsoft i~GitHub (odczyt 2026-09-26):

\begin{itemize}
  \item Własne skille, ich lokalizacje i~priorytet oraz pola \repopath{name},
        \repopath{description}, \repopath{metadata.discovery} i~\repopath{metadata.traits}:
        \emph{Customize GitHub Copilot upgrade}
        (\url{https://learn.microsoft.com/dotnet/core/porting/github-copilot-upgrade/customization};
        źródło \repopath{dotnet/docs},
        \path{docs/core/porting/github-copilot-upgrade/customization.md}, datowane 2026-09-21),
        łącznie z~tym, że instrukcje „From now on…” trafiają do
        \repopath{scenario-instructions.md}.
  \item Instalacja i~wywołanie w~każdym środowisku: \emph{Install GitHub Copilot upgrade}
        (ten sam katalog, \repopath{install.md}, 2026-09-21) i~README
        \path{microsoft/upgrade-agent-plugins}.
  \item Format katalogu skilla (\path{<name>/SKILL.md}, ograniczenia nazwy i~opisu):
        specyfikacja Agent Skills; dokument GitHub \emph{About agent skills} (skille projektu
        w~\path{.github/skills}, \path{.claude/skills} albo \path{.agents/skills}); skille
        dostarczane przez Microsoft
        w~\path{microsoft/upgrade-agent-plugins}, które używają tego samego układu. Nazwa
        stosuje regułę autorską Microsoft dla skilli upgrade (zaczyna się od gerundium:
        \repopath{applying-}).
  \item \repopath{/skills list}, \repopath{/skills info} i~\repopath{/skills reload} w~Copilot
        CLI: dokument GitHub \emph{Adding agent skills for GitHub Copilot CLI}.
\end{itemize}

Niezweryfikowane tutaj i~celowo pozostawione P6, które uruchamia człowiek z~GitHub Copilot:

\begin{itemize}
  \item \textbf{Że agent podejmuje ten skill w~prawdziwym przebiegu upgrade'u.} Nic w~tym
        repozytorium nie potrafi uruchomić GitHub Copilot. Kryterium akceptacji R4 -- „agent
        podejmuje go w~tickecie P6 benchu” -- potwierdzają albo obalają kontrole z~kroku 6.
  \item Jak \repopath{metadata.discovery: preload} zachowuje się dla skilla repozytorium
        w~każdym kliencie. Microsoft dokumentuje \repopath{preload} jako „always available”;
        jeśli klient traktuje go inaczej, jawny prompt z~kroku 5 i~tak wprowadza skill.
  \item Strona Microsoft o~dostosowaniu pokazuje przykładową ścieżkę pojedynczego pliku
        (\path{.github/skills/my-skill.md}), podczas gdy dostarczane skille Microsoft
        i~dokumentacja GitHub używają formy katalogu \path{<name>/SKILL.md}, użytej tutaj. Jeśli
        P6 stwierdzi, że forma katalogu nie jest wykrywana, pierwsze do sprawdzenia są wiersze
        tabeli rozwiązywania problemów, a~znalezisko należy zgłosić jako issue w~tym
        repozytorium.
\end{itemize}

\zrodla{\path{letsgolegacy.secondkey-standards/docs/USING-WITH-MODERNIZE-DOTNET.md},
\path{letsgolegacy.secondkey-standards/docs/adr/0004-skill-for-github-copilot-upgrade.md},
\path{letsgolegacy.secondkey-standards/generated/nuget/README.md},
\path{letsgolegacy.secondkey-standards/generated/config/.globalconfig},
\path{letsgolegacy.secondkey-standards/scripts/verify-package.sh},
\path{letsgolegacy.secondkey-nopcommerce-bench/docs/WORKPLAN.md}}

\subsection{Testy}\label{sec:standardy-testy}

Projekt \path{tests/SecondKey.Standards.Generator.Tests} używa xUnit v3 na
\path{Microsoft.Testing.Platform} (projekt testów jest aplikacją \repopath{Exe}, więc da się go też
uruchomić bezpośrednio) i~referencjonuje generator. Testy uruchamiają prawdziwe polecenia
w~procesie, na jednorazowych repozytoriach w~katalogu tymczasowym: \repopath{TestRepository}
tworzy repozytorium standardów (\path{catalog/pack.json}, szablon skilla, changelog, pliki
reguł pisane przez test), \repopath{Consumer} -- repozytorium konsumenta, a~\repopath{StandardsSource}
-- lokalne repozytorium git z~tagami wydań dla drift checka. Część testów sprawdza samo to
repozytorium: jego reguły, zatwierdzone drzewo generowane i~changelog. W~opisywanej wersji
zbiór liczy 166 przypadków testowych (108 metod \repopath{[Fact]} i~58 wierszy
\repopath{[Theory]} w~16 metodach; 17 klas testowych) i~wszystkie przechodzą.

Testy uruchamia się z~klonu repozytorium: \repopath{RepositoryRoot} szuka
\path{catalog/pack.json} nad katalogiem binariów testów i~bez niego kończy się komunikatem
„\texttt{run the tests from a clone of the repository}”. Testy drift checka tworzą
prawdziwe repozytoria git (z~wyłączonym podpisywaniem i~stałą tożsamością) i~uruchamiają na
nich \texttt{git ls-remote} oraz \texttt{git fetch}, więc wymagają git w~\repopath{PATH}; sieci
nie potrzebują.

\begin{xltabular}{\linewidth}{@{}>{\raggedright\arraybackslash}p{5.1cm} >{\raggedright\arraybackslash}X@{}}
\toprule
Klasa testów & Co sprawdza \\
\midrule
\endfirsthead
\toprule
Klasa testów & Co sprawdza \\
\midrule
\endhead
\path{RepositoryStandardsTests} & Kryterium akceptacji C4-a na regułach tego repozytorium (identyfikator, waga, uzasadnienie, oba przykłady); każda z~czterech migracyjnych diagnostyk Portcullis przypisana dokładnie jednemu standardowi; tematy wymagane przez ticket są pokryte (fragmenty tytułów: \emph{System.Web}, \emph{HttpContext.Current}, \emph{never block on a task}, \emph{IOptions<T>}, \emph{StringComparison}, \emph{culture}, \emph{constructor injection}, \emph{IHttpClientFactory}, \emph{EF6}, \emph{Preserve behaviour}); każda reguła \repopath{SK-ARCH-} ma zasadę i~link do jej kotwicy w~konstytucji \\
\path{RepositoryGeneratedTests} & Zatwierdzone \path{generated/} odpowiada źródłom (to samo porównanie co \repopath{-{}-check}); skill spełnia ograniczenia Agent Skills (nazwa równa katalogowi, opis 1--1024 znaków, \repopath{SKILL.md} poniżej 500 wierszy); każda diagnostyka jest w~\repopath{.globalconfig} z~wagą reguły \\
\path{SkillContentTests} & Skill ma każdą z~ośmiu sekcji; protokół flag w~skillu jest tym z~SK-MIG-011; każda reguła i~każda diagnostyka są w~skillu; każdy identyfikator reguły cytowany przez skill i~przez reguły istnieje; katalog skilla jest samowystarczalny (linki względne wskazują pliki w~katalogu); opis kieruje do skilla upgrade .NET Framework (zawiera m.in. „.NET Framework”, „.NET 10”, „ASP.NET Core”, „EF6”) \\
\path{RuleParserTests} & Parsowanie poprawnej reguły i~każdy błąd walidacji: brak lub niezamknięty front matter, nieznany klucz, nieznana lub źle zapisana diagnostyka Portcullis, zasada spoza konstytucji, analyzer spoza SDK, \repopath{appliesTo} niebędące listą, podpowiedź o~aliasie YAML, nazwa pliku, CRLF i~BOM, brakujące lub puste sekcje, przykład bez kodu, pusty blok kodu, nagłówek \repopath{\#}, nieznana i~powtórzona sekcja, brak stwierdzenia \\
\path{MarkdownBodyParserTests} & Nagłówki w~blokach kodu, ogrodzenia z~tyld i~dłuższe, niezamknięty blok, numery wierszy, nagłówki poziomu trzeciego \\
\path{StandardsLoaderTests} & Kolejność po identyfikatorze i~pominięcie README, zduplikowany identyfikator, diagnostyka przypisana dwóm regułom, wszystkie zepsute pliki w~jednym uruchomieniu, pusty katalog standardów, nieznany klucz w~\path{pack.json} \\
\path{EmissionTests} & Wagi w~\repopath{.globalconfig}, zakresy w~\repopath{.editorconfig}, poprawny YAML front matter skilla, plik referencyjny i~wiersz tabeli dla każdej reguły, dostarczenie konfiguracji przez props pakietu, manifest, skrót treści niezależny od wersji, nieznany i~brakujący placeholder, odrzucenie linku względnego \\
\path{TreeSyncTests} & Pliki brakujące, nieaktualne i~osierocone; zapis usuwa osierocone pliki i~puste katalogi; \path{bin/} i~\path{obj/} nie są osierocone; kopia z~CRLF lub BOM jest nieaktualna; emisja tylko pod \path{generated/} \\
\path{VersionGateTests} & Zmiana treści wydanej wersji odrzucona i~nic nie zapisane; zmiana wersji \repopath{Unreleased} dozwolona; podniesienie wersji przyjmuje nową treść; brak wpisu w~changelogu i~wersja niższa niż zatwierdzona odrzucone; podniesienie bez zmiany treści dozwolone \\
\path{ChangelogTests}, \path{PackConfigTests}, \path{SemanticVersionTests} & Format changeloga (także błędne i~źle uporządkowane nagłówki), wybór wpisów między wersjami, wpis dla wersji pakietu; poprawność \path{pack.json} (wersja bez pre-release, nazwa skilla zgodna ze specyfikacją Agent Skills, opis do 1024 znaków bez nawiasów kątowych, \repopath{discovery}); porządek wersji (liczbowy, pre-release przed wydaniem), tylko wersje trzyczęściowe, metadane builda poza wersją \\
\path{CommandLineTests}, \path{GenerateCommandTests} & \repopath{validate} i~\repopath{generate} na tym repozytorium i~na repozytoriach testowych, adnotacje GitHub, wykrywanie katalogu głównego, uruchomienie poza repozytorium standardów, nieznane polecenie, \repopath{-{}-check} przed i~po generowaniu, ręcznie edytowany plik, niepoprawne standardy nic nie generują \\
\path{DriftCheckerTests} & Na prawdziwych źródłach git (kryterium C4-c): aktualny pin przechodzi, nieaktualny zawodzi i~mówi, co się zmieniło; niezgodne piny, brak pinu, pin przed wydaniami, pomijanie tagów pre-release i~niewersyjnych, źródło bez wydań, nieosiągalne źródło, jawny pin, podana najnowsza wersja, nieaktualny pin, który nie był wydaniem \\
\path{PinReaderTests}, \path{DriftCheckCommandTests} & Odczyt pinu z~\repopath{PackageVersion}, \repopath{GlobalPackageReference}, \repopath{PackageReference} (wielkość liter, notacja \repopath{[x.y.z]}), elementu podrzędnego starego projektu, \repopath{packages.config} i~skilla; ignorowanie wyników builda; polecenie \repopath{drift-check} z~adnotacją, podsumowaniem i~wyjściami w~GitHub Actions, ścieżkami względnymi i~wartościami domyślnymi \\
\bottomrule
\end{xltabular}

Testy skilla (\repopath{SkillContentTests}: 7 metod, 14 przypadków) są strukturalne: dowodzą,
że tekst skilla ma wymagane sekcje i~zgadza się z~regułami i~generatorem, a~nie że agent go
stosuje -- żaden test nie uruchamia agenta (sekcja~\ref{sec:standardy-modernize}).

Fixture dla joba \repopath{drift-check} w~CI (\path{tests/fixtures/drift/}) to repozytoria
konsumentów ograniczone do plików, które czyta drift check:

\begin{tabularx}{\linewidth}{@{}>{\raggedright\arraybackslash}p{2.6cm} >{\raggedright\arraybackslash}p{4.8cm} >{\raggedright\arraybackslash}X@{}}
\toprule
Fixture & Piny & Oczekiwany wynik \\
\midrule
\path{stale/} & pakiet i~skill w~\repopath{0.0.1} & błąd: w~tyle za najnowszym wydaniem, z~opisem zmian \\
\path{disagreeing/} & pakiet w~\repopath{0.0.2}, skill w~\repopath{0.0.1} & błąd: agent i~bramka dostają różne standardy \\
\path{unpinned/} & brak (projekt bez referencji pakietu i~bez skilla) & błąd: repozytorium nie przypina standardów \\
\bottomrule
\end{tabularx}

Aktualny konsument nie jest zatwierdzoną fixture: CI składa go z~wygenerowanego skilla tego
commita i~wersji pakietu, dokładnie tak, jak konsument robi aktualizację, więc nie może się
zestarzeć po podniesieniu wersji. Wersje \repopath{0.0.1} i~\repopath{0.0.2} pozostają
poniżej każdej prawdziwej wersji, więc zatwierdzonych fixture też nie trzeba aktualizować.

\zrodla{\path{letsgolegacy.secondkey-standards/tests/SecondKey.Standards.Generator.Tests/},
\path{letsgolegacy.secondkey-standards/tests/SecondKey.Standards.Generator.Tests/SecondKey.Standards.Generator.Tests.csproj},
\path{letsgolegacy.secondkey-standards/tests/SecondKey.Standards.Generator.Tests/Support/RepositoryRoot.cs},
\path{letsgolegacy.secondkey-standards/tests/SecondKey.Standards.Generator.Tests/Support/StandardsSource.cs},
\path{letsgolegacy.secondkey-standards/tests/fixtures/drift/README.md},
\path{letsgolegacy.secondkey-standards/docs/adr/0002-generator-as-a-dotnet-tool.md}}

\subsection{CI, skan sekretów i~dystrybucja}\label{sec:standardy-ci}

\subsubsection{Workflow CI}

Workflow \path{.github/workflows/ci.yml} jest uruchamiany na każdym \repopath{pull\_request} -- bez filtra gałęzi, bo pull requesty w~tym
repozytorium są piętrowe (bazą każdego jest poprzednia gałąź funkcji), a~filtr na
\repopath{main} zostawiłby bez CI wszystkie poza pierwszym -- na \repopath{push} do
\repopath{main}, na tagach \repopath{v*} (strażnik wydań, bo tag jest tym, co przypinają
konsumenci) i~ręcznie. Uprawnienia \repopath{contents: read}; grupa współbieżności
\repopath{ci-<ref>} z~anulowaniem trwającego przebiegu.

\begin{xltabular}{\linewidth}{@{}>{\raggedright\arraybackslash}p{2.9cm} >{\raggedright\arraybackslash}X@{}}
\toprule
Job & Kroki \\
\midrule
\endfirsthead
\toprule
Job & Kroki \\
\midrule
\endhead
\repopath{build-test} & Ubuntu, limit 15 minut. \repopath{actions/checkout@v7}, \repopath{actions/setup-dotnet@v6} (\repopath{10.0.x}), \repopath{dotnet restore}, build i~testy w~konfiguracji Release. Następnie \repopath{validate} (kryterium C4-a jako krok CI; problemy są adnotacjami na diffie) i~\repopath{generate -{}-check} (kryterium C4-b; odrzuca też zmianę treści wydanej wersji). Na końcu autotest sprawdzenia nieaktualności: dopisanie wiersza do \path{standards/SK-MIG-001-replace-system-web.md} bez regeneracji -- \repopath{-{}-check} musi zawieść; przywrócenie; dopisanie \texttt{dotnet\_diagnostic.CA1310.severity = none} do wygenerowanego \repopath{.globalconfig} -- \repopath{-{}-check} musi zawieść; przywrócenie; \repopath{-{}-check} musi przejść. Oczekiwane porażki biegną z~\repopath{GITHUB\_ACTIONS=false}, żeby nie tworzyły adnotacji na PR \\
\repopath{package} & Po \repopath{build-test}. \path{scripts/verify-package.sh} pakuje \repopath{SecondKey.Standards} i~weryfikuje go na buildzie konsumenta; \repopath{dotnet pack} pakuje narzędzie; oba \repopath{.nupkg} trafiają do artefaktu \path{secondkey-standards-packages} (brak plików jest błędem). Nic nie jest publikowane w~nuget.org \\
\repopath{drift-check} & Po \repopath{build-test}, checkout z~pełną historią. Kryterium C4-c: publikuje źródło wydań -- gołe repozytorium w~\repopath{RUNNER\_TEMP}, do którego trafia \repopath{HEAD} jako \repopath{main}, z~tagiem \repopath{v<wersja z pack.json>}; składa aktualnego konsumenta (kopia skilla i~\path{Directory.Packages.props} z~\repopath{GlobalPackageReference} w~tej wersji); uruchamia \path{./actions/drift-check} dla aktualnego konsumenta i~trzech fixture (fixture z~\repopath{continue-on-error}); ostatni krok wymaga wyników \repopath{success/current}, \repopath{failure/behind}, \repopath{failure/disagree} i~\repopath{failure/unpinned}. Trzy fixture zostawiają oczekiwane adnotacje błędów \\
\repopath{release-guard} & Tylko na \path{refs/tags/v*}, po \repopath{build-test}, limit 5 minut. Tag musi być równy \repopath{v<wersja>} z~\path{catalog/pack.json}; \path{CHANGELOG.md} musi mieć nagłówek \texttt{\#\# [<wersja>] — RRRR-MM-DD} (dopuszczalne są pauza, półpauza lub dywiz); \path{generated/manifest.json} musi mieć tę wersję \\
\bottomrule
\end{xltabular}

\subsubsection{Pozostałe workflowy i~automatyzacja}

\begin{itemize}
  \item \path{codeql.yml} -- analiza statyczna C\# generatora i~drift checka na każdym pull
        requeście, pushu do \repopath{main} i~co tydzień (cron \texttt{41 3 * * 1}); build
        ręczny (\texttt{dotnet build} na \path{SecondKey.Standards.slnx} w~konfiguracji Release),
        uprawnienie \repopath{security-events: write}. Komentarz uzasadnia: generator pisze to,
        co słyszy każdy agent migracyjny i~egzekwuje każda bramka.
  \item \path{secret-scan.yml} -- gitleaks na pełnej historii (\repopath{fetch-depth: 0}; płytki
        klon skanowałby tylko wierzchołek) z~obrazu \path{ghcr.io/gitleaks/gitleaks:latest}
        (\repopath{:latest} celowo: zestaw reguł wykrywania jest wart tyle, ile jego najnowsza
        reguła), konfiguracja \path{.gitleaks.toml}, \repopath{-{}-redact}; na każdym pull
        requeście, pushu do \repopath{main}, co tydzień (cron \texttt{30 5 * * 1}) i~ręcznie.
        Przy porażce podsumowanie kroku podaje kolejność: najpierw rotacja danych
        uwierzytelniających, potem usunięcie ze źródła, dopiero potem czyszczenie historii;
        fałszywy alarm usuwa się przez przeredagowanie przykładu, a~nie przez wyjątek dla
        ścieżki.
  \item \path{scripts/setup.sh} -- onboarding jednym poleceniem: sprawdza .NET~10 SDK (bez
        niego kończy się kodem 1 i~adresem instalacji), szuka gitleaks albo działającego
        Dockera, instaluje hook pre-commit w~katalogu hooków git (inny istniejący hook
        zapisuje jako \path{pre-commit.backup}) i~uruchamia \repopath{validate} oraz
        \repopath{generate -{}-check}; ich niepowodzenie zgłasza jako ostrzeżenie
        z~poleceniem naprawy. Bez skanera kończy się ostrzeżeniem, że commity będą
        odrzucane, dopóki skaner nie zostanie zainstalowany; opcja \repopath{-{}-check}
        zamienia brak skanera w~błąd (kod 1).
  \item \path{scripts/hooks/pre-commit} -- skanuje zmiany w~stagingu gitleaksem
        z~\repopath{PATH} albo, gdy go nie ma, obrazem \path{ghcr.io/gitleaks/gitleaks:latest}
        przez Dockera; celowo odmawia commitu, gdy nie ma gitleaks ani Dockera, bo ochrona
        zależna od narzędzi dewelopera byłaby egzekwowana tylko pamięcią człowieka;
        \repopath{git commit -{}-no-verify} pozostaje głośnym obejściem, a~CI i~tak zablokuje
        push. Znalezisko hook rozpoznaje po kodzie wyjścia 2 (\repopath{-{}-exit-code 2}), bo
        gitleaks zwraca 1 zarówno dla znaleziska, jak i~dla błędu użycia; każdy inny kod
        niezerowy to awaria narzędzia, przy której commit także jest odrzucany („a scanner
        that did not run proves nothing”), z~podpowiedzią
        \texttt{./scripts/scan-secrets.sh -{}-staged}. Przy fałszywym alarmie hook radzi wąski
        wpis allowlisty w~\path{.gitleaks.toml} z~komentarzem, podczas gdy
        \path{.gitleaks.toml} i~podsumowanie joba CI każą przeredagować przykład zamiast
        dodawać wyjątek dla ścieżki. \path{scripts/scan-secrets.sh} odtwarza skan CI
        lokalnie: bez argumentu -- cała historia (to, co uruchamia CI),
        \repopath{-{}-staged} -- zmiany w~stagingu (to, co uruchamia hook), \repopath{-{}-dir}
        -- drzewo robocze bez historii; nieznany tryb kończy się instrukcją użycia i~kodem 2.
  \item \path{.github/dependabot.yml} -- cotygodniowe (poniedziałek 06:00) grupowane
        aktualizacje GitHub Actions i~pakietów NuGet; pakiet \repopath{SecondKey.Standards} jest
        wykluczony, bo występuje tylko w~fixture drift checka, celowo poniżej każdego wydania.
  \item Szablon pull requestu wymaga identyfikatora ticketu i~jego „done when” (dosłownie),
        opisu zmiany z~powodem oraz dwóch list kontrolnych (niżej).
        \path{.github/CODEOWNERS} przypisuje całe repozytorium opiekunowi.
\end{itemize}

Listy kontrolne szablonu pull requestu (\path{.github/pull_request_template.md}). Decyzja
wersji -- zmiana tego, co dostają konsumenci, wymaga decyzji; niepasujące wiersze się usuwa:

\begin{checklista}
  \item \textbf{major} -- a rule is reversed or withdrawn, or a severity is lowered
  \item \textbf{minor} -- a rule is added, or a severity is raised
  \item \textbf{patch} -- nothing required changes (wording, examples, links)
  \item No standards content changed
\end{checklista}

Weryfikacja:

\begin{checklista}
  \item \texttt{dotnet test} passes
  \item The generator's checks pass (\texttt{validate}, and \texttt{generate -{}-check} once the
        generator exists)
  \item \texttt{./scripts/scan-secrets.sh} is clean
\end{checklista}

\subsubsection{Dystrybucja}

\begin{itemize}
  \item \textbf{Skill} -- kopiowany z~repozytorium: z~tagu wydania albo, przed pierwszym
        wydaniem, z~\repopath{main} (sekcja~\ref{sec:standardy-modernize}, krok 3).
  \item \textbf{Pakiety NuGet} -- \repopath{SecondKey.Standards} (konfiguracja analyzerów)
        i~\path{SecondKey.Standards.Generator} (dotnet tool \repopath{secondkey-standards}) są
        pakowane w~jobie \repopath{package} i~dostępne jako artefakt
        \path{secondkey-standards-packages} każdego przebiegu CI. Nie są publikowane
        w~nuget.org; konsument umieszcza je w~feedzie lokalnym lub wewnętrznym. Artefakty
        wygasają po okresie retencji GitHub Actions (workflow nie ustawia
        \repopath{retention-days}), więc trwałego wydania pakietu nie ma.
  \item \textbf{Akcja drift check} -- używana bezpośrednio z~GitHub:
        \path{konradcinkusz/letsgolegacy.secondkey-standards/actions/drift-check@main}.
\end{itemize}

\zrodla{\path{letsgolegacy.secondkey-standards/.github/workflows/ci.yml},
\path{letsgolegacy.secondkey-standards/.github/workflows/codeql.yml},
\path{letsgolegacy.secondkey-standards/.github/workflows/secret-scan.yml},
\path{letsgolegacy.secondkey-standards/.github/dependabot.yml},
\path{letsgolegacy.secondkey-standards/.github/pull_request_template.md},
\path{letsgolegacy.secondkey-standards/.github/CODEOWNERS},
\path{letsgolegacy.secondkey-standards/scripts/setup.sh},
\path{letsgolegacy.secondkey-standards/scripts/hooks/pre-commit},
\path{letsgolegacy.secondkey-standards/scripts/scan-secrets.sh},
\path{letsgolegacy.secondkey-standards/.gitleaks.toml},
\path{letsgolegacy.secondkey-standards/README.md}}

\subsection{Status według WORKPLAN i~CHANGELOG}\label{sec:standardy-status}

Faza 01 (demonstracja). Identyfikatory ticketów odpowiadają backlogowi między repozytoriami.

\begin{xltabular}{\linewidth}{@{}>{\raggedright\arraybackslash}p{1.1cm} >{\raggedright\arraybackslash}X >{\raggedright\arraybackslash}X >{\raggedright\arraybackslash}p{3.3cm}@{}}
\toprule
Ticket & Rezultat & Zrobione, gdy & Status \\
\midrule
\endfirsthead
\toprule
Ticket & Rezultat & Zrobione, gdy & Status \\
\midrule
\endhead
R4 & standardy migracyjne Second Key, razem z~zasadami architecture-standards, które powtarzają, spakowane jako własny skill dla agenta modernizacyjnego Microsoft (modernize-dotnet) & Agent podejmuje skill w~tickecie P6 benchu (P6 uruchamia człowiek z~GitHub Copilot) & scalone (\#3); zrobione, gdy P6 (planowane) potwierdzi podjęcie skilla \\
C4-a & Standardy migracji jako Markdown z~metadanymi reguł (identyfikatory zmapowane na diagnostyki Portcullis) & Każda reguła ma identyfikator, wagę, uzasadnienie, przykład zgodny i~niezgodny & zrobione (scalone \#1) \\
C4-b & Generator: Markdown → katalog skilli, \repopath{.editorconfig} i~\repopath{.globalconfig}, pakiet NuGet & Tryb \repopath{-{}-check} w~CI zawodzi, gdy wynik jest nieaktualny & zrobione (scalone \#2) \\
C4-c & Drift check jako jeden krok CI repozytorium konsumenta (wersja przypięta vs opublikowana) & Repozytorium fixture z~nieaktualnym pinem zawodzi, aktualne przechodzi & zrobione (scalone \#4) \\
\bottomrule
\end{xltabular}

Faza 02 i~dalsze według WORKPLAN: standardy klienta dostarczane jako skille (jedno źródło na
estate) -- faza 02; kanały wydań standardów przez NuGet i~tagi git oraz przypinanie wersji
per repozytorium po stronie odbiorców -- fazy 02--03.

\path{CHANGELOG.md} ma jeden wpis: \textbf{[0.1.0] -- Unreleased}, sekcja \emph{Added}:
\begin{itemize}
  \item trzynaście standardów migracyjnych (\repopath{SK-MIG-001} … \repopath{SK-MIG-013});
  \item siedem zasad konstytucji przeniesionych na migrację (\repopath{SK-ARCH-001} …
        \repopath{SK-ARCH-007}): P4, P5, P9, P10 i~P15;
  \item wyniki generowane: skill \path{applying-dotnet-migration-standards}, wagi
        \repopath{.editorconfig} i~\repopath{.globalconfig} dla jedenastu reguł z~diagnostyką
        oraz pakiet NuGet \path{SecondKey.Standards};
  \item instrukcje skilla dla agenta modernizacyjnego GitHub Copilot: baseline legacy przed
        pierwszym zadaniem, reguły w~każdym zadaniu, bramka i~testy po każdym, format flag,
        kryteria sukcesu i~obsługa błędów;
  \item drift check dla repozytoriów konsumentów: \texttt{secondkey-standards drift-check}
        i~composite action \path{actions/drift-check}.
\end{itemize}

Stan z~README: standardy, walidator, generator, skill i~drift check istnieją; żadna wersja
nie jest wydana (otagowana).

\zrodla{\path{letsgolegacy.secondkey-standards/docs/WORKPLAN.md},
\path{letsgolegacy.secondkey-standards/CHANGELOG.md},
\path{letsgolegacy.secondkey-standards/README.md},
\path{letsgolegacy.secondkey-nopcommerce-bench/docs/WORKPLAN.md}}

\subsection{Ograniczenia i~otwarte punkty}\label{sec:standardy-ograniczenia}

Zebrane z~ADR, README, WORKPLAN, \path{LICENSE}, procedury, reguł, kodu narzędzia
i~dokumentacji Portcullis -- tak, jak opisują je źródła:

\begin{itemize}
  \item \textbf{Brak wydania.} Wersja 0.1.0 nie jest otagowana, więc drift check względem
        prawdziwego repozytorium zwraca \repopath{no-release} (kod 2), a~krok CI konsumenta
        kończy się błędem; dla P6 przed wydaniem krok pomija się.
  \item \textbf{Akceptacja R4 zależy od P6.} Tego, że agent podejmuje skill w~prawdziwym
        przebiegu, nie da się sprawdzić w~tym repozytorium; potwierdza to człowiek
        z~GitHub Copilot kontrolami z~kroku 6 procedury. P6 ma w~WORKPLAN benchu status
        \emph{planned}, więc nikt dotąd nie sprawdził, że agent stosuje standardy.
  \item \textbf{Otwarte do P6}: zachowanie \texttt{metadata.discovery: preload} dla skilla
        repozytorium w~poszczególnych klientach oraz rozbieżność między przykładem
        jednoplikowym na stronie Microsoft a~formą katalogu \path{<name>/SKILL.md}.
  \item \textbf{Pakiet poza nuget.org.} Konsument potrzebuje feedu lokalnego lub
        wewnętrznego; drift check nie porównuje z~feedem, tylko z~tagami.
  \item \textbf{Pin przez właściwość MSBuild} wymaga \repopath{-{}-pinned-version}; drift check
        nie wylicza MSBuild.
  \item \textbf{Akcja buduje narzędzie przy każdym uruchomieniu} (dziesiątki sekund);
        alternatywą jest zainstalowane narzędzie z~własnego feedu.
  \item \textbf{Kontrakt z~Portcullis jest ręczny.} Lista \repopath{portcullisRules}
        w~\path{catalog/pack.json} jest utrzymywana ręcznie; walidator sprawdza tylko, że reguła
        wskazuje identyfikator z~tej listy, a~żaden test nie porównuje jej z~rejestrem
        Portcullis (Portcullis przypina jako kontrakt w~testach tylko cztery identyfikatory
        migracyjne).
  \item \textbf{Diagnostyki Portcullis wyprzedzają analyzery.} Wagi dla identyfikatorów,
        których analyzera nie ma w~buildzie konsumenta, nie mają efektu, dopóki bramka ich nie
        dostarczy.
  \item \textbf{Walidator skilli Microsoft poza CI} z~powodu warunków licencyjnych jego
        repozytorium; jego istotne ograniczenia pokrywają generator i~testy.
  \item \textbf{Dziewięć z~dwudziestu reguł nie ma sprawdzenia wykonywalnego} (SK-ARCH-003,
        SK-ARCH-004 i~SK-MIG-007 do SK-MIG-013): to instrukcje dla agenta, a~w~zakresie
        zachowania ma je weryfikować replay i~porównanie Second Key.
  \item \textbf{Wagi standardów tylko w~buildzie.} \repopath{portcullis scan} (CLI, akcja,
        kontener) raportuje poziomy domyślne Portcullis i~nie stosuje konfiguracji standardów,
        podobnie jak SARIF i~\texttt{sk gate}; dla pięciu diagnostyk różnią się one od wag
        standardów, a~czterech diagnostyk \repopath{CA} skan w~ogóle nie uruchamia
        (sekcja~\ref{sec:standardy-wagi-kanaly}).
  \item \textbf{Licencja.} Wszystkie prawa zastrzeżone, a~model licencjonowania jest decyzją
        otwartą (\path{LICENSE}); dotyczy to także pakietu \repopath{SecondKey.Standards},
        który niesie ten plik.
  \item \textbf{Comparery sortowania poza analyzerami.} CA1310 znajduje wywołania, których
        domyślne porównanie zależy od kultury, ale domyślnych comparerów, do których wracają
        \repopath{OrderBy}, \repopath{List<string>.Sort}, \repopath{Array.Sort},
        \repopath{SortedList} i~\repopath{SortedDictionary}, nie zgłasza żaden analyzer;
        trzeba je wyszukać (SK-MIG-005).
  \item \textbf{Linki w~stopkach przed wydaniem.} Stopki \repopath{SKILL.md} i~plików
        referencyjnych (21 z~27 wygenerowanych plików) linkują źródła pod tagiem
        \repopath{v<wersja>}, a~0.1.0 nie jest otagowana, więc te odnośniki prowadzą
        do nieistniejącego tagu (GitHub zwraca 404); zaczną działać po otagowaniu wersji.
  \item \textbf{Drift check porównuje numery wersji, nie treść.} Wersja, której wpis
        w~\path{CHANGELOG.md} mówi \repopath{Unreleased}, może się jeszcze zmienić, a~kopia
        skilla wzięta z~\repopath{main} przed wydaniem niesie jej numer
        w~\repopath{metadata.version}; dlatego krok~3 procedury każe zapisać commit,
        z~którego pochodzi kopia.
\end{itemize}

\zrodla{\path{letsgolegacy.secondkey-standards/docs/adr/0003-generated-tree-and-versioning.md},
\path{letsgolegacy.secondkey-standards/docs/adr/0004-skill-for-github-copilot-upgrade.md},
\path{letsgolegacy.secondkey-standards/docs/adr/0005-drift-check.md},
\path{letsgolegacy.secondkey-standards/docs/USING-WITH-MODERNIZE-DOTNET.md},
\path{letsgolegacy.secondkey-standards/README.md},
\path{letsgolegacy.secondkey-standards/LICENSE},
\path{letsgolegacy.secondkey-standards/standards/SK-MIG-005-explicit-string-comparison.md},
\path{letsgolegacy.secondkey-standards/src/SecondKey.Standards.Generator/Generation/SkillEmitter.cs},
\path{letsgolegacy.secondkey-standards/src/SecondKey.Standards.Generator/Drift/DriftChecker.cs},
\path{letsgolegacy.secondkey-standards/catalog/pack.json},
\path{letsgolegacy.secondkey-nopcommerce-bench/docs/WORKPLAN.md},
\path{letsgolegacy.portcullis/README.md},
\path{letsgolegacy.portcullis/docs/CONFIGURATION.md},
\path{letsgolegacy.portcullis/tests/Portcullis.Engine.Tests/Rules/RuleRegistryTests.cs}}
